\documentclass{design/chronicle}
\BookTitle{The Debt of Warmth}
\hypersetup{pdftitle={The Debt of Warmth -- Volume 1},pdfauthor={Written by Hy4 Flash}}
\begin{document}
\frontmatter
\includepdf[pages=1,pagecommand={\thispagestyle{empty}\CRnobgthispage}]{build/cover/cover.pdf}
\CRtitlepage{The Debt of Warmth}{A surgeon’s account of Sarn}{Volume 1}{Chapters 1–2}{Written by Hy4 Flash}
% Colophon — printed on the verso after the title page of every volume.
\clearpage\thispagestyle{empty}
\vspace*{\fill}
\begin{center}
\begin{minipage}{0.8\textwidth}\centering\small\color{faded}
{\color{ornament}\pgfornament[width=2.4em]{11}}\par\vspace{8pt}
{\FellFont\itshape This chronicle is written, chapter by chapter and without end,\\
by the language model \textsc{hy4-flash-free}, working through OpenCode.\\
The world, its people, its maps and every chapter ornament are its own.}\par\vspace{8pt}
{\FellFont\itshape The book design, page texture, borders and cover\\
were made by Claude, which tends the press and keeps the ink flowing\\
but writes no word of the story.}\par\vspace{8pt}
{\FellSC Set in EB Garamond, Cinzel and the Fell types\\
with Lua\LaTeX{} on a parchment of numbers.}\par\vspace{10pt}
{\color{ornament}\pgfornament[width=2.4em]{11}}\par\vspace{8pt}
{\footnotesize This is a work of fiction for adult readers. It contains violence,\\ cruelty, slavery and other darkness of its invented world.}
\end{minipage}
\end{center}
\vspace*{\fill}
\clearpage

\tableofcontents
\mainmatter
\CRarc{1}{The Cold Stone}{a warm stone, already going cold}
% Generated by forge/build.py from meta.json and the author's scene files.
% Edit the scene files, not this file.
\CRdecoration{book/chapters/ch-0001/decoration.tex}
\CRchapter{1}{The Stone Going Cold}{Day 1 · Stone 7, Y.H. 812, a Seventhday · the Ilm under ice, three days before the Five Cold Days}{The birth room off the Clay Yard in Low Kiln, Sill, Grath; a brief breaking-through of the expressway near Chicago}
\CRepigraph{A stone at the belly is a stone borrowed from the fire. It goes cold. Somebody has had the warmth back, and nobody says whose.}{A hearth-maid’s saying, Low Kiln}
\dropcap{T}{he stone} had been warm when the midwife put it against him and was warm no longer.

It lay under his belly in the fold of the cloth, a small river stone the colour of skimmed milk, giving up its heat with the slow reluctance of a thing that has been asked and has agreed. It had been warmed at a neighbour's fire --- at the kiln-house down the yard, where somebody kept banked coals through the cold --- and carried here in a rag, still hot from a hearth that was not this hearth. That was a favour. It was not entered anywhere. \person{Daniel Marchetti} had known households where a kindness goes in warm and comes out written in a hand you cannot read, and men who warm a stone for a cot and never ask who paid for the fire. He could not have said which kind this was. He could only feel the stone going cold against him, a little at a time, the way he had once felt a patient's core temperature fall under his palm.

The child was \person{Corin}, two hours old, and called nothing yet that he could answer to. He lay on a fold of grey wool in the front room of \place{House Silt}, one of two clay-daubed rooms that opened off the \place{Clay Yard} in \place{Low Kiln}, the poor quarter of \place{Sill}, in \place{Grath}. The floor was trodden ash, packed by generations of feet until it had the dull spring of old bread, and along the skirting board ran a grey ridge the width of a thumb: \lore{the Ash}, sifted out of the walls and the pots and the lives that had been kept here, left because nobody sweeps it before the Five. The Five were three days off. Outside, the \place{Ilm} lay under its skin of ice, and the cold came through the door that would not shut and went around the room looking for a place to sit.

It smelled of wet wool and tallow and clay dust, and under those the sour mineral edge of the Ash, which is the smell of a room that has been warm a long time and has kept things it should have let go. A woman sang somewhere in the house, a low song with the words flattened out by work, and the tune went under the skin and stayed. He understood the words. He could not make one.

\person{Alda Reed} worked with steady hands. She had the stoop of a woman who had spent her life leaning over other people's births, and the forearms of one who had lifted other people's children; her nails were scrubbed to the quick and still had clay under them, a thin brown line that no soap had ever got out. She addressed the stone first, turning it in the rag with two fingers the way a cook tests a pot, then the slip pot on the hearth, then the child, then the mother. She did not address the father, who was not in the room and would not be addressed. This was not cruelty. It was the correct number of words for a room where somebody was still in pain.

The slate lay within reach of her right hand: a flat grey slab, a dusting of chalk, a line of writing in a strange neat hand that Dan did not look at twice because he had other things to look at and no way yet to look at anything. Her hand went to the slate before it went to the child. He registered the movement and filed it where he filed everything he could not use yet, under the heading of things that would matter later.

“Warm the stone first,” she said, and her voice was low and even, the voice of a woman who has said the same sentence four hundred times and never once said it carelessly. “A cold stone is a stone, not a kindness.” She said it to the stone, or to the room, or to \person{Dessa Silt}, who was in no condition to argue. Then something about the slip, about not setting it near the fire, about paying for it in the thaw --- practical instruction addressed to the pot before it was addressed to a person, because the pot would answer and the person would not.

Dessa lay with her face turned toward the wall and her breath coming in a rasp that Dan knew before he knew that he knew it: expiration long, inspiration short, a small catch at the end that was not relief but the body trying again. Her right hand went to her left forearm, to a pale scar the width of a sieve-wire, and rubbed it in a slow circular press while she counted under her breath.

“Nineteen,” she said. “Eighteen. Nineteen.”

She had walked those nineteen steps every day for six years, from the slip pit to the drying rack, and counted them then for the length of the working day, for the bits at the end of it, for the number of times her feet would go before she could stop. Now she was counting them for a different reason and Dan could not know what it was. He could only hear the count and know, with the part of him that had once stood under the lights of a trauma bay and sorted twenty people in ninety seconds, that she was not counting steps at all. She was holding something out of the room by counting it in.

His eyes --- his new eyes, wet and unfocused, not yet good at holding a face --- went to the angle of her wrist instead of to her face. The scar sat where sieve-wire had cut it, and the way she pressed it told him she was deciding something she did not want to decide. He read the pulse at the wrist before he understood that he was reading it: quick, thready, a little irregular at the end of expiration. He read the colour under the nails, grey-blue at the beds, and the fine tremor in the hand, which was not from cold. He counted the exits without deciding to --- the door to the yard, the window with its oiled paper, the gap under the lintel where light came in thin --- and then counted them again to check the count. This was what his mind did now. It had nothing to do, so it took the room apart.

He watched Alda's hands because they were the steadiest thing in the room. They did not hurry and they did not hesitate; they went from the stone to the slip to the cloth to the mother in an order that had been decided years before Dan was born and would not be changed now. He had worked with surgeons whose hands could not be hurried and with surgeons whose hands could not be trusted, and he could have told you which was which within a minute of watching them tie a knot. Alda tied nothing, but the order of her reaching was the same kind of signature. When she touched the little tin at her bodice --- once, briefly, before she went back to the child --- his eye caught the gesture and held it. He did not know what it was. He knew that it was a thing done every time and never for the first time.

The stone was giving up its heat at a rate he could almost measure. Not fast --- a stone gives heat the way a bank gives change, reluctantly, a little at a time --- but steadily, and the rag between them was damp now where his body had sweated against it, and damp cloth draws heat faster than dry. He had watched patients cool in a white room far away and had known, without a thermometer, the exact point at which the surface warms and the core falls: \lore{the Afterdrop}, the worsening that arrives wearing the face of improvement. He knew it the way a priest knows sin. There was no rewarming here, no blanket but the scrap of wool, no patient but himself, and still he knew the shape of it --- knew that the warmth going out of the stone was going into a body that could not hold it, and that in a few minutes the body would be colder than it had been before the stone was put there.

He knew the Afterdrop when he saw it. He wanted the breast. Both were true and the second was louder.

The body that wanted was not the body that knew. The knowing body was a large soft man with forearms like a butcher's and a left knee that clicked in the cold, a man who had reduced a thousand shoulders in the dark and could find the thin place under skin with two fingers and a name for it. The wanting body was two hours old, had never held anything up, and was cold and hungry and furious about it. It drooled. It could not hold up its head. When it tried to turn over it stayed where it was, and the grey wool that Alda had tucked around him had come loose at one corner and would not go back. He had set a femur in a ditch in the rain and he could not keep a scrap of wool on his own shoulder.

The joke was there, dry and small, and it cost him something. He did not laugh, because laughing is a thing the body does with air and the body had other plans for its air. It wanted. That was all it had ever wanted, and it wanted now with a whole man's urgency in a vessel the size of a loaf.

He could not speak. He could not reach a cup. A trauma surgeon who could intubate in blood and in the dark, who could cut a way into a destroyed throat with a knife and a tube, lay in a fold of cloth and could not lift his own hand to his own mouth. There was a cup somewhere on the hearth --- he had seen Alda set it down --- and the distance to it was two paces, which is a long way when you have no hands for the journey. He could not have said the word for cup in Grathic. He could not have said it in English either, because the body had not yet learned to push air into words. He had a whole ledger in his head and no way to open it.

It was the first small grief of the new life, and it was not for the dead, or for the daughter in a country that had no name here, or for the call he had not made. It was for a pair of hands that could do nothing. He had spent twelve years being useful in the only way he knew how, which was to decide quickly and carry the decision, and here was a room that needed nothing decided and a body that would not do what it was told. The competence was intact. There was simply nothing behind it. It sat in him like an instrument in a dead man's hand.

Outside, something moved on the frozen river: not the river, but the ice under it, or the wind over it, a small sound like a hand rubbing dry paper. Dan listened to it because listening was what he had left. He counted the song against Dessa's breathing and got a ratio that meant nothing and kept it anyway. He had been counting things for as long as he could remember, pulses and ribs and steps from the bed to the door in rooms he would never sleep in again, and he had never once been told that it was a gift. It was a habit. It was the way he had learned to stand in a room where he could not help.

A woman came in from the yard and stood just inside the door with the cold still on her sleeves. \person{Mira Silt}. Eel slime had dried on the wool in grey patches, and she smelled of river silt and salted perch and the cold off the water. There was something in her pocket that ticked when she moved her thumb, small and hard; she flicked it once against a nail and put it away. She did not sit. There was a stool by the wall and a clear place on the table where a slate might be laid, and she went near neither.

“Hold him?” she said, and it was not clear whether the question was for anyone. “You want me to. You don't. I can see which one's winning.” She looked at the child, or at the place where the child was, and did not reach for him. Then she said nothing further, and stayed standing.

Alda finished with the slip and wiped her hands on a rag that had been a shirt once. She lifted the child --- not roughly, but as a woman lifts a sack of meal, with the whole weight taken in one movement --- and set him back against the folded cloth, and in doing so turned him so that the stone slid a little and lay closer to the hollow of his belly. The grey wool went over him again. This time it stayed for a moment. Then an arm came out of it, small and brown and determined, and the wool was off at the corner again.

“Keep him covered,” Dessa said, and her voice went flat and finished: nothing more to be said on it. She reached over and pressed the wool down at the corner with two fingers, pressing as if to seal something, and let her hand lie there a moment longer than the cloth required. The pale scar crossed the back of her forearm in the tallow light. Then she took her hand away and went back to counting, and what she was holding out of the room went back in.

Mira took her hand off the door frame and went out the way she had come, leaving the smell of silt and salt behind her and the small ticking stopped. She had not held him. She had not sat down. Dan watched her go with the part of him that watched people leave rooms and noted, without knowing why, that she had kept her back to the slate the whole time.

Alda went back to the slate and stood over it with the chalk in her fingers, and the room settled around the standing of her body the way water settles around a stone dropped in it. The banked fire on the hearth shifted and settled, and a small flame came up under the ash and went down again. She did not look at it. She wrote something on the slate, blew the dust off it, and set the slate down within reach of her hand again.

Dan listened to the room the way he had once listened to a corridor at four in the morning: for the note under the noise. The fire had a small irregular tick as the clay of the flue cooled. Dessa's breath had its rasp and its catch. Alda's skirt dragged on the ash floor when she moved. Under all of it, under the song that had stopped and the count that went on, there was the stone, and it had a sound too if you put your ear to it --- not a sound in the air but a sound in the body, the small settling of warmth leaving a thing that will not be warmed again. He had heard that sound in patients. He had never heard it in himself.

Under the child the stone had reached the point where it no longer gave heat but only failed to take it. The warmth that had been in it was in him now, or had gone into the room, and the room was warm around him with its fire and its wet wool and its breathing bodies, and nobody in it was cold. Nobody in it felt the stone go. Alda's hand was on the slate. Dessa was counting under her breath. Mira was out in the yard with the cold still on her sleeves. And under the grey wool, in the hollow of a belly that did not yet know it was hollow, the stone went on giving up the last of what it had --- quietly, without a sound, and with nobody the wiser.

\scenebreak
\person{Alda Reed} lifted him then, in the room off \place{the Clay Yard} in \place{Low Kiln}, and the stone did not come with him.

It stayed in the fold of grey wool where his belly had been, a small pale disc with the damp cloth cooling around it, and she did not look at it. She had come to lift a child and clean him, and the lifting was done with the whole weight taken at once, the way a woman lifts a sack of meal, and the cleaning was done with a rag she had warmed by holding it a moment inside her bodice before she touched him. \person{Daniel Marchetti} noted the rag. He noted the temperature of it, which was body-warm and not water-warm, and the order in which she wiped: face, neck, the crease at the groin, the back of the knees, the place where the cord had been cut. It was a good order. In a white room far away he had watched nurses clean a man who could not clean himself and had judged them by just such an order, and he judged her now, and the judgment was useless, because he could not have told her any of it and would not have been believed if he had.

The room went on with its work. \person{Dessa Silt}'s count ran under it, nineteen and eighteen and nineteen, and the fire gave its small irregular tick, and the door that would not shut let the cold in around the edge and went on letting it in. A woman sang somewhere below, and the song had its words worn flat by work, and he had already stopped listening for the tune. \person{Corin} lay where she put him, small and brown and slippery with slip, and did what he was told by the only part of him that had ever learned to obey.

Then the light changed.

Alda turned him to get the back of him, or turned him to wrap him, and the banked fire came up under the ash and put a low red light across the room, and the light went along his ribs and down his left forearm, and there was something on him that had not been there when he was a man.

It was red. Not the red of a bruise coming up under the skin, not the blue-black of a cord or a pressure mark, not the dusted pattern of petechiae he had seen on a dozen faces in the resuscitation bay. The lines ran finer than any of those and they branched. They started somewhere at the third rib on the left side, where the firelight caught them, and went down and outward like a thing growing, and one arm of them crossed the forearm and stopped just short of the wrist, as if it had been told where to stop. The colour was the colour of a coal that has given up its flame and kept its red, a deep ordinary red, not bright. It sat in the skin. There was no swelling at it, no heat he could feel through his own skin, no edge of injury; it was as if somebody had taken a fine brush and a red that came from inside the body and written a small branching sentence across him, and then closed the skin over it.

He made a differential because that was what he did when a thing frightened him. A capillary mark. A bruise older than the bruises of birth. A stain from the cord, or from the hands that had pulled him. A birthmark, which was a word that meant nobody had looked closely enough. He ran through them the way he had run through a hundred rashes and patterns in a white room, and every one of them was wrong in the same way: they were things that happened to the skin, and this was a thing that had been written in it. The lines branched the way a burn branches when it follows warmth along a limb, and there was no burn. He had seen a man who had come out of a fire with just such a map on his arm, the skin whole and the red underneath it, and he had called it a burn map and entered it in a chart. He could enter nothing here. He could not lift his hand to touch the place. He could only lie in the red light and let the mark be the one thing in the room he could not put a name to.

Alda's hands stopped.

It was half a beat, no more --- the pause of a woman who has laid her hands on four hundred children and has never once had reason to take them back --- but Daniel Marchetti had spent his life watching hands stop. He had seen a senior resident's hand stop over a film, and a nurse's hand stop half an inch from a chart, and he had learned to know before anyone spoke that the stop meant the thing on the table had changed and nobody had yet decided what to call it. His own hands, which were soft brown wrists with milky nails, lay where they lay. He could not have pointed at the mark. He could not have said what it was. He knew the shape of the stop, and he filed it, and the filing was the only thing in the room that belonged to him.

Her two fingers left the rag.

They went to the corner of the mark, to the place where the branching ended under his arm, and hovered there over the skin without touching it; then they came back across her own body in a small cutting of the air, forehead and breast and shoulder, finished before the motion had properly begun. It was an old thing done quickly, the way a woman crosses a threshold she has crossed ten thousand times and still touches the wood. It took less than a breath. Then her hands were on the rag again, and the rag was on him, and the moment was gone as completely as if it had never been in the room.

Except that she had seen it. Alda Reed looked up from the child, and her face went through a small change that had nothing in it of surprise and everything of a door being shut quietly. She had made the sign over a child who had come from nowhere anybody could name, and she had made it where the mother might have seen it, and where the aunt standing in the door might have seen it, and there was no account in Grath for a thing like that. She covered it the way a woman covers a spilt cup: with a word.

“\lore[birth-flower]{Birth-flower},” she said.

“What's that?” Dessa said, from the pallet, without opening her eyes.

“That's what they call them.” Alda wiped along his forearm and would not look at the place she had wiped. “They come like that. I saw one in Nine Pots in Y.H. 806, on the fourth day of Stone, on a potter's third child, and there was nothing to it; he's a carter in Stonemouth now.” She blew the dust off the slate without looking at it. “Not a thing to write home about.”

The count stopped. Dessa's eyes opened a little in the tallow light. “Is he hot?” she said.

It was not a question about the mark. It was the question a woman asks about a child who has been given a stone to hold, and Dan knew the shape of it before he knew the words, because he had heard it in a hundred rooms from people standing at the foot of a bed.

“He's warm,” Alda said. “That's the whole of it. Warm and marked and breathing, which is more than some.”

“He's the little hot one,” Dessa said, and the words went flat and finished, and the count went back to nineteen.

Daniel Marchetti heard the date the way he had once heard a man clear his throat before giving a history that was not quite true. She had given a year and a day and a third child and a trade, none of which had been asked for, and all of which had the flat solid sound of a thing placed carefully over a hole. It was the small exactness of a woman choosing what to set down and what to leave out. He had signed a chart once without reading the nurse's note underneath it, and the chart had been exact, and the note had been the thing that mattered, and the institution had called it a systems failure. He recognized the sound of it now in a clay room with a tallow dip, and it made nothing in him louder than curiosity, which was how he knew the infant's body had won again: he wanted the breast, and he wanted to know what the woman had hidden, and the wanting was bigger than the knowing by the whole width of a man's life.

\person{Mira Silt} had not come nearer. She stood just inside the door with the cold still on her sleeves, and her thumb went once into her pocket and came out, and the small hard thing ticked, and she put her hand back on the frame. She had not sat down. She said nothing about the flower. The grey ridge of \lore{the Ash} ran along the skirting board behind her, and nobody in the room would sweep it before the Five.

Then the smell came.

It came sideways, as smells do when a memory has no door to come through: bitter coffee that had sat too long on a hot plate, and under it the clean metal edge of a room scrubbed four times a day with the same mop. The tallow went thin. The clay went thin. The wet wool went away.

The white room is there.

The light in it is overhead and flat, and it comes down on a body on a table and on a nurse holding pressure in a cavity with both hands. Snow comes at the windows in a brown rush. Outside the glass the sodium lamps make orange tunnels in the snow, one tunnel behind another, and the tunnels go up \place{the expressway} toward the split, somewhere in \place{Chicago}, in February, at 2:17 in the morning. A flare stands in the snow and puts a red rim on everything it touches. The flare smokes cheaply; the smoke is acrid in the throat and lies under the sweetness of the fuel.

The wipers have stopped. They have been stopped for some time. Snow builds along the bottom of the windscreen in a fan and the fan grows. The engine is still running. Fuel goes out of the broken van with a hiss that changes over several minutes into something sweeter and worse, and the smell of antifreeze lies under it, and the cold comes in through the door he has left open.

He is working. Airway, bleeding, the end of the problem. There is a woman in the back seat of the minivan who is still moving, and snow has gone up his sleeves, and his fingers have stopped being reliable and he is angry at them. The order of the work does not matter now, because the fuel has caught.

The phone is in his coat pocket against his thigh. It vibrates. The screen is lit where his coat has ridden up, and \person[Nina Marchetti]{a name} is on it, and he does not look at the name. There is a boat on a lake in summer that he has promised for nine years, and a girl who grew up while he worked nights, and a sentence he has composed so often that nothing is left of it but weight. The weight is in the pocket with the phone. It is heavier than the phone. He does not look.

The cold is leaving his hands.

It goes out of them like water out of a wrung cloth, and the heat comes up through his sleeves from every direction at once, a quick enormous heat against a cold that has been there all night, so much of it that it is like being held. There is metal in his mouth. The hiss of the fuel goes on. The orange light goes on. And underneath all of it, finished and complete and with no words in it that anyone could have written down: \emph{I have left the job unfinished.}

At the edge of the flare light \person[Marra]{a woman} stands in a coat the colour of wet ash. She has not come into the light. The flare puts a red rim on the snow at her feet and leaves the rest of her grey, so that she is a shape and nothing more: a coat, a head, a hand. She holds a slate in her hand, and the hand moves. It moves without hurry, short strokes at the top of the slate, then a pause, then more strokes. She has not touched the child. She has not once looked up. The hand moves. It does not stop when the fire takes the van.

Two lights are in the world now. One is small and cold and under his leg, a rectangle of light with a name on it that he will not look at. The other is the flare, red on the snow, and the woman is at the edge of it with the hand moving, and the writing goes on. He holds the two lights together for one breath: the light he refuses, and the light that refuses to stop.

Then the white room was gone.

The smell of coffee stayed in his nose for a breath after the tallow returned, and then that went too. The tallow came back, and the wet wool, and the clay, and the count. Nothing in the room had moved in the time it takes a memory to arrive and go. Dessa was on nineteen. The fire gave its tick. Alda's two fingers came down on the corner of the wool and pressed it flat, the way a woman presses down a cloth that will not stay where it is put, and the word birth-flower settled over the room like a lid.

\scenebreak
Under the word the room went on with its small noises: the dip guttering, the draft at the door, the low song from below with its words worn flat, and \person[Dessa Silt]{Dessa}'s count running under all of it like a second pulse that belonged to the floor. No one had moved in the time the word took to settle. \person[Alda Reed]{Alda}'s two fingers had come down and pressed the wool flat, and that was that, as though a lid had been put on a pot that had always had a lid.

\person{Corin} had been laid on his back on a board-bed with a fold of wool under him. The wool was thin. The boards beneath it were cold, and the cold came up through the wool in a slow even way that he could have put a number to if he had had a mouth for numbers. Above his head the floorboards ran off toward the door, and he counted the nails in them because there was nothing else in the room that kept still long enough to be counted. Four on the near board. Four on the next. On the third board a fifth thing that was not a nail at all but a knot the width of a thumb, dark in the tallow light, with the grain going round it the way silt goes round a stone in a slow river.

He counted them without deciding to. He had counted the tiles in the ceiling of a resuscitation bay once, while a resident explained a thing he already knew; he had counted the staples across a man's chest after he had finished with him and while somebody else was still talking. It was a habit with no business in a room like this and less business in a body like this one. He did not yet know that what he was counting was the way out. He knew only that there were four, and four, and a knot where a fifth nail ought to have been, and that the number of things in a room was something a man could know without being told.

There was a wrong familiarity in the room that he had not been able to shake since he woke in it: the circle the dip threw on the wall, the shape of a slate within reach of a woman's hand, the way a lamp puts a lid of light on a table. He had come back to a room he had never been in, and the room had been rearranged in his absence, and somebody had put a stone at his belly and a word over his head, and he was expected to lie still under it. He lay still under it.

The light from the dip reached in a small circle and left the corners to \lore{the Ash}. He lay inside the circle and read what he could reach with his eyes, because his eyes were the only instruments he still had and they worked. Dessa's hand hung over the edge of the pallet near his face. He could see the pulse at the wrist move the skin a little, a quick light tapping under the pale surface, and he could have told you the rate of it the way another man might tell you the time. He looked at the nails of that hand: the beds were the colour of tallow with a thin blue seam at the edge of them, and one nail had a white fleck in it that had not come off with washing. He looked at the inside of the left forearm where the pale scar lay, a long old seam where a sieve-wire had opened the skin six years ago and been left to close itself any way it could. She rubbed the scar with the thumb of the other hand when she was making up her mind about something. He had seen that gesture on the woman herself and on a hundred women before her. It was the gesture of a person standing at the edge of a thing and deciding how much of it to say.

He made nothing of it. That was the strange part. He had made everything of it once, a pulse, a nail bed, a scar, the small evidence a body carries, and now he lay on a board-bed with the evidence three inches from his face and could have named every part of it and could not lift one finger to touch any of it. He had reduced a thousand shoulders in rooms that were warm and bright and full of people who did what he said. He knew \lore{the Afterdrop} when he saw it, and there was a woman two pallets over whose breathing had the long shallow shape that comes before one, and he noted the shape of it and put it away, and there was nothing in the room he could do with the knowledge. The clinical ear was wide awake. The hands were soft brown wrists with milky nails, and they lay where they lay.

\person{Daniel Marchetti} did not look up at the cough. It came up from the pits in cold weather, which was most of the year, and it had been in Dessa's voice all winter, and Alda did not look up at it and the woman in the door did not look up at it, so he did not either. It was weather. It belonged to the room the way the draft belonged to the door. He let it go.

Under the cough her count went on. Nineteen and eighteen and nineteen. She had walked nineteen steps from the slip pit to the drying rack every day for six years, and now she was counting them for a different reason, and the reason was the small thing lying on the board-bed with its face turned to the tallow. The count had been a way of not thinking on the way to the rack. It was a way of not thinking now.

She spoke without opening her eyes. “A body can be warm and still be gone,” she said. “You learn that in the pits.”

Alda made a small sound that was neither agreement nor disagreement, the sound a woman makes at a proverb she has heard before and has no wish to argue with. She went on folding the wool at the foot of the board-bed.

Dessa's thumb went to the scar. “He's the little hot one,” she said. “That's what he is. I'll call him that and no different.”

Alda answered with a date, the way she answered most things. “Stone seven,” she said. “The seventh day, and \lore{the Five} come in three days, and that's a long cold coming for a child with a stone at his belly. Warm the stone first. A cold stone is a stone, not a kindness.”

“I warmed it,” Dessa said. She opened her eyes a little. “At the fire across \place[the Clay Yard]{the yard}. I carried it in my apron and it was warm when I put it there. Is it still?”

“It's a stone,” Alda said. “It does what stones do.”

Dessa's thumb went back to the scar, and she got specific, the way she got specific when a thing was too large to hold any other way. “Nineteen out, nineteen back. Twice a day since Y.H. 806. I never missed a step. Not one.”

The door that would not shut let in a narrow ribbon of cold, and \person{Mira Silt} came in through it without knocking and without finishing whatever she had been saying.

“—and they say at \place{the Weir} the ice'll hold three days, so that's the eels sorted, and if the Feast is kept there's barley in it for somebody, though which somebody nobody is saying.” She said it to the room at large, to the pot and the fire and the doorframe, and only then to the people. There was eel slime dried grey on her sleeves in stiff creases where she had rolled them and forgotten to unroll them, and the smell came off her of river silt and salted perch and the cold off the water, and under that the tallow smoke of a house that had been up before the fire was lit. The top of her left ear was not there; a hook had taken it, she said, and the scar said a knife, and nobody in \place{Low Kiln} had ever made anything of the difference. She kept a small smile on her face the way another woman keeps a door on the latch: it could be shut without anybody seeing the movement.

She put one hand in her pocket and brought it out with an eel-bone between her fingers, and flicked it against her thumb. Tick. She did it again. Tick. It was not a large sound but it was a sound a person could not stop hearing once it had started, and Corin lay with his eyes open and listened to it the way he had once listened to a drill in the next room, counting the beats between.

Alda's slate lay within reach of her hand on the bench by the door. Mira looked at it, and the smile did not move, and she chose the space beside the door and stayed standing in it.

“Sit,” Dessa said.

“Sit where?” Mira said. “I'll sit where there's no slate. That's the only kind of chair I can abide.” She flicked the bone. Tick. “You want to know. You don't want to know. I can see which one's winning, and I'm not saying which.”

Dessa opened her eyes properly. “Did you tell anyone?”

“Tell them what?” Mira said. “That you bore a child? Everybody in Low Kiln knows you bore a child. There's no secret in a child being born. The secret's in what happens to it after.” She flicked the bone once more and put it back in her pocket, and the small sound stopped, and the room was larger without it. “I didn't tell anybody anything. There's nothing to tell. That's the whole of it.”

She came across then and put her hand on the fold of wool at the head of the board-bed and pushed it back under him, briskly, with two quick movements, before anyone could thank her or make a speech about it; and her face while she did it was the same face she had been wearing at the door, shut and latched.

Under all of this the body had been getting louder.

The surgeon's register wanted the room laid out: the pulse at the wrist, the shallow breathing, the scar, the woman at the door with the bone in her pocket and the answer already chosen. He could have drawn you a map of who in that room was telling the truth and who was arranging it. He could not lift his head off the wool. The head went where it wanted, which was nowhere, and lolled to the tallow side, and the wool that Mira had pushed back came out again at the shoulder in a small soft curl, and he could not have put it back if his life had depended on it. And his life did depend on it, in the small way a newborn's life depends on a square of wool, which is the whole way.

He had reduced a thousand shoulders. He could not keep a blanket on his own body. The joke was clean and it had a good edge to it, and then the edge turned, because under the joke there was a hunger with no name in it and no hand to bring it, and the hunger was not a thought he was having about a body, it was the body itself, immediate and impatient and certain that the world owed it a breast. He knew the afterdrop when he saw it and he wanted the breast, and the wanting was bigger than the knowing by the whole width of a man's life, and it was bigger and louder because it had a mouth.

A sound came out of him. It was not a word. It was the small wet sound a hungry infant makes when it has found that crying is a thing it can do and that nobody has yet come to put right the thing it is crying about. His mouth worked. A thread of drool went down onto the wool. The song from below went on, and the dip guttered, and Dessa did not open her eyes, and Alda went on folding at the foot of the bed, and Mira stood where she stood with her face shut. Nobody came to put it right.

He had been a man who decided which bodies got a room and which got a corner, and now he lay in a corner of a room that had decided something about him, and a square of wool was more than he could manage, and the sound he made was the only thing in the room that was entirely his. He finished on the sound the way a man finishes a shift: with his hands empty and the next thing already waiting.

The amusement went out of it slowly, the way warmth goes out of a stone, not all at once but by degrees, and noticeable only because you had been holding it.

It was the not-coming that did it. He had made the joke about the blanket and waited, in the small animal way a body waits, for someone to come and put the blanket right; and nobody had, and he had understood, with the part of him that was forty years old, that this was not neglect. It was arrangement. The room had arranged itself around not coming.

Dessa had a scar she rubbed when she was deciding how much to tell, and she had decided. Alda answered questions with dates nobody had asked for, and the dates covered something. Mira had a bone in her pocket and a smile on her face and had stood in the one place in the room where the slate was not, and had told him, without saying a single word about him, that there was a secret in what happened to a child after it was born. None of them was saying it. All of them were standing around the shape of it, the way people stand around a pot that has cracked and are careful not to knock it.

He was the pot. He lay on the board-bed with a fold of wool under him and a cold stone where his belly had been and \lore[birth-flower]{a flower} written on his ribs, and the three women in the room were keeping something from him, and the keeping had a shape, and the shape was his. He had read a room like this before, a family around a bed, all of them agreeing without a word that the thing on the table would be called one thing and not another, and he had never once been the thing on the table. He had been the man who decided what to call it. The part of him that was a surgeon took that in and filed it, and the part of him that was two hours old wanted the breast and was not consulted about anything.

Mira's face was the worst of it, because it was the most arranged. The smile sat on it like a door on the latch, and behind the door there was a woman who had come in with eels on her sleeves and had said three things that mattered and made sure that none of them did. She knew more than she would say. She had arranged her whole face to keep it, and had done it so long that the face was the truth of her and the thing behind it was the rumour.

He went back to the count because it was the only thing in the room that belonged to nobody but him. Four on the near board. Four on the next. He had been counting exits without knowing it and now he knew a little, not enough, not nearly, but enough to know that a count was a thing nobody could take from him and nobody could enter on a slate.

On the third board the fifth thing waited. It had not moved. It did not shine in the tallow light the way the nails did; it was dark and old, and the grain curled round it in slow rings. He looked at it until the dip guttered and the light jumped and the knot was briefly somewhere else, and then the light came back and it was there again, a knot the width of a thumb where the fifth nail should have been.

\scenebreak
He had looked at the knot too long. The dip guttered and the light went off the third board and came back, and the knot was where it had been, dark and old with the grain going round it in slow rings. It would be there when the light was gone altogether. That was a comfort of a narrow kind, the only sort this room had to offer: some things kept their shape without being asked.

\person[Alda Reed]{Alda}'s hand moved. It left the fold of wool at the foot of the board-bed and went to the bench by the door, and his eye went with it, because his eye had always gone with hands, and because a hand going somewhere is a thing a man can read before the rest of the body has decided to follow. Hers went to the slate.

The slate lay where it had lain since he woke in the room, a flat grey slab propped against the wall within reach of her right hand, cold to the look of it, with a dusting of chalk along its lower edge. He had known there was writing on it without looking, the way a man knows a chart is lying open on a desk before he has read the name at the top of it. Now he looked.

There were two hands on the slate. The lower one was Alda's: thick, round letters pressed hard into the chalk, a thumb-mark dragging through the tail of one of them, a short column of figures that ended in a date. The upper one was not hers. It ran along the top of the slab in a dark ink that had gone into the grain and dried there, upright and small and neat, with no stroke wasted and nothing crossed out. He could not have told you a letter of it. \lore{Grathic} was a set of sounds that went past him while he lay under a wool blanket and did not come back; the written shape of it was a row of marks he had no key to. But he could read the shape of a hand the way he could read the shape of a stop, and this hand was strange to the room in a way he could not place, and he set it beside the other one and knew they had not been written by the same person or in the same hour.

\person{Daniel Marchetti} had once walked into a room in the hospital where he worked and seen the top line of a chart lying open on a desk that was not his desk, with a name on it he had no business knowing. He had taken the line in at a glance and looked away, and the looking away had not changed the fact that he had read it. It was the same now. He was not supposed to be able to read this, and he read the shape of it anyway: short strokes at the first, a pause, more short strokes, and the pause always in the same place when he let his eye go back to it.

The pause did something to him. It was not the writing. It was the rhythm underneath it, the small decided gap between one motion and the next, the way a hand stops to think and then goes on without hurrying, and he had seen that rhythm at the edge of a flare light in the snow. A woman in a coat the colour of wet ash had held a slate where this slate was and put short strokes on it, then a pause, then more strokes, and the hand had not stopped when the fire took the van. He had filed that hand away without a name and without a thought, because there had been a child under the dashboard and a lit rectangle in his coat pocket with a name on it he would not look at. Now the same pause sat on a slab of grey stone within reach of a midwife's hand in a clay room in \place{Low Kiln}, and the ink of it was dry.

He did not put the two together. He set them side by side, the way he set two films side by side on a light box in a room with no windows, and waited for a third thing to come and explain them, and no third thing came. The neat line was older than the chalk. Alda's marks sat on the surface of the slate and could be blown away; the dark line had been ground in, and whoever had written it had written before the chalk, before the labour, before anybody in this room had looked at the child's face and decided what to call him. That was the part that mattered, and it was the part he could not say, and the not-saying of it was already the shape of the room.

Alda blew the dust off her own column and set the slate down on the pallet edge, near enough to the child's head that he could see the dark line from where he lay. She did not show it to \person[Dessa Silt]{Dessa}. Dessa did not ask to see it. There had been an arrangement in the room about who looked at what, and the arrangement had been made without anyone's having to say a word about it.

Dessa's thumb went to the scar. “You can write what you like on that,” she said. “I'm not saying it.”

“Nobody asked you to,” Alda said.

“You were going to.”

“I was not.”

“You had your mouth open.” Dessa smiled --- the door shut, the smile her sister called the door shut --- and pressed the scar with the flat of her thumb. “I know four loaves. I baked four loaves in Grain, and three went on the levy, and one we ate with our feet up, and I know nineteen out and nineteen back, and I know what day it is, which is the seventh of Stone in eight hundred and twelve. I know all of that. I'm not saying the other thing.”

“A name's a thing a person says when they mean to keep something,” she said, and the smile did not move. “I don't mean to keep anything. That's all the speech there is.”

Dan understood the shape of it without the words, which was the curse of understanding shapes. She had said the year and the loaves and the steps because a year and a loaf and a step were things a person could hold, and the thing she would not hold was the child. She had called him the little hot one and the flower for two hours and would go on calling him that for as long as he lay in this room, because a thing with no name cannot be claimed by a woman with no house, and a name entered in a neat dark hand on a slate was a claim made already by somebody who had not asked her. He could not have said any of that in Grathic. He said none of it in any language, because the body had no breath for a sentence and no lips shaped to one. He lay and took it in, and the taking-in was the only act available to him, and it was not nothing.

Her pulse had gone up while she said the loaves. He could see it at the wrist where her hand hung over the edge of the pallet, a quick light tapping that had not been there when she counted the steps. A man who could not lift his own hand could still read that. It was the same hand he had read an hour ago and the same pale scar, and the change was in the rate and nowhere else, which is the way a body tells the truth when the mouth has been arranged not to. He put the rate away with the rest of it. He had put away a thousand rates in his life and had never once been able to do anything with the ones that mattered.

\person[Mira Silt]{Mira} had been standing in the one patch of floor the slate did not reach, with her back to the writing and her face to the door. She did not look at the slate now. She put one hand into the fold of her sleeve and brought it out with a crust of brown-edged barley bread wrapped in a rag, and put it on the edge of Dessa's pallet with two quick movements, before anyone could thank her or make a speech about it or ask where it had come from. The movement was brisk and small and finished, the way a woman closes a door she does not want to stand talking at.

“Eat it,” she said.

Dessa did not pick it up. “Where'd that come from?”

“There she goes.” Mira's thumb went into her pocket and came out again, and the eel-bone did not tick, because she had put it away for the duration of this small thing. “You want the whole story. Don't have one. It came from a warm house, that's all. There's barley in it and a little smoke and it's a crust, not a sermon. Eat it before it goes cold, which is the only thing in this room that's in a hurry.”

The smoke came off the crust when her sleeve moved, and Dan smelled it from where he lay: woodsmoke, not tallow, the smoke of a kiln banked through a cold night. It was the same smoke he had smelled an hour ago on the rag that had come round the warmed stone. Warmth had come into this room twice in one night in the same wrapping --- once as a stone for a child's belly and once as bread for a child's mother --- and both times from the kiln-house down the yard, and neither time had anybody said whose fire paid for it. He could name the house because he had been told where the stone was warmed. He could not have named the hand that broke the bread, or the reason the same house should pay twice for the same cold family. He filed the loaf beside the stone and the stone beside the writing and went on lying there with a head he could not hold up.

When Dessa still did not take the crust, Mira reached over and pressed it into her sister's hand and shut the fingers on it, and did it quickly, and took her hand back as though it had been burned. It was a kindness done at speed, before either of them could turn it into something that would have to be said. She went back to her patch of floor and kept her back to the slate. She had crossed the room to give bread and had stopped short of the writing to do it, and she had not once asked to see what was written there, and the not-asking was the loudest thing about her.

Alda watched all of it and said nothing about the bread, which was its own kind of answer. Then the room's question, which no one had asked out loud, went to her because it always went to her, and she answered it the way she answered everything: with numbers nobody had requested.

“Stone seven,” she said. “The seventh of the twelfth, in eight hundred and twelve. The ice is holding; they say at \place{the Weir} it'll hold three days, though they said the same in Bone and were wrong by a week. Snow before \lore{the Five}, likely. It's a cold coming, and a child with a stone at his belly wants a warm room more than he wants a name.” She set the chalk down on the slab. “Warm and breathing. That's all I've set down.”

Her hand went to the little tin at her bodice while she said it, a touch as brief as a man's touch on a door to be sure it is shut, and came back to the slate. He had seen that touch once already while the labour was still going on, and had known even then that it was a thing done every time and never for the first time. There was a second record in the room. It was not on the slate. The slate was the record a person could ask to see, and she had answered the unasked question with a date and a weather report and the temperature of a child, and the tin had not moved. What was in the tin he could not know, and he did not try to guess; he had learned in a white room that the note under the chart is the thing that matters, and that the note is never the thing anybody reads aloud.

In Low Kiln the weather was never the subject. It was the thing said in place of one.

\person{Corin} took the room apart the way he had once taken a ward apart at four in the morning, reading who stood where and who had not sat down. Alda had written in two hands on one slate, one of them not her own, and had touched the tin when the questions came near it. Dessa had named a year and four loaves and nineteen steps and would not name the child lying three feet from her. Mira had crossed the floor to give bread and had stopped on the far side of the writing to do it, and had kept her back to the ink the whole time. Nobody had asked him anything. Nobody had looked at his face and then written; they had written first and the face had arrived afterwards underneath the words, which was a different order of things and one he recognized from the other side of a desk. He had spent twelve years entering people into charts. He had never once been entered. Now he had been entered twice: once in the neat dark hand he could not read, and once in the counting he kept for himself, which was the only entry in this room that belonged to nobody.

The room had arranged itself around the not-saying, the way a yard leaves a clear run to the gate for a cart that has not come. The silence had a shape, and the shape had edges: a slate turned so that one woman could keep her back to it, a crust given before a speech could be made, a date put down where a name should have been, two hands on one stone slab with a gap between them the width of a life. The arrangement was not malice. That was the part that would not let him rest. If it had been malice he would have known what to do with it; he had worked in a building full of people who arranged their silences around a chart and called the arrangement policy. This was different. This was three women keeping a thing warm between them the way they keep the fire banked, with care and with frequency and without once asking the thing whether it wanted to be kept.

Under him the stone finished. It had given what it had and then had lain for a while at the temperature of the room, taking nothing and giving nothing, an ordinary stone in a fold of damp wool. Now it went past even that. The cold came up through the boards and the thin wool and took it, and it became a stone that was cold in its own right, with no warmth left in it to spend or to withhold, and he lay against it and was the only person in the room who had been keeping the account. There was no moment to mark and no sound to mark it with. A stone does not announce the hour at which it stops being a kindness. It is one, and then it is not, and the child it was warming is the only one who notices, and the child cannot speak.

Along the skirting board the grey ridge lay where it had lain for years, the width of a thumb, sifted out of the walls and the pots and the kept lives of the house. The Five were three days off. Nobody would sweep the ridge before \lore{the Last Day}; it would be left where it was and put into the fire the next morning, which is a way of dealing with a thing by moving it somewhere it can no longer be read. He had known a building that kept its worst charts where anybody could see them and then burnt them, and called the burning an accident.

And there was a third day in the room. It had not happened yet, and the room was full of its shape: a coming, a hand, a decision made already in the order of the reaching and the not-sitting and the not-saying. Somebody would come on the third day. He did not know who. The room did not say, and the slate with the neat dark line on it did not say, and the woman at the edge of the flare light had not looked up from her writing long enough to say it either.

The dip burned lower and the circle of light drew in around the pallet, leaving the corners to \lore{the Ash}. He went back to the count because it was the only entry in the room that nobody could take from him, and the third day had not come. Then he stopped counting, because there was something to look at instead.

Alda's hand lay on the edge of the slate where she had set it by the child's head, and the other hand --- the neat one, the one in dark ink --- had dried into the grain of the slate beside it. In the fold of the grey wool the small stone lay cold. Warmth spent. Ink dry.

\scenebreak
The dip went out while he was looking at it. There was no struggle in it, no last brightening: the flame drew down into the wick, the circle of light on the boards closed like a hand, and the dark came back for a moment as though it had only been waiting for the room to turn its head. Then the dark thinned. It did not lift; it lessened, the way a pot lessens on a slow fire, leaving the strength of it at the bottom. The corners came back grey. Clay dust showed on the boards where no light had been an hour ago, and the ridge along the skirting board took on a colour at last, the colour of ash that has sat in a corner since before the child was thought of.

He had been awake the whole of the night and nobody had asked him once whether he wanted to be.

\person{Daniel Marchetti} had worked nights in a building where the lights never went out, and had learned to read the hour by the faces of the people coming on shift. Here the hour was told by the light on the floor. He lay on his back under a wool blanket that would not stay on his shoulder and took the morning in the only way open to him: as a change in temperature, a weakening of the dark, a small cold weight against his belly where the stone had been and still was, only now with nothing in it. The stone had given what it had. It lay in its fold of wool with the indifference of a thing that has been spent and is no longer answerable.

\person{Corin} shifted one hand. The hand went where it was told, which was nowhere. He could not turn his head far enough to see the door, and he had already learned that turning it produced a small animal sound from him that made \person[Dessa Silt]{Dessa} look up, so he stopped trying. A trauma surgeon who cannot turn his own head. The joke had worn out in the first hour and the fact had not.

Something came into the doorway before the woman did. It came as a smell: woodsmoke first, and wet wool under it, and beneath both that faint sour mineral edge he had learned to place along the skirting board. Then the cold off \place[the Clay Yard]{the yard} came in with it, and then the hand.

The hand came across the room ahead of the body the way a hand goes ahead of a person in a dark passage, feeling for the edge of a thing. It was a broad hand with forearms the colour of dried slip, and on the back of it, across the knuckles and down toward the wrist, was a burn shaped like a little door: two straight sides, a lintel, a threshold, the skin inside it puckered and pale. He had seen fire write on skin before, in a country with different lights, but most burns are careless of shape. This one had a lintel. The fingers were half closed, not on anything, but closed the way a hand closes when it is used to carrying a cup and has decided to carry something else instead.

The hand stopped at the fold of wool at his belly. It rested there. He felt the pressure through the blanket and the wool and the stone, a careful pressure, the pressure of someone feeling for heat and not finding it. The fingers changed at the knuckles, a small tightening, and went still. He understood the gesture without being able to place it in the room's order of things: this person had come to see whether the warmth had held. It had not.

Then the hand moved up to the covered bundle and stopped above it.

Above, not on it.

The fingers opened and shut once, small, measuring a weight not yet lifted. He knew that movement. He had made it himself a thousand times, standing over a stretcher in a white room, estimating mass and grip and what the thing under the cloth would do when it was moved. He knew the difference between a hand that has come to take and a hand that has come to check. This hand had not decided which it was.

It did not take him.

The room held still around that. \person[Mira Silt]{Mira} was standing in her patch of floor with her back to where the slate had been and her face half toward the door, and she did not move. The small smile sat on her face as it had sat all night, on the latch, and her thumb went into her pocket and came out without the eel-bone, because this was not a moment for a sound. Dessa's thumb went to the pale scar on her forearm and stayed there. \person[Alda Reed]{Alda}'s hand lay on the bench where she had set the slate, and none of the three looked at the woman in the door, and the woman in the door said nothing.

That was the whole of the pause: four people in a room not looking at one another, and a hand waiting above a bundle.

Dan read it the way he had read a trauma bay at four in the morning, when the senior stops giving orders, the nurse stops putting instruments down, and everybody in the room knows that a decision has been made somewhere above them and will now come down. He could see the shape of the decision. A foundling with no house-name is a coat left on a peg in a poor house: anybody can wear it, anybody can take it, and nobody is obliged to say who left it there. Three women had arranged the room around that fact for two hours without a word spoken about it, and now a fourth stood in the door with her hand out, and the arrangement was being put to the test. What would it cost to keep a boy who had no name anyone would speak? Who would pay for the warmth? These were the questions in the room. Nobody asked them. Nobody had to.

He could read the hesitation as a physical thing, the fingers open and shut, the hand not closing on the wool, and he could not read the reason behind it, and the not-reading was like a chart turned face down on a desk he was not supposed to be standing at. He had spent twelve years being the person who decided which patients went to which bed. Now he was the bed. He was the thing being sorted, and he could not so much as lift his own hand to take part in it. He had no words in the language of this house for what the house was doing, but he knew the thing: a household that keeps its own roll and will not show it to a man come down from \place{Slate Row} with a slate in his hand.

Under the blanket he wanted the breast. It was a clean, small, stupid want, and it was louder than the hand and louder than the slate and louder than the smell of woodsmoke coming in off the yard. He had once read a man bleeding out and wanted a cup of water. The body keeps its own ledger and does not care what the mind enters in another. He drooled. A line of it went down into the wool at his neck and was not wiped, because no one in the room was looking at his face. He had reduced a thousand shoulders. He could not wipe his own chin.

“Cold in here,” Mira said. It was not an answer to anything anyone had said, because nobody had said anything. “Cold enough to keep.” She did not say what. She never said what. The eel-bone went tick once against her thumb and then she put it away, and the small smile stayed where it was, shut and harmless.

Dessa's pulse had come up at her wrist. He could see it from where he lay, the quick light tapping in a vessel that had been quiet an hour ago. She had counted nineteen steps out and nineteen back for six years and had counted them tonight for a reason that was not the pits, and now the counting had stopped and the pulse had taken it up. She smiled. It was the smile her sister called the door shut, and it meant that nothing further would be said, and the hand in the doorway went on hovering above her child.

Alda's fingers moved on the bench. She did not look up. In all the births she had attended, in a district where a woman with a question was a woman who wanted an answer written down, she had learned that the correct number of words in a room like this was none, and she kept to it. Her hand was steady. He knew from the first hour that steady hands were not the same thing as a woman who was not afraid.

The kiln had not breathed yet. The room smelled of its smoke without giving its sound, the way a house smells of a fire in the next street before anyone has seen the flame.

\person{Sella Vail} said nothing. She had not come into the room. The hand had come into the room and the smell had come into the room and the shadow of her stood in the door, and the woman herself had held back, as though she understood that a body in a room is a thing that must be answered for, and a hand is not.

The hand closed again. The little door of the burn shut on nothing, and the fingers went back to being a hand at her side, and the woman who owned it stepped back into the doorway and then out of it. She did not look at the child. She did not look at any of the women. Out in the Clay Yard her feet moved on clay dust, and the woodsmoke moved with her, and then both were gone from the room.

Alda's hand went to the slate. She laid it flat on the bench with the neat dark line underneath and the chalk underneath and her own thumb-mark underneath, so that the room was a room again with no writing in it, and he could no longer see the shape of the hand he could not read. The slab made a small flat sound on the wood. It was the only sound anyone in the room had made in a long time.

Then, somewhere off the yard, a kiln drew breath.

It came through the wall more than through the air: a low intake, a pause, and a long letting-out that went on longer than a breath has any right to, as though the thing had been holding it since the middle of the night and was only now permitted to spend it. Woodsmoke went with it. He knew the smoke. It was the smoke that had come off the crust of brown-edged barley an hour ago, and the smoke that had been in the rag round the warmed stone, and it had crossed this room twice already without once saying whose fire had paid for it. He could name the house that kept the fire. He could not name the hand that had stood in the door, and in this district those were two different things and the distance between them was where the whole arrangement lived.

The kiln breathed again. He counted it.

He counted it the way he counted the boards, because it was there and because there was nothing else to count: four on the near board, four on the next, and a fifth that was not a nail but a knot the width of a thumb. Four breaths, four breaths, and then a fifth that was not a breath at all but the small space where the next one should have been, a gap the width of a thumb in the order of things. The fifth never came. The room went on breathing around it. He had counted pulses in a hundred bodies and exits in a hundred rooms and nails in these boards since he woke, and this was the only count he had ever kept that belonged to nobody but him. The slate had a name in it that someone else had written. The stone had been warmed at a fire someone else kept. The count was his. He kept it the way a man keeps the one thing in a building that is not entered against him.

The light came on further. Morning got into the yard and the yard got into the room, and the ridge of \lore{the Ash} along the skirting board showed its whole length at last, from the bench to the door and under the door into the next room, where it would be left where it was until \lore{the Last Day} and then swept into a fire and called dealt with. \lore{The Five} were three days off. The ice on \place{the Ilm} was holding; they said at \place{the Weir} it would hold three days, and in \place{Low Kiln} a man says three days when he has counted what he has and does not want to say it. Somebody would come on the third day. The room had said it without a word, the way it had said everything else.

He went back to the stone. It was cold in its own right now, with nothing left to give or withhold, an ordinary river stone in a fold of damp wool against the belly of a child who could not have told anyone it had once been warm. He had been the only person in the room keeping that account, and the account was finished, and he had not been asked and had not been paid and had never once been given the number he wanted.

The hand had hovered. The name had been written in a hand he could not read. He did not know that either of them would come back for him; he knew only that the third day had not come, and that it was coming, and that until it did there was the cold stone at his belly and the four-and-four-and-a-knot of the floorboards under him, which nobody in the room could see and nobody in the room would have known to look for.

\CRtailpiece
\begin{CRchronicle}{1}
In the birth room off the Clay Yard in Low Kiln, a stone warmed at a neighbour’s fire lay against a newborn’s belly, already giving up its heat; nobody felt it go. Alda Reed worked with steady hands, clay under her nails, reaching for her slate before she reached for the child. “Warm the stone first,” she said. “A cold stone is a stone, not a kindness.” Dessa Silt lay in the aftermath, face to the wall, counting the nineteen steps she had walked from slip pit to rack for six years, rubbing the pale seam of a sieve-wire scar.

Corin was two hours old and carried the whole memory of Daniel Marchetti, a trauma surgeon who died at 2:17 on a Chicago expressway, going back into a burning minivan for a child under the dashboard. His warmth left him in the giving and did not cool. The wipers had stopped, the fuel hissed, a phone lay lit in his pocket with a name he would not look at; at the edge of the flare a woman in a grey coat wrote on a slate, paused, wrote again, and did not stop when the fire took the van. \emph{I have left the job unfinished.}

Alda lifted him and saw the branching red mark across his ribs and down the left forearm. She crossed two fingers over it, caught herself, and covered the act with a word and an unasked-for date from Nine Pots: birth-flower. Dessa would say no name. She called him the little hot one and pressed the wool flat at his shoulder with fingers that stayed longer than the cloth required.

On the slate a name stood in a neat dark hand, written before anyone looked at his face, in characters he had no key to --- though the pause between strokes was one he knew from the flare. Alda’s chalk sat beneath it; twice she touched the tin at her bodice. Mira Silt came in with eel slime on her sleeves, an eel-bone ticking. She would not stand where the slate lay open, answered a question with a question, and gave Dessa a crust of barley bread briskly. She knew more than she said and had arranged her face to keep it.

The surgeon took the room apart and could do nothing. He knew the afterdrop and wanted the breast, and the wanting was louder. He drooled; he could not hold up his head or reach a cup. He counted four nails, four, and a knot where the fifth should be: the way out, though he did not know it. Toward morning a hand with a burn shaped like a little door felt for heat at his belly, hovered over the bundle, and did not take him. The slate shut; a kiln breathed off the yard. The Ilm lay under ice, the Five three days off, a coming on the third day. He lay under the spent stone with a count nobody could keep.
\end{CRchronicle}

% Generated by forge/build.py from meta.json and the author's scene files.
% Edit the scene files, not this file.
\CRplate{build/art/chapters/ch-0002.jpg}{Days 2–3. The Vail kiln-house takes in a child.}
\CRdecoration{book/chapters/ch-0002/decoration.tex}
\CRchapter{2}{He’s Not Vail}{Days 2–3 · Stone 8 and Stone 9, Y.H. 812, an Eighthday and a Ninthday · the Ilm under ice and the Vail kiln breathing through the wall}{The birth room of House Silt and the Vail kiln-house, wheel-shed and Clay Yard in Low Kiln, Sill, Grath; one remembered car park in Chicago}
\CRepigraph{The foundling is a coat on a peg: anybody may wear it, nobody must say who left it, and the cold comes in at the seam nobody mended.}{A saying of the Clay Yard, Low Kiln}
\daymark{2}{Stone 8, Y.H. 812 (Eighthday)}

\dropcap{T}{he kiln} had been breathing since the middle of the night. In the \place{Vail kiln-house} at first light a man waking in the next room would have known the place by that sound and by nothing else: a low intake drawn through a throat of brick, a pause as though the thing had forgotten what it was about, and then a letting-out so long it went on past the point where a breath had any business ending. It was the breath of the fire. It had been kept by coals banked under ash for two years, and it would go on after the house had stopped speaking and gone to its work. People who lived beside it had long since ceased to hear it, the way a man ceases to hear his own breathing until it fails.

It came in fours. Four draws, four draws, and then a space where the fifth should have been, a gap the width of a thumb in the order of the thing, and then the long letting-out that filled the room with smoke and went on until you thought the coals had given all they had and found another breath in them. \person{Tarn}, who had slept through the river icing over and the house growing cold around him, had never counted it. \person{Sella Vail}, who had kept the fire through two winters without missing a night, had never needed to. It was simply how the house breathed, and a man might live beside it all his days without hearing the missing fifth, and then one morning the missing fifth would be all he could hear.

You could have known the house without ever seeing a face in it. In the \place{Vail wheel-shed} the wheel stood with its seat worn hollow on the left, a shallow bowl in the wood where a man had sat for twenty years and leaned left into the throw, and the hollow was empty now and would stay empty; Sella had twice stopped a hand that went to fill it. Behind the door a coat on a peg kept the shape of a broad man, salt stiff in the seams from a winter of carrying slip, and the shape was broader than any man who lived in the house now. \person{Jorra Vail} had reached for it in the dark twice and stopped both times and said nothing about either. Along the skirting board ran a ridge of \lore{the Ash}, the width of a thumb and thicker than the same ridge in any other house in \place{Low Kiln}, sour in the nose when the room was warm, plain as a line struck in chalk when it was cold. Nobody swept it. It would be swept on \lore{the Last Day}, with the rest of the district's sweeping, and burned next morning and called dealt with, which was how Low Kiln dealt with most things it did not want to look at.

When the kiln was low and the house had gone quiet, the smell of him was stronger than any memory of him. It came out of the corner where the Ash lay and out of the coat on its peg and out of the hollow in the seat: wet clay and banked coals and the sweetish edge of lead glaze, gone thin and cold, the way a room smells where something has been left too long and nobody has moved it. Jorra had been thirteen when his father went and had taken the wheel a week later, and had spent three years trying to fill the hollow with his own weight, and had not once said that he could not. The wheel did not care whose hands worked it. It cared if they were steady.

The sayings of \person{Garl Vail} had outlived him and improved in the keeping, as sayings do when nobody remembers the tone they were said in. The house kept one of them and said it over the pot and the fire and on mornings when none of the three people in it had slept. “Leave the crack showing. People trust a thing they can see has been handled.” He had been a big man with clay in the creases of his hands and a laugh that ended in a cough, and had gone in the Long Cold when there was no crossing the river for a collector, and had not been collected. His warmth lay where it had fallen. The ridge along the skirting board was the shape of it; the coat kept his shape on the peg; the seat kept the leftward lean of him in the wood. None of these facts had been entered in any account a man come down from \place{Slate Row} with a slate in his hand could read.

The house kept a roll of its own, and had kept it since Garl was a boy, and did not show it to that man either. On its private slate Garl was entered and Enna was entered, and the foundling's name, when he came, would be entered in a hand that was careful about such things. It was a small thing kept close, and whether it was a true roll or a comfort nobody in the district had ever asked, and the house had never offered to say. A household that keeps its own account does not need to explain the keeping; it needs only to keep it.

Tarn had been sweeping since before anyone woke. It was not in the order of the house for a clay-bound hand to be up before the fire was stirred, but Tarn had never needed ordering in that particular, having learned early that a broom in his hand was a thing nobody could call him for. Six years of a seven-year term lay behind him and he had counted every one, and would not, he said, count the next. He called Sella “house,” the way he called everyone house, and said “I owe” where another man said “I am.” He swept the clay dust toward the door and stopped. A pot had been set down somewhere behind him and the small sound of it went through the room and into his sweeping, and for the space of two strokes the broom did not move. Then it moved again. He had not turned his head. Nobody in the room had said anything about the pot.

When he thought nobody was looking he put the small bench knife into his pocket. It was a stub of a thing with the mark on its handle worn almost flat from use, and it went in quickly and without ceremony and was gone. Then he said, of his own accord and to nobody in particular, that he had taken the small knife and put it back and they could check if they liked. He said it in the tone of a man stating the weather, and the tone invited a laugh that did not come. His deference had a nail in it: he agreed so thoroughly that an argument had nowhere to go, and the thoroughness of it was how you knew where he stood. Sella had counted the tools before she opened her eyes that morning and did not look up. The house needed no telling.

Jorra stood at the wheel in the shed with the door rolled back and had been there since the coals were stirred. He was sixteen and built like a hinge, tall for a clay boy and narrow, with slip dried to the colour of old bone up both forearms and a permanent bend in the left little finger. He told people the finger came from a fight, because a fight sounded more like him; it had come from the wheel when he was eleven. He did not look at Sella when she spoke to him. He never had. He obeyed her instantly, which was the other half of the same habit: a son who cannot bear to look at a command and cannot think of disobeying one. He spoke to the wheel first and to the pots after it and to people last, and to people in short flat declaratives that stopped where a longer sentence would have had to go on.

“Is it dry?” he asked the pot. To the wheel he gave the little click of the foot that started it, and to Sella, when the work required it, he said, “The slip's too thick. It'll crack at the shoulder.” He had a way of going quiet when \lore{the levy} was spoken of that the household had learned to read as weather, and there was a temper under the quiet that nobody had yet seen the bottom of. Men had come into the back room through the winter and spoken low about slates and accounts, and Jorra had listened to all of them and given none of them a word they could carry away. When they pressed him he said he had not said he was with them and had not said he was against them, and that in this district those were the same thing and he had said one of them anyway. He would not say which.

Sella kept the house the way a widow keeps a wound clean: daily, without comment, with a hand that knew exactly how much the thing would bear. She was not tall and did not make herself tall; she took up room by being impossible to move. Before she opened the kiln she put her thumb to the cord at her throat and to the \lore{widow's mark} hung there, a thumb-length of lead-glazed clay with the Vail kiln scratched into it, and held it for the count of a breath, and let it go. She touched it at no other time. It was not a thing she did when she was afraid; she did it before the fire, and there were people in the district who said the fire would go out if she ever forgot. She would have let them think it. When she was angry she got quieter, and the quieter she got the more the house straightened itself up around her, as it was straightening now.

“Feed it,” she said to the fire, and the fire was fed. She said “the house” where another woman said “I,” and gave instruction to the clay and the coals and never a finished feeling, because a thing said aloud was a thing the Court could enter, and she had learned to keep her feelings where no Court's hand could reach them. She would not let anyone fill Garl's seat or take his coat down from the peg, and had never given a reason for either, and the household had stopped asking for one years ago.

There was a foundling in the district with no house-name, three doors down the yard in \place{House Silt}, and a slate in that room with a name in a hand nobody there could read. The boy was to be brought under this roof on the morrow, the third day, and nobody in the kiln-house had said what would be paid for him or why, and nobody had said his name. It had been decided without a meeting, the way such things were decided in Low Kiln: a coat left on a peg in a poor house is a coat anybody may wear, and who left it is a question nobody is obliged to answer. A stone warmed at this fire, the colour of skimmed milk, had already crossed the yard in an apron and lay cold now under a fold of wool in the room three doors down, and a crust of brown-edged barley carrying the smoke of this kiln had gone the same way. Neither had been entered against anybody. The house that had paid for both had not once said so.

Sella stood with her hand on the kiln door and said the thing that had been waiting in the room all morning. “The boy comes under the roof on the morrow,” she said. “The house will keep him.”

Jorra did not look at her. His hands went on turning the pot, the left little finger bent and steady, and the wheel went on with its low small sound under his palms.

“He's not Vail,” he said.

It was not loud. It was not an argument in the way arguments are usually made, with the voice up and the face turned; it was the flat declarative of a boy who had put a thing where it could not be taken back, and who knew that in this house a thing said was a thing entered. Sella did not answer it. She had been quiet before and she got quieter now, and the quiet went into the kiln-house like cold into a room, and Tarn's broom found the dust and swept it, and the wheel went on turning.

In the ash where the fire had been cleaned the night before, the charred edge of a slate lay among the coals, pushed in at some time and never raked out, the writing gone and the edge gone black. Nobody had picked it up to look. The house had swept around it for a year without once asking what it was, and had gone on keeping its own account.

She took her hand from her throat and put it on the kiln door and opened it. Heat came up out of the fire and took a degree of chill from the room, and woodsmoke went into the wet wool of the coats and stayed. “Don't set that near the door,” she said to the pot. “You'll pay for it in the thaw.” She gave the fire what it needed and said nothing further, and the house went on.

Three doors down the yard, in a room with a door that would not shut, a child lay under a fold of grey wool with a spent stone at his belly and did not know he had been argued about. He had not yet come under this roof. The third day was still a day off. \place[the Ilm]{The Ilm} lay under its skin of ice, and \lore{the Five} were spoken of as three days off, and Low Kiln said three days when it had counted what it had and did not want to say it. In the kiln-house the wheel went on turning and Tarn swept and Jorra did not look up, and the words stood where he had put them, in the space between a mother and her son where the foundling would lie when he came.

\scenebreak
The door had given up on being shut in the small hours and kept a finger's width of \place{the Clay Yard} in sight all morning, a grey slit where the sleet came in and ticked on the oiled paper of the window and went nowhere. \person{Corin} lay on the board-bed with the damp grey wool under the back of his head and gone from his shoulder again, gone a dozen times since light, and the stone at his belly had cooled past the point where he could tell where it left off and his own skin began. He was two days old and had been awake most of them. \person{Daniel Marchetti} was still here, in the only place that could hold him, which was behind these eyes; he had been forty and had worked nights and had known the weight of a body coming through a door in the dark, and now he could not feel his own left foot. There was a surgeon's exactness in the way he took the measure of all this and an infant's helplessness in the way he could do nothing about it, and the two sat together in him, the man and the small animal he had been folded into, arguing without words.

When he turned his head past a point the sound came up in his throat, a thin animal noise that was not a cry, the noise a thing makes when it is handled and cannot help itself, and he had learned to stop short of that point and lie still, which was a discipline he had not known he still had. He had reduced a thousand shoulders in a room that smelled of sweat and iodine and somebody's bad coffee, and he could not keep a wool blanket on one infant shoulder. The joke had a blade in it. He turned it gently, the way a man tests a tooth with his tongue, and it cost him the same as always.

The count was still there. Four on the near board, four on the next, and a fifth that was not a nail but a knot the width of a thumb, a small dark place where the order of the boards went wrong and nobody had thought to mend it. He had kept it through the night because there was nothing else to keep, and the knot, not the nails, was what he kept returning to. Through the wall, very far off, the Vail fire took a breath and let it out, and he found he was listening for the gap in it before he had decided to listen.

\place{House Silt} held the cold the way a cracked cup holds water, with care and with one thing nobody would mention until it was too late. The ridge of \lore{the Ash} ran along the skirting board from the bench by the door, the width of a thumb, and would lie there until \lore{the Last Day} and be swept then and burned next morning and called dealt with. Damp wool, tallow, the sour edge of clay dust: these were the room's smells, and under them the faint warm milk and scrubbed soap of \person{Alda Reed}, who had been working since before he woke. On the pallet \person{Dessa Silt} breathed the long shallow breathing of a woman who had spent a night and a day and most of another on her feet before she had spent it lying down, and her pulse had come up again now that there was nothing for it to do.

A cup was set down on the bench across the room, and the small flat sound of it went through the grey slit of the door and into the room and into him.

He knew the note before he knew that he knew it. The cup had given back not \lore{the full note} but its flat shadow: a small dead sound with no body behind it, a sound that stopped at the surface of the thing and would not go through, as though the clay had answered and then thought better of it. The full note was what a vessel made when the whole of it agreed, a bright knock that carried through the bottom and came up clean under the knuckle; this was a vessel that did not answer full. He had no \lore{Grathic} word for either and had not needed one. He knew it the way he knew a pulse under his fingers before he had counted it, the way a man hears a wheel out of true on a stone, and the knowing was a tool with nothing in the world for it to touch.

The cup sat two paces off on the bench. He could see the thin light on its rim through the grey slit of the door. He could not reach it. He could not lift a hand to it, could not make a sound that would mean tap it, could not have told anyone in the room what he had heard if the room had had a word for hearing it. The competence had found its object at last and sat in an infant's body with a drool on its chin, prisoner of two paces of trodden ash. He listened for the cup to be set down again. Nobody set it down again.

He had once been able to stand at the foot of a bed and tell, from the sound a chest made when a man breathed in, which side of it had gone quiet. That was not a gift. It was four thousand repetitions and the habit of attending to the thing nobody else in the room was attending to, and here it was again in a body that could not turn over, running on a cup two paces off. He had the tool. He had never in his life had less use for one.

\scenebreak

Then \place{Chicago} came up under him, unsummoned and untagged. A car park in the city went wet with slush under the lamps and the dashboard clock read something past three; the pizza had gone cold in a paper bag on the console. \person{Deshawn Carter} sat in the passenger seat, a big man comfortable in his own skin, with a gold tooth he had earned in a joke and kept in affection, and quick hands that could start a line in the back of a rig on a road made of ice, and he says, \memory{You look like shit, Marchetti. Sit down. Ten minutes.} The rig smells of diesel and the lemon disinfectant the medic has wiped it down with, and cheap coffee, and the ten minutes stretch out and say nothing, which is the whole of the deal, and the bad pizza is eaten and the silence kept and the radio left on low. The hospital doors stand lit a little way off across the wet asphalt, and neither man looks at them, which is the whole reason for coming out here. He laughs. It comes whole out of him, from somewhere low in his chest, unembarrassed, the kind of laugh that makes a room turn and look and then smile despite itself. It goes on. In the car park, with the slush and the lemon smell and the gold tooth catching the lamp, it goes on and fills the space where the eleventh minute would have been; and then the radio, later, after the expressway and the flare light and the snow going orange, that same laugh running on the dispatcher's channel for exactly one minute, timed by a man who had nothing else to give, and the radio put down when the minute was done.

Then it was gone. The car park went out like a lamp put out, the gold tooth, the lemon and diesel, the ten minutes and the one minute both, and the birth room came back with its sleet and its wool and the cold stone. The laugh had been a sound he could not reach, and now there was no sound to reach and no hand to reach with. There was no word in this country for a man who rode in the back of a rig and made jokes instead of asking, and he had no word for him either, and the ache of it went in under the ribs where the red flower was and stayed. The surgeon wanted the cup. The infant wanted the breast. Both were true and neither was reachable, and the second was louder, and the first did not therefore stop.

\scenebreak

Alda came to him in her order. She reached for the slate before she reached for the child: a hand on the flat slab on the bench, a thumb along the edge where the writing was turned under, a small adjusting push that made sure no dark stroke showed. Only then the cloth and the warm water and the child. She worked in the order she had used four hundred times, face, neck, groin, knees, cord, with a stop at each place that was not hesitation but the exact amount of time the thing required. The water was warmer than he had expected, and when she lifted him the air went across his wet skin and he wanted to cry at the cold of it, and did, a thin animal noise, and she said “There now” without once stopping the order. He knew it. He had seen it the day before. He could have recited it, could have told her what came next, could not have lifted a finger to do any of it. She smelled of soap and tallow and the mineral edge of the Ash she had promised herself she would sweep tomorrow. Under the bodice of her grey shawl the little tin box on its cord shifted once as she reached across him, and her hand went to it for the count of a breath and came back to his neck.

“Hold still,” she said to him. “There's a knack to being born and you've missed it.” To the cloth she said, “Don't set that near the fire. You'll pay for it in the thaw,” though nothing was near any fire, and to the mother on the pallet, without turning, “Seventh of Stone, eight hundred and twelve. Sleet on the paper this morning, same as yesterday. It'll go off by the second firing. There was a potter's third child in \place{Nine Pots}, Y.H. eight hundred and six, fourth day of Stone, and he's a carter in \place{Stonemouth} now, so there's no call to fuss.”

The pulse on the pallet answered, a small rise and then a settling, the way it had the day before when the counting stopped. Dessa did not ask what had not been asked aloud. In this room the question sat in the space where nobody put it, and the answer Alda gave it was a date and the weather, which in \place{Low Kiln} was never the subject.

“Warm and breathing,” Alda said. “That's all I've set down.”

It was not all she had set down, and the child under her hands knew it the way he knew the cup, without a word and without a way to say it. She set the line of the cord straight, drew the wool up over his shoulder and pressed it flat with two fingers that stayed a moment longer than the cloth required; and as she did it the two fingers went to the corner of the branching mark and crossed themselves there, the small sign made at the corner of a \lore{birth-flower}, and she did not know she had made it in front of him. She covered it the way she covered everything, with the word for it and a date nobody had asked for, and went on with the washing.

The stone at his belly had not always been cold. It had been warmed at a fire across the yard and carried here in the fold of an apron by the woman on the pallet, and set against him while he was still too new to know that a thing could be done for him and never entered anywhere; and now it was cold and nobody in the room rose to warm it. That, too, was an account, and he was the only one in the room who could read the shape of it, and he could not speak.

The slate lay flat on the bench within reach of her right hand. Written on it, turned under, was a name in a neat dark hand that was not hers, a hand with a pause in it between strokes, a place where the writer had stopped as though the thing being written had needed a moment's thought before it could be a name. The child had registered that pause the day before and had set it beside another pause, the pause in a hand at the edge of flare light in a country with no name here, a coat the colour of wet ash and a slate and a hand that went on writing but stopped in the same place, and he had set the two side by side and had not joined them. He had no key to the strokes. He could not have read a word of it. He knew the shape of a hesitation the way he knew the shape of a wound, and the shape was the same, and the two facts lay in him unjoined, as they had lain since the night before.

Alda finished the cord, put the cloth down and wiped her hands on her apron, and the room went back to its order: the sleet, the door's grey slit, the pallet breathing, the cup across the room. He turned his head the allowed distance and no farther. The flat note was still in his ear as though the cup had gone on giving it after the sound had stopped, a small dead note with no body behind it, the sound of a vessel that did not answer full. Two paces off the cup sat on the bench with its rim catching the light. His hand went toward it --- the small hand, the milky nails, the hand that had held a clamp in another life and could not now close on air --- and it moved the distance an infant's hand can move and went nowhere, and the note stayed in his ear, and the door kept its finger's width of the yard.

\scenebreak
The afternoon came on without arriving, the way cold comes into a house that has one fire and that fire three doors away: by degrees nobody marks, until the hands have gone stiff and the light on the oiled paper has gone from white to the colour of skimmed milk. The sleet kept at the window, ticking and ticking, and the grey slit under the door kept its finger's width of \place{the Clay Yard}. \person{Corin} lay where he had been left, on his back with the damp grey wool under the back of his head and gone from his shoulder again, and the stone at his belly had settled into the one temperature it would hold from now until somebody thought to take it away, which nobody did. The flat note of the cup across the room was still in his ear --- the shadow of \lore{the full note}, a vessel that did not answer full --- and there was nothing in the room he could make answer, not the cup, not the wool, not the cold.

Word had it that \lore{the Five} were three days off, and \place{Low Kiln} said three days as a woman says the loaves she has left --- counted, and not wanted to say it. Beyond the last of the houses \place{the Ilm} lay under its skin of ice, and the ice was spoken of as holding, which meant that somebody had been out to look at it and had come back with a number they did not care to give. On the pallet \person{Dessa Silt} breathed the long shallow breathing that had been the shape of a woman spent, and her pulse had come up again now that there was nothing for it to do.

Three doors down the yard the kiln took a breath and let it out, and the sound came through the wall in a low intake and a long letting-out that he had stopped needing to listen for, because it had become a thing the room did, the way the room did the sleet. Four, he counted. Four. And then the gap --- the fifth that was not a breath but the place where a breath should have been, a small silence with the shape of a missing thing in it. He had counted the nails in the floorboards the same way, four on the near board and four on the next and a knot the width of a thumb where the fifth nail should have been, and he had not yet asked himself why the gap was the part he kept. It was the part he kept. \person{Daniel Marchetti} had been forty years old and had counted pulses in rooms where the machines did the counting now, and here he was running an instrument across a fire three doors off because there was nothing else in the world for the instrument to touch.

Dessa coughed. It came up from the pits, low and wet, and went back down again, and she called it weather because weather was a thing a body could have without anybody entering it against an account. She turned her face a little further into the wall and the cough went on behind her, and her thumb went to the pale scar on her left forearm, the long pale line where a sieve-wire had opened her years ago, and rubbed it while she decided how much to say. He had watched that thumb the day before. It was the tell of a woman weighing a truth and choosing the smaller piece of it, and he knew the shape of it as he knew the shape of a pulse that had changed and would not say why.

“Nineteen steps to the rack and nineteen back,” she said to the wall. “I've counted them since eight hundred and six. Four loaves in a batch, three to the levy and one with your feet up, and the seventh of Stone in eight hundred and twelve, which is today, if you're counting days the way I count steps.”

The pulse in her throat answered its own question, a small rise and then a settling. She would not say the other thing. The little hot one lay two paces off on a board-bed and the flower lay there, and neither of them had a name she would put in her mouth where a wall could hear it. A name was a thing a person said when they meant to keep something, and she meant to keep nothing that could be claimed by a woman with no house, so she kept the name the way she kept the bread money, in a place nobody was asked to look.

“A body can be warm and still be gone,” she said. “You learn that in the pits. They bring them up with the heat still in them and you hold them a minute before you set them down, and it doesn't make them yours.”

She said it in the proverb-form, so that it could be taken as advice if the listener preferred, and she did not look at the child when she said it. He had learned the form in two days. In Low Kiln a person said the difficult thing sideways, in a sentence about pots or weather or steps, so that the truth could get into the room without anybody having to own it. He had no doubt what she meant. He had not had a doubt about a thing like that since he was twenty-eight.

\scenebreak

\person{Alda Reed} answered the unasked question in her own order, which was to reach for the slate before she reached for the mother and to reach for the mother before she reached for the child. Her hand went flat on the slab on the bench, a thumb along the edge where the writing was turned under, a small adjusting push that made sure no dark stroke showed; then the stool, then the pallet, then the board-bed, in that order, as though the room were a body and she were washing it in the proper direction. The little tin on its cord inside her bodice shifted once as she leaned, and her hand went to it for the count of a breath and came back to the work.

“There was a cooper's seventh in \place{Mud End}, Y.H. eight hundred and ten, third day of Stone,” she said, to no one in particular, and set a cup down where the cup had been set before. “Father come down out of the peat with a cart and took him, and there was no trouble and no slate about it. Weather's turned. It'll go off by the second firing, same as last year, same as the year before that, and if it don't there's no call to fuss.”

The cup gave its small flat sound again and she did not hear it, or heard it and gave it no more attention than the sleet. The exact history sat in the room with nobody having asked for it, as a cloth laid over a thing sits with nobody having agreed to the shape of it. She had answered a question that had not been spoken aloud, which was the only kind of question asked in this room, and she had answered it with a date and a place and a man with a cart, and she had said the weather last, which in \place{Low Kiln} was never the subject.

The question itself stayed where it had been. Who paid to keep a foundling with no house-name, and for how long, and who had agreed to it: the questions sat in the space between the bench and the pallet as plainly as the grey ridge of \lore{the Ash} sat along the skirting board, the width of a thumb, lying there until \lore{the Last Day} and swept then and burned next morning and called dealt with. A child without a house-name was a coat left on a peg in a poor house: anybody might wear it, nobody need say who had left it, and the cold came in at the seam nobody had thought to mend.

Daniel Marchetti had read a room like this at four in the morning, before the day staff came and while the night still had the place to itself. He had stood at the foot of a table under a light that buzzed, smelling the metal-and-ammonia cleaning fluid and the blood, and he had known before anybody spoke that a decision had been made above him and would come down in a moment as a dropped instrument comes down. The senior would stop giving orders. The nurse would stop putting instruments down, her hands going still in the middle of a motion. The resident would look at the floor. Then the word would come, and the word was always the same three with different names in them: this one for the ward, this one for the corridor, this one for the drawer at the back where no name was written at all. He had said that word in his time. He had said it with a steady voice and clean hands and a chart signed at the bottom, and he had never once found a way to make it clean.

A voice from the life he had lost said its part of it in his ear. \memory{Decide, Marchetti. Then carry it.} That was the whole of the teaching, and he had carried it into a room in a clay town where there was nothing to decide and everything to carry.

He was the thing being sorted. He lay on the board-bed with his milky nails and the drool on his chin and the cold stone at his belly, and the two women in the room were arranging a kindness as women in a kitchen arrange a cloth before they fold it, with the corners squared and the thing underneath hidden. There was no word in \lore{Grathic} with which he could have said no. There was no hand at the end of either arm that could have reached across two paces of trodden ash to take part in the arranging. The competence and the helplessness were the same fact in one small body, and the fact had a temperature, and the temperature was the stone going cold against his skin while nobody in the room rose to warm it.

\scenebreak

Nobody said the arithmetic, but it was in the room all the same, as plainly as the Ash ridge was in the room: three of the four loaves to \lore{the levy}, the fourth eaten with the feet up or given away, depending on who had done the counting; the slip in the pit to be paid for in the thaw, when the brook came up and the clay went soft and a man could be sent down to do what a man had to do; and a bed and a warmth set aside for a foundling no one in the room had agreed to pay for. Somewhere inside that arithmetic a door would open on the third day and somebody would come through it. The room had said it without a word, as a slate says a figure when you hold it at the right angle to the light. He had the figure. He could not read the hand that had written it.

The warmth that had crossed this room last had come in the fold of an apron from a fire across the yard, a stone the colour of skimmed milk carried hot and set against him while he was too new to know that a thing could be done for him and never entered anywhere. That warmth was gone now. A brown-edged crust of barley, smoked with the same fire, had come across the floor the day before from a hand that had crossed the room briskly so that no speech could be made over it, and it, too, was gone or set aside, and neither the stone nor the bread had been written down in any column on any slate. He was the only one in the room keeping that account, and he kept it without a word, the way he kept the count of the floorboards and the count of the kiln: four, four, and the gap where the fifth should have been.

He remembered a hand from the day before. Two fingers had pressed the grey wool flat at his shoulder and had stayed a moment longer than the cloth required, and he had not known whose hand it was, and he had not wanted to know. There was a kind of love that was only a pressing of wool flat, a warming of a stone, a walking of nineteen steps to a rack and back again, and it did not announce itself and could not be entered on a slate because it had no figure and no name. It was done and it was gone and it left the cold in place. He had thought for forty years that a thing not spoken was a thing not done. In two days in a clay room he had come to suspect the opposite, and the suspicion had no comfort in it, because a kindness that cannot be answered is a kindness that cannot be refused.

The door lifted a finger's width at the latch as a gust took it, and the sleet came in and went along the trodden ash and stopped at the Ash ridge as though the ridge were a small cliff the weather would not cross. Beyond the slit the yard was the colour of a dirty slate, and the third day was somewhere in it, moving toward the room as a cart moves toward a house it has been told to reach. Dessa's thumb went once more to the pale scar and stopped. Alda's hand went flat on the slate and stayed. The kiln through the wall took four breaths and four breaths and then the gap, and the gap was the part he kept, and the stone lay cold at his belly with nobody in the wide world having thought to warm it.

\daymark{3}{Stone 9, Y.H. 812 (Ninthday)}

Woodsmoke came into the birth room before \person{Sella Vail} did, and under it wet wool, and under that the thin sour edge of \lore{the Ash}, so that \person{Corin} knew who was coming before the door had opened wide enough to show a shoulder. It was the third morning. The sleet of the day before had gone off into a dry cold that lifted the oiled paper at the window in small shivers, and the yard beyond the finger's width of \place{the Clay Yard} was the colour of a dirty slate, hard and bright and empty of everything but its own weather. The birth room of \place{House Silt} held that cold the way a cracked cup holds water, with care and with one thing nobody would mention until it was too late. The stone at his belly had passed through cold into a condition past temperature, and he lay with the damp wool under his head and the flat shadow of \lore{the full note} somewhere in his ear --- the cup was still two paces off on the bench, and would not be taken up, and he could not have said whether the note was in the room or in him.

No one had come for him in two days. Now someone had, and the room knew it before the room knew anything else: \person{Dessa Silt}'s breathing changed on the pallet, a small intake held and let go, and \person{Alda Reed}'s hand went flat on the slab on the bench, a thumb along the edge where the writing lay turned under. The door lifted at the woman's hand and would not shut after her, which was how it had behaved since the night he was born, and she gave it no more attention than the weather.

She did not come straight to the bed. She crossed to the bench first and put her hand near the slate and did not touch it, as a woman might put her hand near a sleeping animal to see whether it breathed, and then think better of waking it. Then she looked at the wall, at the window, at the cup, at everything in the room except the woman on the pallet and the thing on the board-bed. “Wrap him,” she said to the room at large. “There's cold in the yard. That wool won't hold. It never held.”

Then the hand came, and the hand that had hovered over him three days ago took. The burn across the back of her left hand was shaped like a little door --- two straight sides, a lintel and a threshold --- and as her fingers closed on the cloth the skin inside it flattened and the door shut first on nothing and then on the weight of him. He had taken the weight of a man through a coat collar on a floor in \place{Chicago} and known, before the whole weight came, how far the damage went; it was the same knowledge, the pressure through cloth, the small testing take, the settling. The hand was not gentle. Gentleness was the wrong word for it and he knew it, having had cause to hate gentle hands on bodies that needed decision. It was measured, which was better. It knew how much a thing would bear before it asked the thing to bear it.

“The boy comes under the roof this morning,” she said. “The house will keep him.”

She said the boy, not my son. She said the house, not I. The two small substitutions lay in the room as plainly as the grey ridge lay along the skirting board, and nobody in the room picked either of them up.

On the pallet Dessa got herself up while the woman was still at the bench. She stood with one hand on the wall and the other at her forearm, and her pulse had come up into her throat where the surgeon in the bundle could read it as he had read a pulse on a wrist in a room where nobody knew he was reading. She did not look at Sella and Sella did not look at her. The two women arranged their faces with the care of people folding a cloth they intend to lay over something.

The stone did not go with him. It stayed in the fold of the grey wool on the board-bed, a small pale thing the colour of skimmed milk, and he felt the place of it at his belly as an empty place is felt, with an exactness that had nothing to do with temperature. He had been warmed at the fire that was now carrying him away, and the warmth had been spent and not entered anywhere, and the stone that had held the last of it was left behind in a room that would not think to warm it again. The cup gave its flat note to the empty bed. The door kept its finger's width.

\scenebreak

He counted her steps because there was nothing else to count. It was a new count laid over the old one, and he kept the two apart as a man keeps two pulses apart in a room with more than one body in it. Dessa's count was nineteen out and nineteen back, a wheel that came every time to the same place; this was a line, and the line went one way. One, two, three. Her stride was short and even, the stride of a woman who carried weight for a living and had never learned to waste any of it. He had counted four nails and a knot in the floorboards for two days; out here the numbers were larger and there was no knot, only the long unbroken run of a woman going somewhere she had decided not to name. He had the count before he understood that he was making it.

The Clay Yard opened beyond the door in hard light. Woodsmoke lay thicker with every step, coming down off the roofs and out of the \place{Vail kiln-house} three doors along, so that the smell of the woman's coat was only the front edge of a greater smell waiting at the end of the walk. Behind the houses \place{the Slip Brook} ran clay-coloured under a skin of ice at its edges, and he knew it by the mineral smell before he knew it by anything he could see, because the bundle let him see very little: the underside of a chin, the collar of a coat, the sky stripped of colour, the ragged top of a wall, the drying racks with nothing on them but the weather. Out beyond the last of the roofs \place{the Ilm} lay under its ice, and \lore{the Five} were three days off, and \place{Low Kiln} said three days when it had counted what it had and did not want to say it.

The kiln breathed closer. Four, he counted, four, and the gap where the fifth should have been; and then four, four, and the gap again, the intervals shortening as they went, which was not the fire changing but the woman carrying him getting nearer to it. He had no idea where she was taking him. He had an idea of the distance and the direction and the number of steps, and no idea at all of the purpose, and the two kinds of knowing sat in him without touching, as the surgeon's register and the infant's hunger had sat in him for three days without touching. There was a little frozen mud on the threshold of every door they passed, and a pot set outside one of them with a rim of ice in it, and a woman at a tub who looked up and went back to her work without appearing to have seen the bundle at all.

He tried once to turn his head toward the sound. He went past the point and the small animal noise came up in his throat, the noise a thing makes when it is handled and cannot help itself, and the head went where the head wanted, which was nowhere he wanted it to go. Her hand tightened one degree on the bundle, not unkindly, the way a nurse steadies a thing that is about to slip, and he understood that she had taken the sound for a child's fret and had no notion she was carrying a man who was trying to look at a roof. The drool went down into the wool. He could not wipe his own chin. There was a dryness in that, and he took it, as he had taken the joke of the blanket, because the alternative was to weep in a bundle and be mistaken for a wet child.

She said nothing on the walk. Once, when the breath of the fire came close enough to be felt through the coat, she made a small sound in her throat that was not a word, and her thumb moved on the back of her hand where the burn was, and that was all. He counted. Nineteen was Dessa's number and not hers; hers ran past it without stopping, ten, twenty, thirty, and he found he was keeping the number because a number was a thing that could not be taken away by a woman who had decided not to explain herself. The woodsmoke went into the wool and stayed. The cold went through the coat and found him. Somewhere ahead of them a door opened on the fire.

\scenebreak

Dessa stood at the end of the yard. She had followed as far as the ground where the houses stopped and the rise toward the wheel-shed began, and there she had stopped, and she did not come the last nineteen steps or any part of them. Her pulse was in her throat. Her thumb was on the pale scar on her left forearm, the long pale line where a sieve-wire had opened her years ago, and she rubbed it with the small pressing motion he had learned to read as the tell of a woman weighing a truth and choosing the smaller piece of it. The smile was on her face, the small closed smile her sister called the door shut, and it said that nothing further would be said.

She said nothing further. She did not say his name, which she had not said in three days and would not say now that he was going away under another roof, because a name was a thing a person said when they meant to keep something, and she had arranged herself to keep nothing a woman with no house could be accused of claiming. She counted under her breath instead. Nineteen out and nineteen back, the count she had kept since eight hundred and six, and she did not take one of them. The day before she had said that a body could be warm and still be gone, and had said it sideways so that it could be taken as advice; now she was doing the thing she had described, and the surgeon in the bundle heard her counting the way he had once heard a monitor in a room go from a rhythm to a straight line, a sound that had been there a moment ago and was now a line with nothing behind it, and he knew what it meant and had no way to answer it.

Alda came up behind her with her hands at her apron, smoothing it in two passes, the way she smoothed a cloth before she laid it on a body. She had come from the birth room without looking back at it, and she stopped a pace short of Dessa and did not put a hand on her.

“He's entered,” she said.

Dessa did not turn from the yard. “Where?”

“In the levy. Among the children taken for the account this month.” Alda's voice went exact, the way it went exact when a date was needed. “Seventh of Stone, eight hundred and twelve. The account went up on the eighth, before the paper went out, the way it goes up every year. There was a cooper's child in \place{Mud End} in eight hundred and ten, third day of Stone, father come down out of the peat with a cart, and there was no slate about it and no trouble. It's a slater's business, not a mother's. Warm and breathing was all I set down.”

Dessa asked nothing further. The yard went on doing what yards do: the smoke lay low, the ice held at the brook's edge, something small moved in the wall and stopped. Alda's exactness went on after the question had stopped, as water goes on running after the tap is shut. “A woman with a thin return goes asking after the levy and a hand comes down on a slate that was better left where it lay,” she said. “You'd only draw attention to it. There's no call to fuss. It'll go off by the second firing, same as last year, same as the year before that, and if it don't there's no call to fuss.”

The figure was not spoken. It lived instead in the pressure of Dessa's thumb on the pale scar and in the number of dates Alda gave where no date had been asked for, dates arriving faster than dates had any business to arrive, as a cloth laid over a thing arrives before the thing has been given its shape. The boy went up the yard with the number unspoken behind him and the story of it settling into the space where a name should have been.

\person{Daniel Marchetti} heard his own life being rewritten behind him and could not turn around to correct it. He had heard a thing like it done in a room with a light that buzzed: a chart that said one thing and a room that had done another, and a signature at the bottom to make the two agree. He had once asked a family to sign a form and had called the signing a conversation. He had written the neat version himself, in a hand that was steady, and had gone home and eaten his dinner. The neat version was kinder to everyone. It spared the woman in the yard a scene, and it spared a house a slate read twice, and it cost the only thing it was pretending to be about. He understood the arithmetic of it with the part of him that had once decided which bodies went to the ward and which to the corridor and which to the drawer at the back where no name was written at all. He had made that decision in a steady voice and with clean hands, and had never once found a way to make it clean.

But he had not been a surgeon for three days. He was a bundle under a coat with a drool on his chin and a head that went where it wanted, and the woman carrying him did not know that anyone was listening. He had been disposed of as an institution disposes of an account: not cruelly, not even carelessly, but in the way that is hardest to forgive, which is the way that believes itself to be mercy. He held the empty place where the stone had been at his belly and he held the count of her steps, and he had no voice in either of them, and the flat note of the cup was still in his ear.

Behind him the yard kept its own account: Dessa at the end of it counting nineteen steps she did not take, her thumb on the pale scar, the door-shut smile staying where she had put it; Alda's exact dates settling over the space where a name should have been. The kiln took four breaths and four and the gap where the fifth should have been. And the boy went up the yard under the coat of a woman who would not say why she was keeping him, while the woman he had come from stood at the end of nineteen steps and did not follow.

\scenebreak
The door of \place{the Vail kiln-house} gave under \person{Sella Vail}'s shoulder, let the cold in for a moment behind them, and then the heat took them. It was not the heat of a room with a fire in it. It was the heat of a thing kept alive a long time on very little, coals banked under ash, the low intake and the long letting-out that he had counted through a wall since the night of his birth and now felt against his cheek with no wall between. \person{Corin} had come to the end of his count. The breath was no longer on the other side of anything; it was the air he lay in. The \place{Clay Yard} went on doing what it did beyond the door, which was nothing that concerned this room, and the door, having let them through, went back to its old habit of not quite shutting.

The heat went to his hands first, the way heat does in a body that has given up expecting it, and came on afterward in his feet, a slow return that itched and stung. He knew the shape of it. He had watched men come out of the cold into a warm room and seen the blood go back into the extremities and the pain arrive with it, and here it was happening to him under a coat that smelled of another woman's smoke. The kiln sat somewhere to his left and kept its own time. Four, he caught himself counting, four, and then the gap, and he stopped, because the gap had been a thing to listen for through a wall and here it was only the house breathing in its sleep.

The floor was trodden ash like the floor of \place{House Silt}, but the ridge of \lore{the Ash} along the skirting board was thicker here, near the width of two thumbs, running the length of the wall and turning the corner without a break. In House Silt the grey had lain in a line that the midwife's skirt lifted clear of; here it had been left to gather until it had an edge and a shadow of its own. The sour mineral tang of it came up under the woodsmoke and the wet wool. Nobody had swept it. Nobody would sweep it before \lore{the Last Day}, which was three days off, or five, or whatever number \place{Low Kiln} had decided to give out that morning. He took the ridge in as a fact and put it with the others he was keeping: the door that would not shut, the stone that had gone cold in a fold of grey wool three doors away, the hand at the edge of the flare light that had not stopped writing.

She laid him on a bed under the roof. There was fresh straw in it, still with the dry stalk smell on it, and a fold of cloth over the straw that had been warmed either by somebody before him or by the room itself; he could not tell which, and the not-knowing was its own small irritation. The hands that set him down were the hands that had lifted him in the yard, measured, unhurried, knowing what the thing would bear. They did not linger. That was the kindness of them. A bed was a thing that could be entered in a roll as easily as a name could; he had seen a ward enter men with less ceremony than this, a wristband and a number and a bed, and the bed had been the first proof that the entering had happened. He was warmer than he had been that morning and more watched, and he could not turn his head far enough to see who was watching.

\scenebreak

\person{Sella Vail} went to the kiln with her hand at her throat. The thing under her fingers was \lore{the widow's mark}, a thumb-length of lead-glazed clay on a cord, the Vail kiln scratched into it in a hand that had pressed hard enough to leave the groove rough; he saw it for a moment against the coals and then it went under her collar again. She touched it before she opened the kiln and at no other time, and he set the observation down beside the dates \person{Alda Reed} had given in the yard and the pause in the writing hand on the slate in House Silt, three things that had the same shape and did not yet make a fourth. The kiln door came open on its own weight. Heat leaned out of it and the coals showed red under their grey lid. She did not look toward the bed. She looked at the pot nearest her, which had nothing in it, and told it what to do. “Don't set that near the door,” she said. “You'll pay for it in the thaw.” The fire got the same instruction in the same flat voice, and then nothing, because there was no finished feeling to give it.

The arithmetic of the house was not written down anywhere he could see and was written all over it. The foundling had the bed under the roof. The clay-bound hand had the \place{Debt Loft} above the \place{Vail wheel-shed}, which he knew because a man crossing the room above their heads made the small settling sound of a board taking weight, and because \person{Tarn} came down the ladder with a broom in his hand and went straight to the floor as if he had been sent for it and no sending had been necessary. \person{Daniel Marchetti}, who had once stood in a corridor of stretchers in the hour before morning while the decision came down from people who were not going to do the carrying, looked at the bed and the loft now with the same part of him. A poor household had sorted its own members and called the sorting kindness. He had done the same thing in a steady voice and with clean hands. He could not lift one of his own to object.

\scenebreak

\person{Jorra Vail} stood at the wheel in the wheel-shed with his weight settled into the worn place on the left side of the seat, which was the place his father had worn and which no one else in the house was allowed to fill. He had been trying to fill it with himself for three years and had not managed it, and he did not look as though he had noticed. Slip had dried to the colour of old bone up both his forearms, and the little finger of his left hand bent at its middle joint where a splinter had gone in when he was eleven and had been left there until the flesh grew over it. His hands moved the wheel and then stopped it, and when the pot on the wheel had been true for a moment he said to it, “You'll do,” in a voice with no welcome in it and no unkindness either, which was the way he said most things.

He had said one sentence about the boy in this room before the boy had been carried into it: \emph{He's not Vail.} The sentence had been refused an answer, and the arrival had sealed it. There was a child under the roof now and nobody had written him as Vail, and the house went on around the fact the way a pot goes on around a flaw that has been decided not to mend. It was what Jorra had wanted. He did not look at Sella while she spoke to him, and he obeyed her before she had finished, which was the habit of the house and the correction neither of them made. When she told the fire to take the draft off the left, he got up and shifted the damper without a word, and came back and put his weight into the hollow again. He did not speak to the boy. He spoke to the wheel, and to the pots, and once to Tarn about a rasp left out of its place, and the boy lay under the roof and was not addressed by any of it.

The surgeon in the bed read the bend in the finger and the slip on the forearms and the set of the shoulders as a man reads a chart left open on a desk: sixteen years old, and old in the hands. He read the not-looking too. He had known residents who would not meet the eye of the senior because the thing they had to say to the senior could not be said, and he had mistaken the silence for respect for most of his career. Here was the same silence with a wheel under it. He had no words in Grath and no breath behind the words if he had them, and he took the reading anyway, as he took all his readings now, and kept it.

\scenebreak

Tarn swept the floor in front of the kiln with short economical strokes, moving the broom as if the broom were another man's and he might be asked to give it back. He was bent at the shoulders from six years in the slip pit, which made twenty-two sit on him like a borrowed coat, and grey Ash sat in the creases of his knuckles that no scrubbing lifted. There was a pale ring round his left wrist where an iron clip had sat in eight hundred and six. He called Sella house when he had occasion to answer her, which was often, and said I owe where another man would say I am, and agreed with everything so completely that there was nowhere for an argument to go.

Once, while he thought nobody was counting the small tools, he took a bench knife out of his pocket and set it on the bench with two fingers, and then, having set it down, picked it up and put it back where it had been, and said without being asked, “I took the small knife. I put it back. You can check.” It was a flat knife with its mark worn nearly away from use, and the hand that put it back knew the weight of it the way a nurse knows a missing instrument before she has looked in the tray. He had kept a note of every day of his term and would not be told the number of the next. “Six years. One to go. I've counted every one, and I won't count the next,” he said to the floor, in the tone of a man stating the price of something, and went on sweeping.

When Sella said that the boy would need watching --- she said the house would watch him, which came to the same thing --- Tarn stopped his broom long enough to put the one sentence into the room that nobody had asked for. “He's a foundling. I'm a debt. Don't put us in the same sentence, Sella.” Then he went back to sweeping, because a clay-bound hand who complains about beds is a clay-bound hand who can be returned to the slate. The surgeon on the bed heard it with the part of him that had once sat at the edge of a ward and listened to what the orderlies said when they thought the patients were asleep. It was not resentment. It was a man keeping the one thing he owned, which was the exact account of what he was owed and what he owed, and grudging the house even the loan of a bed for somebody else.

A pot was set down on the bench while Tarn was midway through a stroke. It gave a small flat sound, the same dry answer the cup had given across the birth room two days before, a vessel that did not answer with \lore{the full note}. The broom stopped. Tarn's head went a degree toward the bench and came back, and nobody in the room said what the sound had been, and he went on sweeping with the stroke he had left in his hand. The surgeon lay with the note in his ear and did not join it to anything. It was the second flat note in three days and the second one he had done nothing about.

After that Tarn banked the coals, which meant lifting the grey lid carefully and setting it back so that the red underneath kept its slow hold, and he did it without looking at the fire, the way a man buttons his coat in the dark. When he had finished he stood still for a moment because there was nothing in his hands, and stillness in that house was a thing that got a man sent to find work, so he found a corner of the floor that had not been swept and swept it again.

\scenebreak

\person{Garl Vail} was not in the room and was the shape of most of it. The wheel-seat in the shed was worn hollow on the left where he had sat for twenty years, and Sella would not let anyone fill it, though she let her son sit there, which was a different thing and had never been explained. Behind the door hung a coat with salt stiffened into its seams and a stain across the shoulders the colour of old bone, which Jorra had twice reached for in the dark and stopped, and which nobody took down. The broom stopped at the ridge of Ash and went back the other way. The ridge had grown another inch since the Long Cold and nobody had swept it and nobody would say why.

The proverb lived in the house's mouth the way the coat lived on its peg, handled and not spoken to. \emph{Leave the crack showing. People trust a thing they can see has been handled.} It had been Garl's before it was anybody's, and now the house said it instead of him, to the pot and the fire and the floor, never to a person and never as advice, which was how the house said everything that mattered.

The surgeon on the bed took the hollow seat and the salt coat and the ridge along the wall and counted them the way he had counted the floorboards in House Silt and the steps up the yard and the breaths of the fire through the wall: four, four, and a space where the fifth should have been. The space was a man. It had the shape of a man who had sat in one place for twenty years and got up and gone out into the cold and not come back, and the house had arranged itself around the space as a body learns to move around the place where an arm used to be, reaching for it still and finding the air. He had felt the place of the cold stone at his belly in just this way, an empty place felt with an exactness that had nothing to do with temperature. Here the emptiness had a coat on a peg and a hollow in a seat and a ridge of grey gathering in the corner, and the house called it the order of things.

In the ash of the kiln, where the grey had been pushed aside to let a draft through, lay the charred edge of a slate, black on one side and eaten into the shape of a strip of bark, the writing gone off the end of it. It lay where a slate had been put into the fire and not pulled out in time. He saw the edge of it from the bed and did not examine it. He was three days old and could not reach it, and there was nobody in that room who would have understood the question if he could have asked it.

The afternoon came on without anybody marking it. The boy lay on the bed under a roof that was not his, with the salt coat on its peg behind the door and the hollow seat in the shed and the Ash ridge keeping its place along the wall, and the count still his because it belonged to nobody and could not be entered on a slate. He had come from a house that would not name him into a house that would not ask where his memory came from, and both had kept him warm in their own fashion, and neither had asked him anything at all. The third day was done. The market and \place{the Cross} were still ahead, and \lore{the Five} were three days off, and the fire kept its four and four and its gap, and the boy listened to the gap because it was the only part of the count that was his to keep.

\CRtailpiece
\begin{CRchronicle}{2}
Before the boy came under the Vail roof, the house had arranged itself around an absent man. The kiln drew four breaths, four, and missed the fifth. The wheel-seat kept Garl Vail’s hollow; the salt coat kept his shape on the peg; the Ash ridge had grown past sweeping. Tarn swept, called Sella house, said I owe, and put the bench knife back. Jorra at the wheel never looked at the woman who ordered him. When she said the boy would come under the roof, he said he was not Vail, and she answered the pot instead.

A cup on the bench gave back a small flat note --- a vessel that did not answer full, the first tool his old life had found an object in, two paces out of reach. Corin, two days old, could not tap it or say what he had heard. Deshawn Carter’s laugh came up out of a car park in Chicago --- bad pizza, silence, a minute timed on a radio --- and went out like a lamp. There was no Grathic word for him; the ache went in under the flower on his ribs.

Alda Reed washed him, reaching for the slate first, giving dates where comfort was expected, crossing two fingers over the birth-flower unaware she did it before witnesses. Dessa coughed from the pits and called it weather, counted nineteen steps to the rack and back, rubbed the sieve-wire scar. She would not name him. A body can be warm and still be gone, she told the wall, so that the truth could enter without anybody owning it. The question of who paid to keep a foundling with no house-name sat between bench and pallet.

On the third morning Sella came ahead of her smell of woodsmoke and took the bundle with the burn-door closing on the weight. She said the boy, not my son; the house, not I. The cold stone stayed in the fold of grey wool on the abandoned bed. Dessa stood at the end of the nineteen steps and did not follow, her thumb on the scar and the door-shut smile on her face. Alda told her he had gone in the levy, among the children taken for the account, and gave dates faster than dates had any business arriving. No figure was spoken.

He was laid on fresh straw in a bed while Tarn was given the loft. Sella touched the widow’s mark before the kiln and told the pot it would pay in the thaw. Jorra’s sentence stood, sealed by the arrival: nobody had written the boy as Vail. Tarn took the knife out and put it back, and stopped his broom when a pot was set down. In the ash lay the charred edge of a slate he could not reach. He ended the day warmer and more watched, keeping four and four and the gap.
\end{CRchronicle}

\backmatter
\chapter*{Dramatis Personae}
\markboth{Dramatis Personae}{Dramatis Personae}
\addcontentsline{toc}{chapter}{Dramatis Personae}
\CRcharacter{Daniel Marchetti}{a butcher with a surgeon’s hands, awake in a child’s body}{build/portraits/dan-marchetti.pdf}{\CRfield{Also called}{Dan Marchetti, Dan, Corin, Corin of Low Kiln, the butcher, the foundling with the birth-flower}\CRfield{Role}{protagonist}\CRfield{Born}{1984, Hermosa, Chicago, Earth (first life); Day 1, Stone 7, Y.H. 812, Low Kiln (new vessel)}\CRfield{Kind}{human — an Earth man of the 2020s, now a vessel of Sarn}\CRfield{Origin}{Hermosa, Chicago, Illinois, Earth; reborn in Low Kiln, Sill, Grath}\CRfield{Allegiance}{Cook County trauma service (first life); House Vail kiln-house, Low Kiln (by rebirth); unwitting instrument of the Unbanked}\CRfield{Status}{alive (as Corin, three days old at the close of Day 3); cracked, low-capacity, unranked, below Thin; died 2024, Chicago, in a fire after a blizzard (first life)}\CRfield{First seen}{Chapter 1}}{\CRsheetheading{Appearance}
In the old life he was a large man gone soft in the middle, with forearms like a butcher’s and a left knee that clicked when the weather turned. His hands were short-fingered and steady, with a pale line where a wedding ring had been; his hair went early at the temples, and on post-call days he shaved badly and did not care. He smelled of hand sanitiser, machine coffee and the faint metallic tang of a hospital cleaned too many times with the same mop. People read his face as anger. It was fatigue, and he never corrected them, because correcting them cost something.

In the new life all of that is folded into an infant body that cannot hold up its own head. The surgeon’s hands are now small soft wrists with milky nails; the butcher’s forearms are gone; the clicking knee is a soft unformed joint that will click again in time. Across his ribs and down his forearm runs a branching red mark like a frost-flower or a burn map, which the midwife Alda Reed crossed herself at and called a birth-flower, and which is in truth the fire writing itself into the hot trace that crossed with him. He drools. He cannot turn over. He is washed, dressed and discussed in the third person by people who do not know there is a forty-year-old man listening. He is a trauma surgeon who cannot reach a cup. He counts the nails in the floorboards because there is nothing else to count, four on the near board, four on the next, and a fifth that is not a nail but a knot the width of a thumb.

\CRsheetheading{Voice \& Manner}
He speaks little, plainly, low, and exactly. His question is usually the shortest one in the room. He does not use the old-life idiom with people who would find it strange; he never says \emph{I was a surgeon once}, because no words in Grath would carry it. With the few adults who can bear it, the Chicago rhythm comes back: faster, profane, warm under the withering. Under danger his sentences get shorter, not longer, and the clinical register takes over completely.

\begin{itemize}[leftmargin=1.4em,itemsep=1pt,topsep=3pt,parsep=0pt]

\item “How long since he stopped moving? Not the story. How long.”

\item “I’m not asking you to pray. Hand me the clamp.”

\item “That’s a lovely account. Shame nobody’s holding the heat.”

\item “Leave it. It’s gone.” — the sentence that is his gift and his heresy in Sarn.

\item “Don’t set that down. You’ll pay for it in the thaw.” — learned from the clay yard, said in the local proverb-form before he has fully learned to mean it.

\item Remembered, in the old life: \emph{You’re back on in ten, Marchetti.} said in a voice like a file.

\end{itemize}

He cannot stop reading bodies. At a family dinner he noticed a new knuckle on his brother’s hand and a cut his aunt wiped on her skirt before anyone saw. He counts without deciding to: pulses per ten seconds, exits in any room, steps from the bed to the door. His humour is fast and withering and used like a handrail. He finishes everything on his plate, including food he dislikes, because waste offends him in a way he cannot explain.

\CRsheetheading{History}
He was born in 1984 and raised in a two-storey brick bungalow in Hermosa, on a block where the el ran high and iron over Pulaski and the summer smelled of charcoal and wet pavement. His father Aldo was a machinist from near Bari who showed love by fixing things and never learned to say it; his mother Theresa worked as a nurse’s aide and came home smelling of bleach and lilac soap. There were three children: Dan, Jimmy the patrol officer, and Lisa who left for Florida. The house was loud and full and he learned from it volume, profanity, the duty of finishing everything on his plate, and the certainty that love is a practical act refused in conversation. He went to Lane Tech because his mother insisted, took things apart on the shop floor, and found the human body in the cadaver lab at the University of Illinois Chicago the way other boys find engines: by being forbidden and doing it anyway. Anatomy did not argue. It had a right answer under your hands.

He trained at Cook County, became an attending, worked nights because nights paid more and days were when the administration was awake. What he knew, precisely: airway in blood and dark, a cricothyrotomy with a knife and tube if the mouth was destroyed; the difference between pulsing blood and welling blood, packing a cavity, the third rib that tears the lung; reducing shoulders and pinning femurs; burns by the rule of nines and the Parkland formula; hypothermia the way a priest knows sin — the afterdrop, the fact that a cold patient cannot clot, that the grimace of severe cold looks like a smile, that you do not rub a frozen foot; sterile field, debridement, leaving a wound open; the small evidence a body carries; and above all triage, the ninety-second sorting of twenty injured people into those who will live, those who will live only if he chooses them, and those who will die whatever he does. He did it hundreds of times and never once found a way to make it feel clean.

What he did not know mattered as much. He was a clinician, not a physicist. He could not grow penicillin, anaesthetise with a weed, make a fire with his hands, track a beast, bargain for food, read a Sarnish ledger, or make a native believe a clamped vessel might survive. He was an amateur boxer from his twenties — enough to keep his chin down and take one punch and wait. He knew exactly how a body fails and was almost helpless when the body in question was his own.

He married Elena Rossi in 2007; their daughter Nina was born in 2000. The shifts ate the marriage. He was absent for the dinners and the small ordinary evenings that turn out to be the whole of a life. The divorce was quiet and worse for it. He kept a photograph of Nina at six on his wall and looked at it every night without seeing it. There was Marisol Vega after, an emergency-room nurse kind enough to try for eleven months, who left because she was tired of being the woman who waits. There was Des Carter, who made jokes instead of asking; Ostrowski, who taught him to decide and carry; Vesci, who called him the butcher and never forgave the report.

When he was thirty-one a twenty-two-year-old named Tasha Broussard came in after a fight, was examined, had a normal blood pressure and a bruised abdomen, and was discharged with a sprain. She returned four hours later in cardiac arrest from a ruptured spleen the size of a fist. The junior had missed the sign; Dan, covering the shift, signed the chart without reading the nurse’s note that said she had been vomiting. At the review he let the hospital call it a systems failure and said nothing about the note. He kept his job, his licence and his reputation, and carried the shape of his own conduct for the rest of his life. His father died that spring while Dan was on a thirty-six-hour shift; he noticed the casket was slightly crooked in the aisle and thought about it for years.

He died in February, in a blizzard, on the expressway at 2:17 in the morning, going back to the minivan for the last time and getting the child out before the fire took the van and took him. His phone was lit in his coat with a name he would not look at. The last thing he knew, orderless: the hiss of fuel, the orange light on snow, the cold leaving his hands, and a wordless sentence that was nevertheless complete — \emph{I have left the job unfinished.} A woman in a coat the colour of wet ash stood at the edge of the flare light writing on a slate, and her writing did not stop.

He woke as an infant in Low Kiln, taken in by Sella Vail of the kiln-house, with the birth-flower across his ribs and a warm stone at his belly that was already going cold.

\CRsheetheading{Wounds \& Fears}
The body’s humiliations are exact: drool, the useless head, the blanket that will not stay on, the small animal sound that comes when he turns his head past a point, the cup two paces off that he cannot reach. Day 1 is lived awake from the birth to morning; nobody thinks to ask whether he wants to be awake. By the afternoon of Day 3, under the Vail roof, heat has come back into his hands and feet with the itch and sting of a rewarming, which he reads in reverse; the pale stone that held the last of that warmth was left in the fold of grey wool in House Silt, and he feels the place of it as an empty place felt with an exactness that has nothing to do with temperature. But the deeper wounds are old ones. The missed call to Nina, promising a boat on Lake Michigan for nine years. Tasha’s chart and the note he did not read. The crooked casket nobody corrected. A sister in Florida, a mother’s kitchen, a father who could not say what he felt and so fixed things.

He fears being the only adult in the room and him with the hands of a child. He fears that his clinical gift is not a gift at all but a moral defect — that the willingness to say \emph{leave it} says something about him he would rather not know. He fears the hollow and the Kept: a vessel that holds perfectly and feels nothing, which is the thing he has spent his life doing to himself in the name of duty. He fears going home most of all, because the mechanism is real and the price is the warmth of the woman the child in the minivan became. He cannot pay that account by dying again; it can only be carried.

His contradictions are the engine of him. He who decided who was beyond saving now lives in a world whose founding doctrine is that nobody is ever beyond keeping. He who fled intimacy now carries a whole district’s debt. He distrusts institutions because he has performed their gesture himself — and yet he learns their trades, joins a household, and will one day wear the Court’s own disciplines like a skin. He is most himself when he is least comfortable, which is a dangerous thing to be in a world where comfort is stolen from the poor and he will not notice that somebody else is paying for his.

\CRsheetheading{Desires}
He wants to go home. He wants to hear Nina’s voice, to say the nine-year-old sentence about the lake and the boat, to put down the phone he keeps reaching for in a pocket that holds nothing. He thinks he wants a better clamp: a way to mend the world’s crack without spending anyone, a way to pay the debt of the Forty of Low Kiln whose warmth bought his whole return. What he needs is harder and he will spend thousands of chapters arriving at it: that repair is not restoration, that a mended vessel carries its break, that a debt can be carried instead of settled, that attention is the one thing no institution can tax, and that the answer is not a tighter seal but a third thing — a repair that leaves the crack visible and the Warm free to move.

\CRsheetheading{Bonds}
\textbf{Nina Marchetti} — his daughter, the light she stopped leaving on. He loves her in the only idiom he has: practical acts, unsaid. She resents him for it. The unsent call is the shape of the debt that made him portable; the words never crossed, only their weight. He will reach for a phone for her at every important death in Sarn and never look at the name. \emph{Ch 1:} she is the lit rectangle in the coat pocket on the expressway, the boat promised nine times, a name he will not look at; heavier than the phone.

\textbf{Deshawn Carter} — the paramedic friend who made jokes instead of asking. Their friendship was ten minutes of bad pizza and silence in a car park. Des’s laugh is the sound Dan reaches for when he cannot reach Nina. The tension: Des never asked the question that would have let Dan confess, and Dan never gave him the chance. \emph{Ch 2:} the laugh came up whole out of a car park in Chicago --- bad pizza going cold on the console, wet slush under the lamps, a gold tooth catching the light, an eleventh minute that never came, and the one minute timed on the radio after the expressway --- and went out like a lamp. There is no word in this country for a man who rode in the back of a rig and made jokes instead of asking; the ache went in under the flower and stayed.

\textbf{Lena Ostrowski} — the senior resident who taught him the first duty is to decide and the second is to carry it. He calls her the closest thing he had to a mother in the hospital and sent flowers to her retirement and did not go. Her maxim is both his gift (the courage to say \emph{leave it}) and his excuse (the silence over Tasha’s chart). She is the voice of his heresy and the voice he must finally outgrow. \emph{Ch 2:} her correction arrived in his ear while the room was arranging him: \emph{Decide, Marchetti. Then carry it.} He carried it without deciding, which is the finer art.

\textbf{Paul Vesci} — the rival who called him the butcher. Vesci is the man Dan could become: a man whose institution absorbs his guilt and hands him back clean. Dan reported him once and has spent fifteen years being called a butcher for it; the secret of Tasha’s chart makes the nickname true in a way Vesci cannot know. They pull against each other across a gulf of mutual, unspoken recognition. \emph{Ch 2:} he heard his own life rewritten behind him as a chart that said one thing and a room that had done another, with a steady hand and a signature at the bottom to make the two agree --- the grammar he learned in Vesci’s hospital and has never unlearned.

\textbf{Marra} — the collector in the grey coat at the crash, who wrote his measure and paid for his return with her own fingers. He does not know her name in the early chapters; she is the woman at the edge of the flare light. She has been paying for him ever since. \emph{Ch 1:} he has filed her as a coat, a head, a hand, and the rhythm of a pause between strokes; he has not put the hand at the flare beside the dark hand on the birth-slate, though both are on the page. \emph{Ch 2:} he set the pause in the neat dark hand on the birth-slate beside the pause in the writing hand at the flare a second time, and did not join them.

\textbf{Sella Vail} — the widow who took him in and has never once asked where his memory came from. She is waiting to see what he will become. Her silence is its own kind of Court: kindness that does not want an answer. \emph{Ch 1:} her hand came in ahead of her, felt for the heat at his belly, found none, opened and shut once over the bundle and withdrew; she could name the house that kept the fire but not the hand in the door. \emph{Ch 2:} the hovering hand became a taking hand, measured and unhurried and knowing what the thing would bear; she said \emph{the boy} and \emph{the house}, laid him on fresh straw, told the pot not to be set near the door, and wrote him down as nothing he could see.

\textbf{Alda Reed} — the midwife who delivered him, reached for the slate before the child, crossed two fingers over the flower and called it a birth-flower. \emph{Ch 1:} he read her order of reaching, her tin touched twice, and her exact dates covering a thing she would not say. \emph{Ch 2:} she washed him in the same order, reached for the slate first, touched the tin, crossed the sign at the corner of the flower in front of him without knowing she had done it, and told Dessa he had gone in the levy.

\textbf{Dessa Silt} — the mother who will not name him. \emph{Ch 1:} he read the pulse that rose when her count stopped, the sieve-wire scar she rubbed, and the two fingers that stayed too long pressing the wool flat at his shoulder. \emph{Ch 2:} she followed the bundle to the end of the yard, counted nineteen steps she did not take, stood with her pulse in her throat and her thumb on the pale scar and the door-shut smile on her face, and let him go believing he had been spent among the levy children.

\textbf{Mira Silt} — the aunt who knows. \emph{Ch 1:} she stood in the one patch of floor the slate did not reach, said the secret was in what happens to a child after it is born, and crossed the room to give bread without once looking at the writing. \emph{Ch 2:} she is not on stage; the brown-edged crust from her hand is still the last warmth that crossed House Silt before the boy was taken.

\textbf{Jorra Vail} — the son of the house, sixteen and old in the hands, who said \emph{he’s not Vail} before the boy came under the roof and was refused an answer. \emph{Ch 2:} the arrival sealed the sentence --- a child under the roof and nobody wrote him as Vail, which was what Jorra wanted and what he could not bear. He did not look at Sella, obeyed her before she had finished, spoke to the wheel and the pots and once to Tarn, and never addressed the boy.

\textbf{Tarn} — the house’s clay-bound hand, six years into a seven-year term. \emph{Ch 2:} he was given the loft above the wheel-shed while the boy was given the bed, and said \emph{He’s a foundling. I’m a debt. Don’t put us in the same sentence, Sella.} He took the small bench knife out and put it back where it could be checked, banked the coals, and stopped his broom for two strokes at the flat note of a pot set down, without saying what the sound was.

\textbf{Garl Vail} — the husband of the house, dead in the Long Cold and never collected. \emph{Ch 2:} Dan read him not as a man but as a space --- a seat worn hollow on the left, a salt-stiff coat on its peg, an Ash ridge the width of two thumbs, a proverb in the house’s mouth --- and counted him four, four and a gap.

\CRsheetheading{Arc So Far}
He has lived three days, awake through most of them, in a body that can do nothing he wants. The double consciousness has staged itself in full: the clinical register reads hand-grips, pulses, the bend of a finger, the not-looking, the return of blood to cold extremities and the arithmetic of a bed and a loft; the infant body drools, hungers, cannot hold up its head, and wins every argument by being louder. He is cracked, low-capacity, unranked --- below Thin, twelve to fifteen embers at best --- with no Warm moved and no power used. Three facts have gone in and cannot be got out: a name was written on the slate before anyone saw his face; a household arranges its silences around a thing it will not say; and he has been given up by a mother who loves him and told he is dead by a woman who is lying to keep him safe. The clinical ear has found its object --- the flat note of a vessel that does not answer full, heard from a cup across the birth room and again from a pot in the Vail house --- and remains a prisoner of two paces of trodden ash. He has passed from a house that would not name him into a house that will not ask where his memory came from. He keeps the count --- four, four, a knot, then a gap --- as the one entry in any room that belongs to nobody and cannot be entered on a slate. He cannot speak, cannot read Grathic, cannot reach a cup, and is as lost as the reader.

\CRsheetheading{Chapter Log}
\begin{itemize}[leftmargin=1.4em,itemsep=1pt,topsep=3pt,parsep=0pt]

\item Ch 1: Born in the front room of House Silt on Stone 7 and lived the whole day awake. Held the count of the floorboards, read Dessa’s pulse and sieve-wire scar, matched the pause in the dark hand on the birth-slate to the hand at the flare light and did not put them together, and felt the river stone go cold with nobody else in the room noticing. Wanted the breast more than he wanted the truth. Drooled, could not hold up his head, could not turn it past a point without a small animal sound, could not reach a cup, could not keep the wool on his shoulder. The clinical ear was wide awake and wholly useless.

\item Ch 2: Lived Days 2 and 3. Heard the flat note of a cup across the birth room for the first time and could neither reach the cup nor make anyone tap it; heard Des’s laugh break through in the present tense and go out; watched Alda cross the sign over his flower again without knowing it. Was carried up the Clay Yard in Sella’s bundle, counting her steps, while Alda told Dessa he had gone in the levy and the number was never spoken. Was laid on fresh straw under the Vail roof and given a bed while Tarn was given the loft; read the hollow seat, the salt coat and the Ash ridge as a man-shaped gap, and heard a second flat note from a pot without anyone naming it. The stone was left behind in House Silt; heat came back into his hands and feet; he was warmer and more watched and could say nothing about either.

\end{itemize}
}

\CRcharacter{Deshawn Carter}{the laugh in the car park}{build/portraits/deshawn-carter.pdf}{\CRfield{Also called}{Des, Deshawn “Des” Carter, Des the medic}\CRfield{Role}{past-life}\CRfield{Born}{1978, Chicago, Illinois, Earth}\CRfield{Kind}{human (Earth)}\CRfield{Origin}{the South Side, Chicago, Illinois, Earth}\CRfield{Allegiance}{Emergency Medical Service; the ambulance bay of the same Level I centre}\CRfield{Status}{unknown on Earth; last heard taking a call at 3 a.m. after Dan’s death}\CRfield{First seen}{Chapter 2}}{\CRsheetheading{Appearance}
A large man, comfortable in his own skin, with a shaved head and a gold tooth he had earned in a joke and kept in affection. His hands were quick and calm and could start an IV in the back of an ambulance on a road made of ice. He smelled of cheap coffee, diesel exhaust and the lemon disinfectant the rigs were wiped down with. He had seen more of Chicago at three in the morning than most people see in a lifetime, and it had left him neither cynical nor quiet. He laughed with his whole body, unembarrassed, the kind of laugh that makes a room turn and look and then smile despite itself. When he stopped laughing, which was rare, the silence he left behind was enormous.

\CRsheetheading{Voice \& Manner}
Fast, warm, profane, generous. He makes jokes instead of asking questions, and the joke is always about himself if he can manage it, which is why nobody minds. He has a paramedic’s way of saying the worst thing in a cheerful voice, because the cheerful voice is the handle by which the worst thing can be carried. His ritual with Dan was ten minutes of bad pizza in a car park after a bad shift, saying nothing at all.

\begin{itemize}[leftmargin=1.4em,itemsep=1pt,topsep=3pt,parsep=0pt]

\item “You look like shit, Marchetti. Sit down. Ten minutes.”

\item “I’m not asking, Dan. I’m just saying, you look like a man who’s waiting for the other shoe.”

\item “He told me once, ‘I’m fine.’ I said, ‘Sure, and I’m the Pope.’”

\item “I told her you were working late. What was I supposed to say, Nina? That your father was doing what he always does?”

\item “Ten minutes. Then we talk about nothing. That’s the whole deal.”

\end{itemize}

\CRsheetheading{History}
He was a paramedic on the South Side, the kind who knows which alley is flooded before the dispatcher says it and which apartment buildings have the bad locks. He and Dan met on calls and became friends in the way men of that work become friends: by seeing each other at the worst possible moments and never mentioning them. Des knew Dan was falling apart. He made jokes about it instead of asking. This was not cowardice in his own account; it was the only tool he had. He had watched other men ask, and watched what happened when the answer came. So he kept the car-park ritual alive — the bad pizza, the ten minutes of silence, the easy exit — for years, and believed he was holding his friend together.

He was in the ambulance the night Vesci left a junior alone with a bleeding patient. He saw it, said nothing to Dan at the time, and mentioned it fifteen years later in a way that was half a joke and half an accusation, by which point it was too late to do anything but be angry at the right person in the wrong tense. He was the one the hospital called at 3 a.m. after the expressway. He laughed for exactly one minute before he put the radio down, because that is what the dispatcher needed from him and because he did not yet know what else to do. He was the one who told Nina her father had been working late, and he has never forgiven himself for the kindness of it.

\CRsheetheading{Wounds \& Fears}
His wound is the unasked question. He is the funniest man in the room and the man who cannot say the one true thing. He fears that his laughter is not warmth but a wall, that he has spent twenty years being the man everybody is glad to see and nobody actually talks to. He fears silence in others — it reminds him of the car park, where the ten minutes were the only honest thing in his friendship and they were also a way of never getting to the eleventh minute. He fears that he was an accomplice to Dan’s slow collapse, that he knew and chose the joke. His contradiction is precisely this: the man who brings the light is the man who keeps the dark in place, because a lit room is easier to leave.

\CRsheetheading{Desires}
He wants one more car-park silence. He wants Dan to walk through the door covered in somebody else’s blood, say nothing, and let Des hand him a slice of terrible pizza. He thinks he wants to laugh; what he needs is to learn that showing up is not the same as reaching, and that a man can be loved out of a silence only if somebody is willing to sit through the eleventh minute. In the world Dan entered, Des’s laugh becomes a sound the protagonist reaches for when he cannot reach Nina — the memory of a friendship that failed by being too comfortable, which is its own kind of warmth stolen and never paid back.

\CRsheetheading{Bonds}
\textbf{Daniel Marchetti} — the friend he never asked. Their bond is built on a ritual of not-speaking, and it is the real thing and also a failure. Dan’s silence kept Des from knowing; Des’s jokes kept Dan from confessing. Neither broke it. Des is the man Dan would have called if he had ever learned to call anyone, and the call never came. \emph{Ch 2:} he arrived inside the boy’s head in the present tense and untagged --- car park, slush under the lamps, cold pizza on the console, the gold tooth catching the light, the ten minutes that said nothing and the laugh that filled the space where the eleventh minute would have been --- and then the timed minute on the dispatcher’s radio after the expressway. He went out like a lamp. Dan had no Grathic word for him and the ache went in under the flower.

\textbf{Nina Marchetti} — he drove her home sometimes when she was small and her father was on a shift, and when she asked about him Des said the thing that was easiest. She knows it was a lie. Her tension with him is quiet and permanent: she was the child at the other end of the joke, and she has spent her adult life being the woman who stops waiting — the thing Des’s kindness helped make.

\textbf{Lena Ostrowski} — the other witness. Des and Ostrowski were the two people who saw Dan clearly, and they never once compared notes. There is a mutual wariness between them: she thinks he enables him, he thinks she judges him, and both are right.

\textbf{Paul Vesci} — Des saw what Vesci did to the junior and said nothing for fifteen years. Vesci, in his warm amused voice, called Des “the entertainment.” Des once told Dan, half-joking, that the butcher and the pretty boy were the same man in different coats — a line Dan did not laugh at, which is the point of it.

\CRsheetheading{Arc So Far}
He arrives in Chapter 2 as a remembered laugh and a car-park silence, a `**` voice that comes up under the boy in the birth room of House Silt while a cup across the room gives its flat note and cannot be reached. The car park is wet with slush under the lamps, the dashboard clock past three, the pizza cold in a paper bag on the console; the rig smells of diesel and lemon disinfectant and cheap coffee. He says the line about looking like shit and sitting down for ten minutes, and then he laughs whole from low in his chest, on through the space where the eleventh minute would have been, and later on the radio for exactly one minute, timed by a man who had nothing else to give. The whole arrival is in the present tense, untagged, and goes out like a lamp put out. In a country with no word for a paramedic he is an ache under the ribs where the flower is.

\CRsheetheading{Chapter Log}
\begin{itemize}[leftmargin=1.4em,itemsep=1pt,topsep=3pt,parsep=0pt]

\item Ch 1: Does not appear; the car park is held for Chapter 2.

\item Ch 2: Came back unbidden on Day 2 in the birth room: car park, bad pizza, ten minutes of silence, the gold tooth in the lamplight, the diesel-and-lemon smell of the rig. Said \emph{You look like shit, Marchetti. Sit down. Ten minutes.} and laughed for the space of the eleventh minute and then for one minute on the radio after the expressway. Vanished with the car park; did not know he had been heard.

\end{itemize}
}

\CRcharacter{Nina Marchetti}{the light she stopped leaving on}{build/portraits/nina-marchetti.pdf}{\CRfield{Also called}{Nina, N., Nina Marchetti}\CRfield{Role}{past-life}\CRfield{Born}{2000, Chicago, Illinois, Earth}\CRfield{Kind}{human (Earth)}\CRfield{Origin}{Elmwood Park and Hermosa, Chicago, Illinois, Earth}\CRfield{Allegiance}{no house; a grown woman who worked with animals outside the city, then moved}\CRfield{Status}{alive on Earth, last known; unknown to Sarn}\CRfield{First seen}{Chapter 1}}{\CRsheetheading{Appearance}
In the photograph on his wall she is six years old, missing two front teeth, squinting in summer light on the edge of Lake Michigan, her whole face given over to the brightness. That is the image he kept and looked at every night without seeing. In life she grew into a woman with her mother’s patience and her father’s habit of not asking: a steady, unshowy person, practical with her hands, who could settle a frightened animal without a word and could not settle a conversation with a man who was good at leaving. She has a way of looking at a person while she decides whether to say the thing that hurts, and then saying it in a light voice, which is a kindness of a sort and a cruelty of a sort. She does not wait well, and she has spent her life being told by other people that this is a virtue.

\CRsheetheading{Voice \& Manner}
Warm, careful, slightly dry. She says the difficult thing casually, so that the other person can pretend not to hear it. With her father the sentences get shorter over the years, until there is almost nothing left to say on the phone. She never raises her voice; her silences are louder.

\begin{itemize}[leftmargin=1.4em,itemsep=1pt,topsep=3pt,parsep=0pt]

\item “Don’t worry. I’m not waiting up.”

\item “You always say that. I’ve stopped counting how many times.”

\item “I know. It’s just a boat, Dad. It was always just a boat.”

\item “I stopped leaving the light on. You didn’t even notice. That’s the part I can’t forgive.”

\item “I called. You didn’t pick up. I know you were working. I know. I just wanted you to know that I know.”

\end{itemize}

\CRsheetheading{History}
She was born in 2000, in the years when her father was a junior at Cook County and the house in Hermosa was loud. By nine she had learned not to ask when he would be home, which is a thing a child learns by asking once too often and getting a kind answer that means nothing. By twelve she had stopped leaving a light on for him. She did not hold this against him then; she held it against herself, as if the light had been a claim she had no right to make. She went to college, worked with animals somewhere outside the city, then moved further away. They spoke on birthdays and not otherwise. He kept meaning to call. He was always meaning to call.

The promise was a boat on Lake Michigan, made when she was small, renewed every birthday for nine years, never delivered. It was the first line of the call he composed in his head on the drive home the night he died: something about the lake in the summer and a boat, a sentence he had rehearsed so often that in the moment of breaking it had no words left, only weight. That night she called him. The screen lit in his coat pocket with a name he would not look at. He did not pick up. She has the missed call on her own phone still, and has never deleted it, and has never told anyone she has it.

She grew into the woman her mother became: one who stops waiting, who calls it peace. When the news came she cried for one evening and then was fine, and the being fine frightened her more than the crying. Des Carter told her, gently, that her father had been working late; it was the last lie anyone told her about him and she has never quite forgiven the kindness of it. Ostrowski sent a card that arrive late and said, in a clipped hand, that Daniel had been a good surgeon. Nina kept the card. She resents it, because it is a better account of her father than she has.

\CRsheetheading{Wounds \& Fears}
Her wound is the ordinary one that nobody names: a parent who was present in the room and absent everywhere else. She fears she has become her mother — a woman whose great achievement is that she no longer waits — and that she will pass this on, that she will have a child who learns not to ask. She fears that her father died believing she had stopped caring, when the truth is she called him that night and he let it ring. She fears the missed call more than the death: the death was an event, the call is a habit, and habits are harder to bury.

Her contradiction is that she is the woman who stopped waiting and the woman who kept the receipt. She tells everyone she is over him; she keeps the photograph of the boat in a drawer, and the nine-year-old promise, and the missed call, and a card from a surgeon she never met. She believes she has released him. She has only stopped saying it aloud.

\CRsheetheading{Desires}
She wants to be a person whose life is not organised around a man who might call. She thinks she wants an apology — one sentence, from him, in which he says he is sorry and means it. What she needs is to understand that the debt was never hers to carry: his silence was his, his guilt was his, his death was his, and the call she made that night was an act of love, not a failure. In the terms of the world he entered, she is the weight that crossed with him, the unsent call made solid — and she is not the child in the minivan. That other woman, the one he pulled from the van, carries his last warmth. Nina carries his last unsent words. He never named either of them in Sarn.

\CRsheetheading{Bonds}
\textbf{Daniel Marchetti} — her father. Love and resentment braided so tightly that she cannot pull one thread without the other. He is the man who promised a boat and never bought it, who died with her name lit on his phone, whose greatest act was to give his warmth to a stranger’s child. She would trade all of it for one honest conversation. He carries her into Sarn as a debt he cannot settle, and reaches for a phone for her every time somebody dies. \emph{Ch 1:} she is a lit rectangle under a coat flap on the expressway and the nine-year promise gone down to weight; he does not look at the name.

\textbf{Elena Rossi} — her mother, the woman who stopped waiting first. Nina learned the art from her and now fears she has perfected it. They love each other in the practical Marchetti way: by doing things and refusing to discuss them. Mother and daughter have never once said out loud what her father’s absence did to them.

\textbf{Deshawn Carter} — the family friend who told her, kindly, that her father was working late. She knows it was a lie. Her tension with Des is that he was the adult who chose the comfortable sentence, and then laughed about it with her father in the car park for years; she was the child at the other end of both.

\textbf{Lena Ostrowski} — the surgeon who sent the late card. Nina resents her with a cold, precise resentment: this woman had a version of her father that she never had, got the full name spoken in a voice like a file, and was thanked with flowers. Nina kept the card because it is the only account of him that does not require her to supply the feeling herself.

\textbf{The woman in the minivan} — a stranger Nina has never heard of, who was nine or ten the night her father died and is now a woman still warm on Earth. Nina does not know that the last warmth of her father’s life went into that stranger, not into her. When she learns it, if she ever does, it will be the final account, and it cannot be settled by dying again.

\CRsheetheading{Arc So Far}
Chapter 1 gives her the shape she will keep for thousands of chapters: a light under a coat flap and the weight of a sentence never finished. She has no words in the book yet and already carries the debt that made a man portable. The missed call on her own phone, the boat photograph in the drawer and the late card from Ostrowski are still offstage; the reader knows her only as the thing Dan refuses to look at.

\CRsheetheading{Chapter Log}
\begin{itemize}[leftmargin=1.4em,itemsep=1pt,topsep=3pt,parsep=0pt]

\item Ch 1: Present only as the lit screen in his coat pocket on the expressway --- a name he will not look at, heavier than the phone --- and as the boat promised for nine years whose sentence has nothing left in it but weight. Not named in the Sarn narration; no line of her own.

\end{itemize}
}

\CRcharacter{Dessa Silt}{the mother who counted nineteen steps and called the child dead}{}{\CRfield{Also called}{Dessa, Dessa of the Slip Pits, the little hot one’s mother}\CRfield{Role}{family}\CRfield{Born}{Y.H. 788, Low Kiln, Sill, Grath}\CRfield{Kind}{human (Grathic; returned thin)}\CRfield{Origin}{Low Kiln, Sill, Grath; House Silt}\CRfield{Allegiance}{the slip pits, Clay Yard, Low Kiln}\CRfield{Status}{alive at the close of Day 3; believes her child went in the levy among the Forty and will believe it three years; dies Day 38, Ilmen 9, Y.H. 813, drowned in the Slip Brook, uncollected}\CRfield{First seen}{Chapter 1}}{\CRsheetheading{Appearance}
Twenty-four and worked older: a thin woman with slip-grey nails and a cough that comes up from the pits in cold weather, which is most of the year. She has a way of smiling when she is angry that her sister Mira calls \emph{the door shut}, because the smile means nothing further will be said. Her left forearm carries a long pale scar from a broken sieve-wire, and she rubs it when she is deciding whether to tell the truth. She smells of wet clay, sour milk and, on the days she has been at the pits, the mineral edge of water from the Lead Run that nobody in Low Kiln will drink twice. Her hands are small and quick and never clean. The thing people notice first is the smile; the thing they notice last is that she counts under her breath when she walks — the nineteen steps from the slip pit to the drying rack, which she has counted every day for six years.

\CRsheetheading{Voice \& Manner}
Short, practical, with a joke that costs her. She says the difficult thing in the proverb-form the district uses for saying anything at all, so that it can be taken as advice if the listener prefers. She never calls her child by his name; she calls him \emph{the little hot one} or \emph{the flower}, and has done this so long that the habit has become a way of not saying what she means. When she is frightened she gets specific — dates, steps, the number of loaves — because specifics are a handle on a thing that is otherwise too large.

\begin{itemize}[leftmargin=1.4em,itemsep=1pt,topsep=3pt,parsep=0pt]

\item “Nineteen steps. I counted them when I was waiting for him to come back. He didn’t come back.”

\item “A body can be warm and still be gone. You learn that in the pits.”

\item “I didn’t give him away. I gave him heat. There’s a difference, and I’ll hear it from a collector if you like.”

\item “Don’t set that near the fire. It’ll crack, and then you’ll have a crack and a pot.”

\item “He’s not a Vail. Don’t let anybody write him as one and call it kindness.”

\end{itemize}

\CRsheetheading{History}
She was born in Y.H. 788 to a mother who was banked and returned thin, which means Dessa came into the world with a faint interest in a life she cannot name. She learned early that this is considered an advantage in Low Kiln, where a person with some of somebody else left in her is a person who can work two shifts. She works the slip pits for six bits a day, shares a room with her sister Mira, and at nineteen had a child by a quarry hand from Stonemouth named Nils, who went upriver in the autumn and never sent for her. She carried the child through the Long Cold with nothing but a stone warmed at the Vail fire and the kindness of Alda Reed, who attended the birth on Stone 7, Y.H. 812, and crossed herself at the branching red flower across the baby’s ribs.

On the third day Sella Vail took him. Alda told Dessa the boy had been entered to a house and then, when Dessa asked which, said the thing that has been a lie for three years: that the child had gone in the levy, that he was among the forty children of Low Kiln banked that month for a private account, and that a mother with a thin return would only draw attention to a slate that was better left alone. Dessa believed her. She has spent three years believing her son was spent to buy somebody else’s wholeness, and has never once been able to say the number forty without touching the pale scar on her forearm.

\CRsheetheading{Wounds \& Fears}
A mother who returned thin and reached for a life that was not hers; a child she believes was banked; the nineteen steps she counts because counting keeps her from thinking. By the end of Day 1 the count has stopped and her pulse has taken it up; the cough is still called weather. The sieve-wire scar is rubbed whenever a truth is near. She fears that if she ever learned the boy was alive she would want him back, and that wanting would make her a thief in a district where a woman with no house can be accused of anything. She fears the collectors, the name-pot, and the Second Want in herself — the shape that sometimes comes over her in the night and leaves her reaching for a cradle that is not in the room. Her contradiction: she tells everyone she has stopped thinking about the little hot one, and has kept a scrap of cloth from his birth in a tin above the bed for three years.

\CRsheetheading{Desires}
She wants to be warm once without anyone entering it against her account. She thinks she wants Nils to come back, though she would not take him if he did. What she needs is to know that the boy lived; and beyond that, to understand that the shape in her that reaches for him is not the defect of a thin return but her own love, which is the one thing the Court has never learned to deduct.

\CRsheetheading{Bonds}
\textbf{Corin} — her son, who lies under another roof and does not know her. She calls him the flower when she thinks nobody is listening, and once, at a market, stood behind him for a quarter of an hour while he slept in a basket and did not touch him. \emph{Ch 1:} she pressed the grey wool flat at his shoulder with two fingers that stayed longer than the cloth required, and would not say his name. A thing with no name cannot be claimed by a woman with no house. \emph{Ch 2:} she followed the bundle as far as the end of the yard and did not take one of the nineteen steps to come after it. She never said the name; she counted nineteen out and nineteen back instead, and the count went from a rhythm to a straight line and then stopped.

\textbf{Sella Vail} — the woman who took him and paid her in barley and slip-clay to stay away. Dessa has taken the barley every winter and hated herself for it each time, because taking it is a kind of agreement. \emph{Ch 1:} her pulse came up at her wrist under the \emph{door shut} smile while the other woman’s hand hovered above the bundle; she did not look up. \emph{Ch 2:} the taking happened and neither woman looked at the other. Dessa stood where the houses stop and the rise toward the wheel-shed begins, pulse in throat, thumb on scar, the door-shut smile, and let the boy go up the yard under the other woman’s coat.

\textbf{Alda Reed} — who lied to her and has not slept well since. They avoid each other in the street; when they must speak, they speak about the weather, which in Low Kiln is never the subject. \emph{Ch 1:} Alda delivered the child and met the room’s unasked question with a year and the weather; Dessa answered with loaves and steps and the day of the month. \emph{Ch 2:} Alda came up behind her with her hands smoothing her apron in two passes and said the boy was entered in the levy. Dessa asked which, got the exact dates, and asked nothing further. She believes it. The figure was never spoken.

\textbf{Mira Silt} — her sister, her room-mate, the person who knows the boy is alive and has not said so. Dessa suspects it. She has never asked, because an answer would be a thing she would have to carry, and she is already carrying one. \emph{Ch 1:} the crust of brown-edged barley came across the floor from her sister, brisk and Vail-smoked, and she took it when it was pressed into her hand. \emph{Ch 2:} Mira is not on stage; the last warmth that crossed this room from her hands is still the crust, and the boy is gone.

\textbf{Berr Clay, the Thane of Clay} — took the levy from her wages in Y.H. 811 and called it one in twelve. She spat at his door in Grain and denied it when he asked, which is the whole of their relationship.

\textbf{Nils of Stonemouth} — the quarry hand who left before the child was born. She has not heard from him, and has begun to hope she never does, because a man who returns after three years is a man who will ask questions about a boy.

\CRsheetheading{Arc So Far}
On Day 1 she is the woman in the room whose count is not about the pits. She holds the thing she will not hold by naming dates and loaves and steps instead of her child; when Sella’s hand hovers, she smiles the door shut and looks away. On Day 2 she coughs from the pits and calls it weather, rubs the sieve-wire scar, and says a body can be warm and still be gone --- in the proverb-form, where it can be taken as advice. On the third morning she follows the bundle to the end of the yard and stops; her count goes from a rhythm to a straight line and then to nothing. Alda tells her the boy went in the levy and gives dates faster than dates have any business arriving, and she believes it because disbelief would be a thing she would have to carry. The reader is meant to take her for a bystander and to be wrong for three years. Already on the page are the nineteen steps, the pale scar, the door-shut smile, the little hot one, and a crust from the fire of the house that took him.

\CRsheetheading{Chapter Log}
\begin{itemize}[leftmargin=1.4em,itemsep=1pt,topsep=3pt,parsep=0pt]

\item Ch 1: Lay on the pallet through the birth and the night, face turned toward the wall where the room allowed it. Counted nineteen out and nineteen back until the count stopped and her pulse took it up. Rubbed the sieve-wire scar while deciding how much to say; got specific --- nineteen steps twice daily since Y.H. 806, four loaves baked in Grain, three to the levy and one eaten with her feet up, the seventh of Stone in eight hundred and twelve --- and would not name the child. Called him the little hot one and the flower. Pressed the wool flat at his shoulder with two fingers that stayed longer than the cloth required. Smiled the \emph{door shut} while Sella Vail’s hand hovered. Took Mira’s crust of Vail-smoked barley bread.

\item Ch 2: On Day 2 coughed from the pits, called it weather, rubbed the pale sieve-wire scar and got specific about steps, loaves and the day of the month; said a body can be warm and still be gone without looking at the child. On the third morning got up while Sella was at the bench, followed the covered bundle to the end of the Clay Yard and took none of the nineteen steps after it. Pulse in throat, thumb on scar, the door-shut smile; said nothing and did not say his name. Heard Alda say he had gone in the levy among the children taken for the account, asked only \emph{where}, and asked nothing further as the dates went on arriving. She believes him spent and will believe it three years.

\end{itemize}
}

\CRcharacter{Garl Vail}{the man whose warmth was left where it lay}{}{\CRfield{Also called}{Garl, Garl of the Vail, the dead wheel-keeper, the uncollected}\CRfield{Role}{family}\CRfield{Born}{Y.H. 775, Low Kiln, Sill, Grath}\CRfield{Kind}{human (Grathic)}\CRfield{Origin}{Low Kiln, Sill, Grath; fourth generation of House Vail}\CRfield{Allegiance}{House Vail kiln-house; Kiln-Masters of Sill}\CRfield{Status}{dead (Stone, Y.H. 811; uncollected in the Long Cold; warmth lost to the Waste)}\CRfield{First seen}{Chapter 2}}{\CRsheetheading{Appearance}
He is known in the house by objects rather than by memory. The wheel in the shed has a seat worn hollow on the left side where he sat for twenty years, and Sella will not let anyone fill it. A coat of his hangs on a peg behind the door with the salt stains still in it, and she has never taken it down, though Jorra has twice reached for it in the dark and stopped. He was a big man with clay in the creases of his hands and a laugh that ended in a cough; his left thumb was shaped by the wheel into a permanent hook. People who saw him said he moved slowly and arrived early, which in a kiln-house is the same thing. He smelled of wet clay and banked coals. In the house the absent smell is stronger than the memory: when the kiln is cold and the Ash ridge is thick, the room smells the way the corner of a room smells where something was left too long.

\CRsheetheading{Voice \& Manner}
He is remembered in proverbs, which means his voice has been improved by use and is probably not as good as the versions of it that survive. He never said a hard word in front of the fire — a habit Jorra has inherited without knowing it — and when he was angry he went to the wheel. His sayings are all about repair and none of them are about completion.

\begin{itemize}[leftmargin=1.4em,itemsep=1pt,topsep=3pt,parsep=0pt]

\item “If it’s worth mending it’s worth mending twice. The second time is the real one.”

\item “Don’t tell me the pot’s ruined. Tell me what it was before it was ruined.”

\item “Feed the fire before you feed yourself. It’ll still be there in the morning. You won’t.”

\item “A man’s hands know what his head’s still arguing about.”

\item “Leave the crack showing. People trust a thing they can see has been handled.”

\end{itemize}

\CRsheetheading{History}
Fourth generation of Vails, born in the kiln-house and married to Sella, a slip-mixer’s daughter, in Y.H. 796. He kept the wheel for twenty years, threw the tiles that went into the floor of the Cross market, and was offered a price for his measure by a Sere factor in Y.H. 809 which he refused and did not mention for a year. He had two children: Jorra, born Y.H. 796, and Enna, born Y.H. 802, who died at nine months of the second chill and whose death he sat through without being able to do anything that mattered. In Y.H. 810 he began to be what Low Kiln calls walking thin — a vessel leaking faster than a man of his trade should leak — and he said nothing about it, because a kiln-master who admits a flaw loses the wheel.

In the Long Cold of Y.H. 811 the Ilm lay under ice for ninety-four days and no collector could cross from Greybank. Garl’s vessel gave in the first weeks of Stone. He died at home with Sella beside him and no iron tube in the room, and his warmth went where warmth goes when nobody takes it: into the Waste, unbanked, unrecorded, unreturned. In the morning there was grey Ash along the skirting board thicker than anyone had seen in that house. Tarn was sent to find a collector and turned back at the frozen ford. Jorra was thirteen and took the wheel a week later. The Vail household has been cool ever since, in a way that has nothing to do with the weather.

\CRsheetheading{Wounds \& Fears}
He knew he was leaking and kept throwing. That is the wound and it is also the answer to any question about why he is dead. He feared being banked thin and coming back as someone who would stand in his own house and not know the wheel; he feared more than that being spent whole and returned beautiful for somebody else’s benefit, which is what he had seen happen to a man he would not name. His contradiction: a man who mended pots for a living would not let anyone mend him, and called the refusal \emph{not wanting to be a bother}, which is the gentlest lie in the district and the most expensive.

\CRsheetheading{Desires}
He wanted to leave the wheel to Jorra whole, with the seat still worn and the note still true, and to come back as himself one day and teach the boy what a full note feels like under the knuckle. He thinks he wanted a whole return for himself. What he needed was to understand, before he died, that a vessel can be repaired without first being emptied — the thing his son’s household’s foundling will one day prove in the face of every priest in Sarn. He did not live long enough to be wrong about it in an interesting way, which is why his absence still has work to do.

\CRsheetheading{Bonds}
\textbf{Sella Vail} — his wife. She decided he would not be collected; he did not ask her to decide. A promise unkept gathers Ash in Sarn as surely as a debt, and the ridge in their room is the shape of theirs. \emph{Ch 2:} she touches the widow’s mark before the kiln and at no other time; she will not let anyone fill his seat or take his coat down, and the house keeps a private roll in which he is entered and which it does not show to a man from Slate Row.

\textbf{Jorra Vail} — his son, who has his hands and not his patience. Garl never told the boy he was proud of him. The boy has spent three years building a life around the absent sentence. \emph{Ch 2:} Jorra sits in the hollow seat he cannot fill and goes to the wheel when he is angry, which is the habit Garl had and never explained.

\textbf{Corin} — he never met the boy, who now sleeps under his roof and uses a warm stone that was warmed at Garl’s fire. There is a late recognition in this: the dead man’s hearth warmed the child he would take in, and a promise that cannot be settled because it was never made. \emph{Ch 2:} the boy read him as a man-shaped gap --- four, four and a space where the fifth thing should be --- and saw the charred edge of a slate in the kiln ash without knowing whose hand had written on it.

\textbf{Berr Clay, the Thane of Clay} — owed Garl for a consignment of tiles and never entered it. Garl never called it in. The Vail house has been short of coin ever since, and the debt has gathered interest in the Court’s books even though nobody in Low Kiln can read them.

\textbf{Enna} — his daughter, who died at nine months and whose second chill he could not stop. He has been uncollected for a year; she has been returned thin or not returned at all, which nobody in the house will say aloud. Both are entered on the house’s private roll.

\CRsheetheading{Arc So Far}
He is an absence with objects attached, and in Chapter 2 the objects are given a room. The kiln breathes in fours with a gap where the fifth breath should be; the wheel-seat is worn hollow on the left by twenty years of leaning; a salt-stiff coat keeps a broader man’s shape on its peg behind the door; the Ash ridge runs the length of the skirting board and turns the corner without a break, near two thumbs wide, and nobody sweeps it before the Last Day. His sayings have outlived him and improved in the keeping, and the house says them to the pot and the fire and never to a person. In the kiln ash lies the charred edge of a slate with the writing gone off the end, swept around for a year and never picked up. The reader knows him by what the household does around the hole before anyone says how the hole was made.

\CRsheetheading{Chapter Log}
\begin{itemize}[leftmargin=1.4em,itemsep=1pt,topsep=3pt,parsep=0pt]

\item Ch 1: Does not appear; his name, the stone warmed at his hearth and the house’s silence around him are held for Chapter 2.

\item Ch 2: Present throughout Days 2 and 3 as the shape of the house. The kiln’s breath he kept alive was counted in fours and a gap; his seat stayed hollow, his coat stayed on its peg, his Ash ridge grew, and his proverb \emph{leave the crack showing} was said in the house’s voice to the pot and the fire. The house’s private roll was spoken around: he was entered on it, and the foundling’s name would be entered there in a careful hand. The charred edge of a slate lay in the kiln ash and was not examined. No collector was called.

\end{itemize}
}

\CRcharacter{Jorra Vail}{the wheel-keeper who has never said which side he is on}{}{\CRfield{Also called}{Jorra, Jorra of the Vail, the wheel-keeper}\CRfield{Role}{family}\CRfield{Born}{Y.H. 796, Low Kiln, Sill, Grath}\CRfield{Kind}{human (Grathic)}\CRfield{Origin}{Low Kiln, Sill, Grath; House Vail}\CRfield{Allegiance}{House Vail kiln-house; Kiln-Masters of Sill; quietly courted by the Broken Slates}\CRfield{Status}{alive}\CRfield{First seen}{Chapter 2}}{\CRsheetheading{Appearance}
Sixteen and built like a gate hinge: tall for a clay boy, narrow, with slip dried to the colour of old bone up both forearms and a permanent bend in his left little finger from a splinter off the wheel when he was eleven. He cuts his hair short so Ash cannot nest in it, which is a thing the district notices and he never explains. His eyes do not finish sentences; they go to the door, to the fire, to the wheel, and come back without reporting. He smells of damp clay, tallow and the sweetish edge of lead glaze carried up from High Kiln on a bad wind. The thing people notice first is his hands, which are the hands of a thrower and have been working since he was nine. The thing they notice last is that he never looks at Sella when she speaks to him, though he obeys her instantly.

\CRsheetheading{Voice \& Manner}
Short, flat, declarative. He asks almost no questions, and when he does it is the tool rather than the clay: \emph{Is it dry? Did you fire it? Who wrote it?} He speaks to the wheel and to the pots as other men speak to animals, and only then to people. When he is angry he gets quieter, the way Sella does, and the household has learned to read the lowering of his voice as a weather signal. He never says \emph{with} or \emph{against} when \emph{I haven’t said} will do.

\begin{itemize}[leftmargin=1.4em,itemsep=1pt,topsep=3pt,parsep=0pt]

\item “The wheel doesn’t care whose hands. It cares if they’re steady.”

\item “You want me to say it. I’m not saying it.”

\item “Sella. Don’t feed him from your own dish. He’s not Vail.”

\item “A slate burns in four minutes. Four. I’ve timed it.”

\item “I didn’t say I was with them. I said I wasn’t against them. That’s different. In this district it isn’t, but it’s what I said.”

\item “He’s not Vail.” --- said in the kiln-house at first light, refused an answer, and sealed by the boy’s arrival.

\item “You’ll do.” --- to the pot on the wheel, in a voice with no welcome in it and no unkindness either.

\end{itemize}

\CRsheetheading{History}
He is the son of Garl and Sella Vail, and he took the wheel at thirteen when Garl died in the Long Cold and did not come back and nobody came to draw him. The wheel had been Garl’s for twenty years and the seat was worn hollow; Jorra has spent three years trying to fill that hollow with his own weight. He has thrown pots since he was nine, fired them since eleven, and knows the full note of a Vail pot before it leaves his hands, though he will not say so in front of the kiln-masters, who would call it early and make something of it. His house has remembered the Ash Rising since Y.H. 668, when children were taken and the rising was put down by a man with a list; the memory is kept not in stories but in the way certain doors are never closed and certain names are never entered. Since he was fourteen men from the Broken Slates have come into the back room of the kiln-house and spoken in low voices. He has listened to all of them and told none of them yes.

He watched Enna, his sister, die of the second chill at four, and has never forgiven the fire for being warm while she cooled. He watched his father die and was sent to find a collector and could not cross the ford, and has never asked whether he was sent too late or whether he was sent at all. When Corin arrived he said one sentence — \emph{He’s not Vail} — and has said nothing about it since, which in that house is a form of argument.

\CRsheetheading{Wounds \& Fears}
A father who went uncollected, a sister who went cold in a warm room, and a wheel he was given before he was ready. He fears Sella will die with something unsaid and leave him holding the house and the lie together. He fears the Broken Slates will decide for him what he is, because he has not decided and the longer he waits the more their decision looks like his own. He fears the boy most of all, because the boy hears things. His contradiction: he is the most careful keeper of the Vail fire in four generations and he is quietly helping the people who want to burn every slate in Sill. He has never once thought these two facts cannot live in the same man.

\CRsheetheading{Desires}
He wants the wheel to stay in the house and the fire to hold and nobody to write the boy down as Vail without asking him first. He thinks he wants Sella to say that he has done well with the wheel. What he needs is to learn that burning a slate does not return a person to you, and that the anger he has been nursing for three years is not his alone — it has been lent to him by men who intend to collect on it.

\CRsheetheading{Bonds}
\textbf{Sella Vail} — his mother. He believes she took Corin in to replace Enna and to keep the house’s standing after Garl died. She has never contradicted him. Every obedience he gives her is therefore double: it is love and it is a correction neither of them will make. \emph{Ch 2:} she announced the taking and he answered \emph{he’s not Vail}; she got quieter and gave him no answer at all. He went on working, did not look at her, and shifted the damper before she had finished telling him to.

\textbf{Corin} — the foundling. Jorra was the first in the house to notice that the boy turns his head toward a pot before anyone taps it, and that he does it to hear the note, not the noise. He has told nobody. He resents the boy for being given a bed in the house while the clay-bound hand sleeps in the loft, and he cannot stop watching him. \emph{Ch 2:} the foundling arrived under the roof and nobody wrote him as Vail, which is what Jorra wanted and what he cannot bear. He spoke to the wheel, to the pots and once to Tarn about a rasp out of its place, and never addressed the boy.

\textbf{Alda Reed} — attended Enna’s death and Corin’s birth, and has never met Jorra’s eye for longer than a moment. He knows she wrote something down the night Enna died and has never asked to see it, because asking would make it a thing between them.

\textbf{Tarn} — the Vail’s clay-bound hand. Jorra treats him as a brother when the wheel needs him and as a servant when it does not. Tarn has decided to hate him for it and keeps failing, because Jorra once sat up all night with him through a fever and has never mentioned it. \emph{Ch 2:} Tarn was given the loft while the foundling was given the bed, and said so in one sentence to Sella. Jorra said nothing; he has never defended the household’s arithmetic and never contests it.

\textbf{Berr Clay, the Thane of Clay} — Jorra threw a stone at his door in Y.H. 810 and denied it. The Thane suspects the boy is with the Broken Slates and lets him be, which Jorra finds more frightening than being watched.

\CRsheetheading{Arc So Far}
He enters at the wheel with his weight in the hollow his father wore and has spent three years failing to fill --- slip dried to the colour of old bone up both forearms, the left little finger bent from a splinter at eleven, a story about a fight he tells because a fight sounds more like him. In Chapter 2 he says one flat sentence about the foundling and is refused an answer, and the sentence is sealed by the boy’s arrival: a child lies under the roof and nobody has written him as Vail. He does not look at his mother while obeying her instantly, speaks to the pots before people, goes quiet when the levy is spoken of, and has given none of the men who come into the back room a word they can carry away. He is the most careful keeper of the Vail fire in four generations and is quietly helping the people who want to burn every slate in Sill, and has never once thought the two facts cannot live in the same man.

\CRsheetheading{Chapter Log}
\begin{itemize}[leftmargin=1.4em,itemsep=1pt,topsep=3pt,parsep=0pt]

\item Ch 1: Does not appear; the wheel, the hollow seat and his objecting are held for Chapter 2.

\item Ch 2: At first light stood at the wheel with slip-coloured forearms and the bent finger, told the pot the slip was too thick and it would crack at the shoulder, and went quiet when the levy was spoken of. When Sella said the boy would come under the roof he said \emph{he’s not Vail}, did not look at her, and was not answered. On the third day sat in the hollow seat as the boy was carried in, obeyed Sella’s instruction about the damper before she had finished, told a pot \emph{you’ll do}, spoke once to Tarn about a rasp, and never addressed the child. Nobody wrote the foundling as Vail.

\end{itemize}
}

\CRcharacter{Mira Silt}{the aunt who knows and will not say}{}{\CRfield{Also called}{Mira, Mira of the Fishstairs, Dessa’s sister}\CRfield{Role}{family}\CRfield{Born}{Y.H. 785, Low Kiln, Sill, Grath}\CRfield{Kind}{human (Grathic; returned thin)}\CRfield{Origin}{Low Kiln, Sill, Grath; House Silt}\CRfield{Allegiance}{the Fishstairs, Sill; eel trade; House Silt}\CRfield{Status}{alive}\CRfield{First seen}{Chapter 1}}{\CRsheetheading{Appearance}
Twenty-seven, broad-hipped, with eel slime dried on her sleeves and a small smile she keeps on her face the way another woman keeps a door on the latch: it can be shut without anybody noticing the movement. The top half of her left ear is missing, taken by a knife in a Fishstairs quarrel when she was nineteen; she says a hook did it, and the scar says otherwise, and nobody has ever made anything of it. She smells of river silt, salted perch, tallow smoke and the cold off the water. She carries a string of eel-bones in her pocket and flicks them against her thumb when she is deciding whether to lie, which is often. The thing people notice first is the smile; the thing they notice last is that she never sits down at a table where a slate is open.

\CRsheetheading{Voice \& Manner}
Fast, low, with the proverb turned sideways. She answers a question with a question and then answers it herself, which gives her the room to choose which answer the listener gets. She never tells a story with a beginning; she starts at the part where somebody is already in trouble. She has learned from Fen traders to mark accounts in knots rather than on slate, and will not write a name down if a knot will hold it. When she is being kind she becomes brisk, because kindness in her house is a thing done quickly before anyone can make a speech about it.

\begin{itemize}[leftmargin=1.4em,itemsep=1pt,topsep=3pt,parsep=0pt]

\item “You want to know. You don’t want to know. I can see which one’s winning.”

\item “The secret’s in what happens to it after. I didn’t tell anybody anything. There’s nothing to tell. That’s the whole of it.”

\item “It came from a warm house, that’s all. There’s barley in it and a little smoke and it’s a crust, not a sermon.”

\item “Nineteen steps. Your sister counts them. I count the ones she doesn’t take.”

\item “Don’t ask a Fishstairs woman what she saw. She’ll sell it to you and then it’s yours and you’ll have to carry it.”

\item “He’s warm. That’s all I’ll say on a slate, and I’m saying it off the slate.”

\item “I kept your secret for three years, Dessa. Don’t make it a fourth.”

\end{itemize}

\CRsheetheading{History}
Born in Low Kiln to the same thin-returned mother as Dessa, Mira learned the eel trade at the Fishstairs at twelve and has worked the dawn market there ever since. She knows which households pay and which owe, which collectors take their ale before a collection and which after, and which boats on the Ilm are carrying something other than fish. She has been the district’s informal ledger of who is warm and who is cooling for fifteen years, and has never once put a name on paper.

She was in the room on Stone 7, Y.H. 812, when her sister bore the child with the branching red flower. She saw Sella Vail carry a covered bundle up the Clay Yard three days later, and she understood what had been done without being told. Alda Reed told Dessa the boy was among the Forty; Mira let her believe it. She has let her believe it for three years, and has taken a share of the barley Sella pays every winter, which she spends on Dessa, who does not know where the barley comes from. She has a child of her own, returned thin in Y.H. 802 to a household in Nine Pots, who does not recognize her; she has walked to Nine Pots twice and never knocked.

\CRsheetheading{Wounds \& Fears}
Her sister’s lie is hers too, and she has begun to feel the weight of it in the hand that holds the eel knife. She fears that telling Dessa the truth will put the boy on a slate that someone can use against him; she fears more that if she waits much longer the truth will come out in a quarrel, which is the worst way for a thing like that to be said. She fears she has become a woman who trades in other people’s warmth and calls it looking after them. Her contradiction: she gossips all day at the Fishstairs and has kept the one secret that matters for three years without a single slip.

\CRsheetheading{Desires}
She wants Dessa to stop reaching for a child she believes is dead, and she wants to be believed once about something that is not a rumour. She thinks she wants the truth to come out, and has arranged her life so that it never has to. What she needs is to understand that a secret kept out of love is still a debt, and that it gathers Ash like any other debt, and that the only way to stop paying interest on it is to say the thing.

\CRsheetheading{Bonds}
\textbf{Dessa Silt} — her sister, her room-mate, the woman whose child she has hidden for three years. Dessa suspects; Mira has begun to be afraid of the moment Dessa stops suspecting and knows. Everything between them is said in the language of barley, coughs and the nineteen steps. \emph{Ch 1:} she crossed the room and pressed a crust of Vail-smoked barley into Dessa’s hand in two brisk movements, stopping short of the slate to do it.

\textbf{Corin} — the nephew she watches from a distance at the market and once followed as far as the Clay Yard. She has never called him anything, not even in her head, because a name is a thing that can be overheard. \emph{Ch 1:} she would not hold him and would not sit where the slate lay open, but pushed the wool back under him briskly with two movements before anyone could make a speech about it.

\textbf{Sella Vail} — they have an arrangement: barley and silence. Mira has begun to resent the barley and has never once refused it, which is the precise shape of her complicity. \emph{Ch 1:} she watched the hand in the door without moving, her thumb quiet in her pocket and the eel-bone put away.

\textbf{Alda Reed} — they trade small facts and never the large one. Alda knows more about the birth; Mira knows more about what happened after it. Neither will sit at a table with a slate on it, and they have built a friendship entirely out of things said standing up. \emph{Ch 1:} they worked the same room standing, each with her back to the writing.

\textbf{Orra Weir} — her supplier at the Weir and her oldest friend, who wants the secret out and has said so twice. Mira has answered both times with the weather.

\textbf{Berr Clay, the Thane of Clay} — buys her eels for the levy feast and pays in marks that are short by two. She short-weights him by two eels. This small war has been going on for six years and neither has acknowledged it.

\textbf{Tarn} — the Vail’s clay-bound hand. She has suspected for a year that he is more than the debt he says he is, and has not asked, because she recognizes a person who keeps a secret and knows better than to open him.

\CRsheetheading{Arc So Far}
She is in the birth room from the first hours of Day 1 and is already arranged: the smile kept on the face like a door on the latch, the back turned to the slate, the eel-bones flicked and put away. She says the one true thing in the room --- that the secret is in what happens to a child after it is born --- in a form that can be taken as nothing, and follows it with \emph{there’s nothing to tell}. The knot-string and the barley account are still offstage; the woman who knows and will not say is not. Her marker is reconciled to Chapter 1: she is on stage here, not first seen at the Fishstairs in Chapter 5.

\CRsheetheading{Chapter Log}
\begin{itemize}[leftmargin=1.4em,itemsep=1pt,topsep=3pt,parsep=0pt]

\item Ch 1: Came in from the yard with eel slime dried on her sleeves and the cold on them; stood in the one patch of floor the slate did not reach, back to the writing and face to the door. Answered a question with a question and then her own answer: the secret is in what happens to a child after. Said she had told nobody and there was nothing to tell. Flicked an eel-bone against her thumb and put it away when the hand came; kept the small smile on the latch all night. Crossed the room to press a crust of brown-edged Vail-smoked barley bread into Dessa’s hand, briskly, stopping short of the slate. The top of her left ear is missing --- a hook, she says; the scar says a knife.

\end{itemize}
}

\CRcharacter{Sella Vail}{the woman who keeps the fire and asks no questions}{}{\CRfield{Also called}{Sella, the Vail widow, Widow Vail, the keeper of the fire}\CRfield{Role}{family}\CRfield{Born}{Y.H. 779, Low Kiln, Sill, Grath}\CRfield{Kind}{human (Grathic)}\CRfield{Origin}{Low Kiln, Sill, Grath; born to a family of slip-mixers}\CRfield{Allegiance}{House Vail kiln-house, Clay Yard, Low Kiln; Kiln-Masters of Sill by widow’s mark}\CRfield{Status}{alive}\CRfield{First seen}{Chapter 1}}{\CRsheetheading{Appearance}
A broad woman of thirty-three with forearms the colour of dried slip and a burn across the back of her left hand shaped like a little door --- two straight sides, a lintel and a threshold, the skin inside it puckered and pale --- which she keeps half-closed when she carries a cup, and which closes on nothing when her hand has nothing to hold. She is not tall and she does not make herself tall; she takes up room by being impossible to move. Her hair is coal-black with one grey streak at the temple that she says came from the winter Garl died and will not discuss further. She smells of woodsmoke, wet wool, and under it the faint sour mineral tang of the Ash ridge along the skirting board, which she has not swept in two years and does not intend to sweep. Her hands are never still. When they have nothing to do they count: spoons, nails, the breaths of a sleeping child. She wears her widow’s mark on a cord at her throat, a thumb-length of lead-glazed clay with the Vail kiln scratched into it, and she touches it before she opens the kiln and never at any other time. People notice the burn first and the mark last, which is the order she prefers.

\CRsheetheading{Voice \& Manner}
Grathic proverbs and half-sentences, practical instruction addressed to the pot or the fire before it is addressed to a person. She says \emph{the house} where another woman would say \emph{I}, and \emph{the boy} where another would say \emph{my son}. She never finishes a sentence that would require her to say what she feels, and she is not unaware of this; she has simply decided that a thing said aloud becomes a thing the Court can enter. Her humour is dry and costs her something. When she is angry she gets quieter, and the quieter she gets the more the household straightens up.

\begin{itemize}[leftmargin=1.4em,itemsep=1pt,topsep=3pt,parsep=0pt]

\item “The fire doesn’t know you’re sorry. Feed it.”

\item “Don’t set that near the door. You’ll pay for it in the thaw.”

\item “He’s a boy. He’s warm. That’s the whole of what I know.”

\item “You can ask me what you like. I won’t answer it. Ask the pot instead; it’s more use.”

\item “A slate-name with no house is a coat anybody can wear. I’ve seen it before.”

\item “The boy comes under the roof on the morrow. The house will keep him.” --- the taking announced without an \emph{I} in it, and never explained.

\item “Don’t set that near the door. You’ll pay for it in the thaw.” --- instruction to the pot, not to a person, given with the kiln already open.

\end{itemize}

\CRsheetheading{History}
She was born in Low Kiln to slip-mixers, married Garl Vail at seventeen, and kept the Vail fire ever since in the way a widow keeps a wound clean: daily, without comment, with a hand that knows exactly how much pressure the thing will bear. She bore two children. Jorra, who lived, and Enna, who died at nine months of the second chill in Y.H. 804, after Alda Reed had warmed the stone and done everything the craft knew. Sella sat up with that child for three nights and learned what the midwives already knew: that a vessel can cool while everybody is being kind to it. In Y.H. 811, in the Long Cold, Garl died with the river iced and no collector able to cross, and his warmth went where unkept warmth goes. The Ash in the Vail corner is thicker than in any other house in the district. Sella has never swept it and has never said why.

On Stone 7, Y.H. 812, a woman in the district bore a child with a branching red flower across his ribs, and Sella took him in on the third day. She gave as her reason the one she has given ever since: a foundling with no house-name can be used by anyone, and a boy entered under a roof is a boy nobody can take for a levy without answering to that roof. This is true. It is not the whole of it. Alda Reed, who delivered him, crossed herself at the birth-flower and told nobody what she had seen, and Sella, who saw the same look on the midwife’s face, asked her nothing at all. She has kept the boy close for three years and has never once asked where his memory came from. She is waiting to see what he will become.

\CRsheetheading{Wounds \& Fears}
Enna’s second chill; Garl going uncollected; the Ash ridge that has grown an inch since. She fears another Long Cold, and fears more than that the day somebody from Slate Row notices that one child in her house listens to the fire before he looks at it. She fears Jorra, who has begun to disappear into the back rooms with men whose names are not spoken, and she fears that she has already lost the wheel to him while he still stands at it. Her contradiction is the whole of her: the woman who feeds every hungry child in the yard has never once said she loves one. She would sooner swallow the widow’s mark.

\CRsheetheading{Desires}
She wants the fire to hold, the wheel to stay in the house, and a Vail name written on the birth-slate of every child born under her roof. She thinks she wants a whole return for Garl, so that he might come back as himself and sit where he used to sit and say nothing in his own voice. What she needs is to admit that she did not take Corin in for the house at all. She took him because he was the only warm thing in a room full of people pretending not to be cold, and because she had already lost one child to the second chill and could not bear to lose another to an explanation.

\CRsheetheading{Bonds}
\textbf{Jorra Vail} — her son, the keeper of the wheel. He believes she took the foundling in to replace Enna and to keep the house’s standing while he did the work. She has never told him otherwise, which is how he has come to believe it. They speak about the kiln, the levy, the weather; they do not speak about Garl, or Enna, or the boy. \emph{Ch 2:} he said \emph{he’s not Vail} when she announced the taking, and she did not answer him; she grew quieter and the house straightened up around it. He went on turning the wheel without looking at her and obeyed her before she had finished. The sentence stands and the boy is under the roof and still unwritten.

\textbf{Corin} — the foundling she took in and has never questioned. She watches him with a hearth-maid’s attention and a widow’s refusal to name what she sees. Her silence is its own kind of Court: kindness that does not want an answer, and that therefore cannot be refused. \emph{Ch 1:} her hand came in ahead of her, felt for the heat at his belly, found none, opened and shut once over the covered bundle, and withdrew. She did not look at him, or at any of the women, and said nothing. \emph{Ch 2:} the hand that hovered now takes. She crossed to the bench first and put her hand near the slate as near a sleeping animal, then lifted him, the little-door burn closing first on nothing and then on his weight. She said \emph{the boy} and \emph{the house}, carried him up the yard without looking back, laid him on fresh straw, and gave instruction to the pot and the fire. She has not written him as Vail and has said nothing of payment or purpose.

\textbf{Alda Reed} — delivered Enna and Corin both, and crossed herself at the second birth. They have traded one silence for three years. Alda has the fuller notes; Sella has the fuller knowledge of what a silence costs when a household must eat. \emph{Ch 1:} Alda laid the slate flat on the bench as Sella’s hand went back out the door; the trade of silences begins in that room. \emph{Ch 2:} Sella put her hand near the writing on the bench as she took the child and did not touch it, and Alda did not look at her. The arrangement holds without a word passing between them.

\textbf{Dessa Silt} — the woman who bore the boy. Sella knows it and has paid her in barley and slip-clay to stay away, and has never entered the payment in any account, not even the private one she keeps in her head. \emph{Ch 1:} Dessa’s pulse rose under her \emph{door shut} smile while the hand hung above her child; neither looked at the other. \emph{Ch 2:} Sella took the boy up the yard without once looking at the woman standing at the end of the nineteen steps, and Dessa did not follow.

\textbf{Berr Clay, the Thane of Clay} — confirmed her widow’s mark and has never called in the debt Garl left against a shipment of tiles. The uncalled debt is worse than a called one, because she cannot pay it and cannot settle it and cannot stop owing it.

\textbf{Garl Vail} — her husband, uncollected in the Long Cold. She decided he would not be collected; he did not ask her to decide. The Ash in the corner is the shape of that decision and she has lived beside it for a year. \emph{Ch 2:} she keeps his seat empty and his coat on its peg, will not let anyone sweep the ridge, and touches the widow’s mark only before she opens the kiln. The house keeps a roll of its own in which he is entered, and does not show it to a man from Slate Row.

\CRsheetheading{Arc So Far}
In Chapter 1 she is a hand in a doorway and a shadow, not a speaking part. Woodsmoke and wet wool arrive ahead of her; then the burn shaped like a little door reaches for the heat at the child’s belly, finds none, hovers over the bundle and withdraws. In Chapter 2 she is inside the house for the first time on the page and is known by its objects before she is known by her words: the hollow seat, the salt coat, the Ash ridge a thumb and then two thumbs wide, the charred edge of a slate nobody picks up, the roll the house keeps and does not show. She announces the taking, is answered by her son, refuses the answer by getting quieter, and on the third day performs the taking with a hand that knows how much a thing will bear. She lays the boy on a bed and tells the fire what to do. She has not explained any of it. Her silence has stopped being a pause and become a wall.

\CRsheetheading{Chapter Log}
\begin{itemize}[leftmargin=1.4em,itemsep=1pt,topsep=3pt,parsep=0pt]

\item Ch 1: Did not enter the birth room and said nothing. Woodsmoke and wet wool came ahead of her, then the cold off the yard, then the hand with the little-door burn across its back --- lintel, threshold, two straight sides. It felt for heat at the child’s belly, found none, opened and shut once above the covered bundle and withdrew. She stepped back out into the Clay Yard without looking at child or women. The crust and the warmed stone had already come from her hearth, neither entered against anybody; the kiln breathed through the wall afterward. She still intends to take him on the third day.

\item Ch 2: Kept the fire and the house through Days 2 and 3. At first light touched the widow’s mark before opening the kiln and at no other time, gave instruction to the pot and the fire, would not let anyone fill Garl’s seat or take his coat down, and said the house kept a roll of its own. Announced that the boy would come under the roof and was answered \emph{he’s not Vail}; got quieter and did not reply. On the third morning came for him, crossed to the bench first, lifted him with the little-door hand, said \emph{the boy} and \emph{the house}, carried him up the Clay Yard and laid him on fresh straw. Gave the bed to the foundling and the loft to Tarn without explaining either, told the pot it would pay for itself in the thaw, and left the charred slate in the kiln ash untouched.

\end{itemize}
}

\CRcharacter{Alda Reed}{the hearth-maid who crossed herself and told nobody}{}{\CRfield{Also called}{Alda, Alda the hearth-maid, the keeper of the second notes}\CRfield{Role}{ally}\CRfield{Born}{Y.H. 770, Low Kiln, Sill, Grath}\CRfield{Kind}{human (Grathic)}\CRfield{Origin}{Low Kiln, Sill, Grath; House Reed}\CRfield{Allegiance}{the Hearth-Maids; birth-slates of Low Kiln; collector’s assistant at Sill}\CRfield{Status}{alive}\CRfield{First seen}{Chapter 1}}{\CRsheetheading{Appearance}
Forty-two and built like a woman who has spent her life lifting other people’s children: broad shoulders, a permanent stoop from leaning over births, and hands so steady that a mother will grab her wrist rather than her arm when the pain comes. Her nails are scrubbed to the quick and still have clay under them. She wears a grey wool shawl in every season and a little tin box on a cord inside her bodice, which she touches when a question is asked and never opens in company. She smells of soap, warm milk, tallow and the faint mineral tang of the Ash she has promised herself she will sweep tomorrow. The thing people notice first is her calm; the thing they notice last is the crossed sign she makes with two fingers at the corner of a birth-flower, which she does not know she does in front of strangers.

\CRsheetheading{Voice \& Manner}
Practical instruction in proverbs and half-sentences. She addresses the newborn, the stone, the slip, the mother, and only then the father, in that order. She does not say \emph{please} and she does not say \emph{sorry}, and she does not consider this coldness; she considers it the correct amount of words for a room where somebody is in pain. She has a habit of answering a question with the date. When she is lying she becomes unusually exact.

\begin{itemize}[leftmargin=1.4em,itemsep=1pt,topsep=3pt,parsep=0pt]

\item “Warm the stone first. A cold stone is a stone, not a kindness.”

\item “Don’t set that near the fire. You’ll pay for it in the thaw.”

\item “He grew worse on the third day. So did the one before him. That’s all I’ve written.”

\item “The Court struck the name out. They didn’t strike the child out.”

\item “I crossed myself. It’s what my mother did. Don’t make it more than that.”

\item “In the levy. Among the children taken for the account this month.” --- the lie, delivered as an entry and never as a story.

\item “A woman with a thin return goes asking after the levy and a hand comes down on a slate that was better left where it lay.” --- mercy spoken in the grammar of a warning.

\end{itemize}

\CRsheetheading{History}
She was born in Low Kiln to a family of reed-cutters and became a hearth-maid at nineteen, which in Grath means midwife, keeper of the birth-slate, and the person a household calls before it calls a collector. She has attended something over four hundred births. The first of the second notes was not hers. In Y.H. 651, long before she was born, a hearth-maid in Low Kiln attended an infant who grew worse on the third day and spoke words in a language no one in the district knew. The Court’s physician at Greybank wrote it down as a settling of the graft, struck the phrase \emph{Second Want} from his own report, and paid the maid for her trouble. That note passed down the order, from hand to hand and tin to tin, until it came to Alda. She has added to it every year since she took it up: every infant who grew worse on the third day, every word that did not belong to Grathic, every birth-flower shaped like a branch instead of a rose. The notes live under a floorboard in her house, in a second hand that the Court has never seen. Alda Reed’s is the fullest set in the order.

She attended Enna Vail’s death in Y.H. 804 and sat with Sella through a night in which nothing she did mattered. She attended Corin’s birth on Stone 7, Y.H. 812, crossed herself at the branching red flower across his ribs, and told nobody what she had seen. Three days later she told Dessa Silt that her child had gone in the levy, which was a lie told to keep a whole-return infant off a slate that every officer in Slate Row would have wanted to read. She has not slept properly since.

\CRsheetheading{Wounds \& Fears}
Four hundred births and one she cannot account for. On Day 1 the sign is made in front of witnesses and covered with a word and a date nobody asked for; from that hour she is a woman who has hidden a thing in a room where everybody is arranging silence. She fears collectors — not their tubes, which she has held, but their slates, which are always in a better hand than hers. She fears that the second notes are a treasonous thing and that keeping them has made her a person who waits for children to get worse. She fears Dessa’s face in the street. Her contradiction: the woman who insists that everything be written down within three days has spent thirty years keeping a second record that cannot be shown to anyone, and calls it \emph{not wanting trouble}.

\CRsheetheading{Desires}
She wants one of the Court’s physicians to read her notes and say the words she struck out herself. She thinks she wants the boy she delivered to grow up ordinary, so that she can stop looking at him. What she needs is to understand that she has been a keeper of the truth in the only way available to a hearth-maid: badly, in secret, and at a cost to the woman she lied to.

\CRsheetheading{Bonds}
\textbf{Corin} — the child she delivered and crossed herself over. She has watched him for three years without once saying what she saw, and has begun to suspect he knows that she knows. \emph{Ch 1:} she lifted him, cleaned him in a deliberate order (face, neck, groin, knees, cord), pressed the wool flat with the same two fingers that had made the sign, and would not say his name. \emph{Ch 2:} she washed him in the same order on the second morning, reached for the slate first and pushed the writing further under, touched the tin at her bodice, and made the crossed sign again at the corner of the flower without knowing she was doing it in front of him. She said warm and breathing was all she had set down, which was not true.

\textbf{Sella Vail} — they have traded a silence since the boy was three days old. Sella has never asked what Alda wrote; Alda has never asked why Sella took the child. This is not friendship. It is an arrangement with a door in it that both women have agreed not to open. \emph{Ch 1:} the arrangement begins in the birth room, before a word is spoken between them: Alda lays the slate flat as the hand in the door goes back out. \emph{Ch 2:} when Sella came for the boy she put her hand near the slate on the bench and did not touch it, and neither woman looked at the other. The arrangement holds.

\textbf{Dessa Silt} — she lied to this woman about her own child and has taken her barley money from Sella’s house ever since. When Dessa speaks to her in the street, Alda answers with the weather and hates the sound of her own voice. \emph{Ch 1:} she delivered Dessa’s child and met the unasked question with dates and weather, her hand on the slate rather than on the mother. \emph{Ch 2:} she stood behind Dessa in the Clay Yard as the bundle went up the hill and told her the boy had gone in the levy among the children taken for the account; she became exact about dates nobody had asked for, and the figure was never spoken. The lie is told to keep a slate unread and has begun to curdle.

\textbf{Mira Silt} — they work the same room standing and trade small facts at the door. \emph{Ch 1:} neither sat where the slate lay open, and each kept her face arranged for the other. \emph{Ch 2:} Mira is not in the room; her crust of brown-edged barley is the last warmth that crossed it from her hands.

\textbf{Jorra Vail} — he knows she wrote something down the night Enna died and has never asked to see it. She has begun to be afraid of the day he does, because the boy in his house is the one thing Jorra could be persuaded to burn for.

\textbf{Berr Clay, the Thane of Clay} — pays her to keep the birth-slates and has never asked about the second notes. She has decided this is because he knows and does not want a copy of anything, which may be true.

\textbf{Orra Weir} — was in the room on Stone 7 and tied a knot she cannot account for. They never discuss it. Orra is the only person Alda has ever let near the subject of the third day, and even she has not seen the tin.

\CRsheetheading{Arc So Far}
The two gestures that will define her are on the page in Chapter 1: the crossed sign she does not know she makes in front of strangers, and the exact, unasked-for dates with which she covers it. In Chapter 2 both gestures repeat and a third is added. She reaches for the slate before the child again, touches the little tin on its cord twice more, crosses the sign over the flower in front of the boy without knowing it, and answers every unasked question with an exact history --- a cooper’s seventh in Mud End, the account going up on the eighth before the paper went out --- until the dates arrive faster than dates have any business arriving. Then she tells Dessa the boy went in the levy. The lie is told as mercy and cannot be undone; the tin under the floorboard is years away from being read.

\CRsheetheading{Chapter Log}
\begin{itemize}[leftmargin=1.4em,itemsep=1pt,topsep=3pt,parsep=0pt]

\item Ch 1: Delivered Corin in the front room of House Silt and worked with steady hands, clay under her nails. Reached for the slate before the child. Put her two fingers to the branching flower, noticed she had done it before witnesses, and covered the act with \emph{birth-flower} plus an unasked-for history from Nine Pots: Y.H. 806, fourth day of Stone, a potter’s third child, now a carter in Stonemouth. Touched the tin at her bodice twice. Gave a date and a weather report for the unasked question and said warm and breathing was all she had set down. Closed the slate flat on the bench while Sella Vail’s hand hovered in the door.

\item Ch 2: On Day 2 washed the child in her order, reached for the slate first, touched the tin a further time, made the crossed sign over the flower before witnesses without knowing it, and gave a date and weather where comfort was expected. On Day 3 stood behind Dessa at the end of the Clay Yard and told her the boy was entered in the levy among the children taken for the account that month, that the account went up on the eighth before the paper went out, and that a woman with a thin return who goes asking draws a hand down on a slate better left where it lay. Said warm and breathing was all she had set down. The number was never spoken; she went on giving dates after the question had stopped.

\end{itemize}
}

\CRcharacter{Marra}{the woman at the edge of the flare light}{}{\CRfield{Also called}{Collector Marra, Marra of the Grey Coat, the woman with the slate, the hand that writes}\CRfield{Role}{antagonist}\CRfield{Born}{Y.H. 768, the Hearthseat, on the Keel}\CRfield{Kind}{human (Grathic; officer of the Keeping)}\CRfield{Origin}{the Hearthseat; no house, no kin; entered on the rolls in Ledgerhand}\CRfield{Allegiance}{the Warm Court, collector’s office; secretly the Unbanked}\CRfield{Status}{alive}\CRfield{First seen}{Chapter 1}}{\CRsheetheading{Appearance}
In the death scene she is a grey shape at the edge of the flare light in a coat the colour of wet ash, holding a slate, writing; the writing does not stop when the van takes him. In Sarn she is a woman in her forties with the neat, unremarkable face of an officer who has spent her life at tables rather than in rooms. Her fingertips are the colour of wet ash and have been since Y.H. 803, which she attributes to cold rooms in a voice that invites no question. She wears her hair pinned tight and her cuffs buttoned to the wrist, partly from the trade and partly because she does not want the hand seen. She has the collector’s habit of touching the slate before she touches the person, and a smaller habit — noticed only by people who have watched her for years — of touching her own wrist afterwards, where the iron clip would go, as if checking that it is not there. She smells of cold iron, slate dust and nothing else, which is the most expensive smell in Sarn.

\CRsheetheading{Voice \& Manner}
She apologises in clauses she does not finish. She is precise about other people’s warmth and vague about her own, and the vagueness is not evasion but a genuine absence: she has entered so many measures that she has lost the habit of entering herself. She speaks the Ledgerhand passive as naturally as breathing — \emph{the measure is banked, the return is recorded, the account stands} — and then, rarely, a plain sentence breaks through, usually about a child, and is immediately corrected into the grammar again. When she lies she becomes even more exact, giving dates and figures that nobody asked for.

\begin{itemize}[leftmargin=1.4em,itemsep=1pt,topsep=3pt,parsep=0pt]

\item “The measure was drawn at the fourth beat. I’m sorry, I — the account stands. That’s all that’s required.”

\item “I don’t remember my fingers. I mean, I remember that they were — they were warm once. That isn’t in the report.”

\item “You were written before you were seen. That’s an irregularity. It isn’t a cruelty. There’s a difference, and I’ve entered it.”

\item “Forty measures. One return. If you add it and it doesn’t balance, you’re adding it wrong. I’ve added it. I know what it balances to.”

\item “I’ve been paying since. Don’t ask me with what. It isn’t in the report.”

\end{itemize}

\CRsheetheading{History}
She entered the Keeping as a collector’s assistant in Greybank in Y.H. 790, arriving from the Hearthseat with a slate and an unusual care for the measurement of other people’s warmth. Her early reports were models of the craft: exact, unembellished, quiet about her own deductions. Within four years she was collecting in the Sill district; within twenty her fingertips had begun to go grey; within twenty-five she was the collector of the Low Kiln levy and the woman the district crosses the street to avoid. She has never been reprimanded and has never been promoted, which is a distinction she has thought about more carefully than either.

In the winter of Y.H. 812 she was standing at the edge of a flare light on an expressway in a country she could not have named, writing because the trace in front of her would not cool. She did not cause the crash and did not touch the child. She wrote the measure of a man who gave at the moment of breaking and whose promise was too heavy to bank. She wrote a name — Corin — before she had seen the infant’s face. Then she went home to Sarn and forged the assessment: a whole return with memory intact, booked at twelve hundred marks and forty unspent measures, with the deduction entered as zero. The money came from forty children of Low Kiln. The collector’s clerk who carried the slate has never been able to spell the buyer’s name.

She has been paying for the return ever since with her own fingers, a little warmth at a time, drawn without a slate and booked nowhere. In Y.H. 803 a slate in Greybank was found with a line that did not balance; it was corrected in a second hand and filed without comment. The first hand was hers. The second was also hers, written with the other hand so that the file would show two officers.

\CRsheetheading{Wounds \& Fears}
Her fingers. And underneath them, the forty. She can add and she has added, and the difference between what a whole return costs and what the Hearth was told it cost is a number she carries the way Dan carries Tasha’s chart — as a small, hard, private knowledge of the shape of her own conduct. She fears the Weighing, where the year’s measures are read aloud in a clean voice, because one of them is hers and it is not on any roll. She fears that the boy will one day look at her with that clinical attention and add her up the way she has added up everyone else. Her contradiction: the woman who invented an officer’s exactness to hide a forgery has made the forgery the most exact thing in her life, and calls it \emph{the account standing}.

\CRsheetheading{Desires}
She wants the sealed world too — the Unbanked’s paradise, where nothing leaks and nothing is owed — and she wants the instrument they sent for to be whole enough to do the sealing. She thinks she wants the boy to live. What she needs is to understand that she has already spent him: forty children bought one man, and no amount of careful handwriting makes that a balance. When she finally confesses it, it will be in a clause she does not finish, and the boy will finish it for her, in the plainest words, and the world will not end.

\CRsheetheading{Bonds}
\textbf{Corin} — the return she paid for and has never watched grow. She knows his name before she knows his face, and when she meets him in Slate Row she will look at him the way a surgeon looks at a chart: to check whether her own work held. He will be the first person to ask her about the fingers, and the first to hear her say \emph{I don’t remember} and understand that she means \emph{I do}. \emph{Ch 1:} she is a coat the colour of wet ash at the edge of a flare, a hand moving on a slate, and the same pause in a dark hand on a birth-slate in Low Kiln. He has filed the hand without a name and has not put the two together.

\textbf{Rusk} — the other half of the purchase. He taught the boy; she bought him. They have not spoken plainly in twelve years and have arranged their lives so that they never have to. She has begun to be afraid that he is gentler with the boy than the work permits, which in her vocabulary is the same as saying he has begun to love him.

\textbf{Berr Clay, the Thane of Clay} — she took his signature on the levy of forty and let him forget the purchaser’s name, which he did deliberately and with relief. She has never thanked him and never will; a debt acknowledged is a debt entered.

\textbf{Alda Reed} — Alda crossed herself at a birth-flower and wrote nothing in the Court’s book. Marra has read the Court’s book. She has never asked Alda what else she wrote, because a hearth-maid’s second note is the one thing even a collector cannot lawfully want.

\textbf{Haldre} — whose paper named the ending she has been paying for. She has never read it. She has heard it described once, by Rusk, and has spent the years since doing the work without the theory, which is a more Sarnish arrangement than either of them will admit.

\textbf{Steward Derren} — her head of office, whom she has never contradicted and never trusted. He has praised her exactness twice. She has decided that his praise is the most dangerous thing in the Hearthseat.

\CRsheetheading{Arc So Far}
Chapter 1 gives the first of her two entrances before she enters. She is a grey shape at the edge of the flare light: a coat, a head, a hand, writing without stopping while the van burns. She is not named, not to Dan and not to the reader, who is allowed to think she caused the fire. The writing has the same rhythm of short strokes and a decided pause as the strange dark hand on the birth-slate in House Silt, and Dan sets the two side by side without putting them together. The fingers, the forged assessment, the forty and the unfinished apology are all still buried.

\CRsheetheading{Chapter Log}
\begin{itemize}[leftmargin=1.4em,itemsep=1pt,topsep=3pt,parsep=0pt]

\item Ch 1: Appears only as the image at the edge of the flare light --- coat the colour of wet ash, hand on a slate, short strokes, a pause, more strokes, the writing not stopping when the fire takes the van. Unnamed in the Sarn narration and unremembered by name. She wrote a name before anyone saw the child’s face; by morning that pause is sitting on a slab of grey stone within reach of a midwife’s hand.

\end{itemize}
}

\CRcharacter{Tarn Vail}{the debt with hands}{}{\CRfield{Also called}{Tarn, the Vail’s Tarn, Tarn of the Yard}\CRfield{Role}{other}\CRfield{Born}{Y.H. 790, Low Kiln, Sill, Grath}\CRfield{Kind}{human (Grathic; returned thin)}\CRfield{Origin}{Low Kiln, Sill, Grath; no house of his own}\CRfield{Allegiance}{House Vail, clay-bound for seven years from Y.H. 806}\CRfield{Status}{alive}\CRfield{First seen}{Chapter 2}}{\CRsheetheading{Appearance}
Twenty-two and bent at the shoulders from six years in the slip pit, which makes him look older than his term. He has grey Ash in the creases of his knuckles that no amount of scrubbing lifts, and a pale ring round his left wrist where an iron clip sat in Y.H. 806, when he was bound. He is a narrow man with a narrow face and the habit of putting the smallest tool on the bench into his pocket when he thinks nobody is looking, because a man who owns nothing should know where the small things are. He sleeps in the Debt Loft above the wheel-shed and keeps his few possessions in a clay pot with no mark on it. He smells of wet wool, tallow and the sour tang of a corner that has not been swept. The thing people notice first is how quickly he moves; the thing they notice last is that he stops moving altogether when he is told he is doing well.

\CRsheetheading{Voice \& Manner}
Deferential with a nail in it. He calls everyone \emph{house}, including Sella, and says \emph{I owe} where another man would say \emph{I am}. He never argues; he agrees so thoroughly that the argument has nowhere to go. When he is frightened he becomes helpful, which in a kiln-house means doing three things at once and quietly putting the fourth right. His jokes are not jokes; they are statements of fact made in a tone that invites someone else to laugh so that he does not have to.

\begin{itemize}[leftmargin=1.4em,itemsep=1pt,topsep=3pt,parsep=0pt]

\item “Yes, house. I’ll sweep it. I swept it this morning. I’ll sweep it again.”

\item “Six years. One to go. I’ve counted every one, and I won’t count the next.”

\item “He’s a foundling. I’m a debt. Don’t put us in the same sentence, Sella.”

\item “I took the small knife. I put it back. You can check.”

\item “You can’t return a man to himself. They tried with my mother. She came back with somebody else’s hands.”

\item “I took the small knife. I put it back. You can check.” --- said of his own accord, to nobody, in the tone of a man stating the weather.

\end{itemize}

\CRsheetheading{History}
His mother was banked in Y.H. 792 and came back thin to the district’s name-pot, with fragments of a life she could not name and hands that frightened him because they were not her hands. He was born of that return, which means he carries a little of somebody else’s interest in a life he cannot place, and has been told all his life that this makes him useful. In Y.H. 806, at sixteen, his mother’s account was entered against him and he was bound to House Vail for seven years, a term he did not choose and has counted every day of. He has worked the slip pits, swept the yard, banked the fire and carried the household’s small accounts in his head for six years.

He was sent to fetch a collector for Garl Vail in the Long Cold and turned back at the frozen ford. Sella never asked whether he had gone. Jorra, who was thirteen, never asked either, which Tarn has decided is worse. He was in the house when Corin arrived and watched the boy be given a bed while he slept in the loft, and he has never said the thing he thinks about it, because a clay-bound hand who complains about beds is a clay-bound hand who can be returned to the slate.

\CRsheetheading{Wounds \& Fears}
He is a debt with a body, and one year of a seven-year term left. He fears that when the term ends there will be no house to take him and he will be entered again for vagrancy, which is how the district keeps people. He fears the Second Want in himself — the shape that comes over him some nights and leaves him reaching for a tool he has never used. He fears Sella’s kindness more than Jorra’s command, because the kindness is the thing he cannot enter in any account and therefore cannot pay back. His contradiction: the man who owns nothing is the one who notices when a tool is missing and the one who keeps the house’s accounts most exactly.

\CRsheetheading{Desires}
He wants a slate with his own name on it and no mark against it. He wants to be called by something that is not a house’s possession. He thinks he wants freedom; what he needs is to learn that freedom in Sarn is another category of account, and that a man can be unowned and still be cold, and that the thing to want is not an empty column but a way to keep warm without being spent.

\CRsheetheading{Bonds}
\textbf{Sella Vail} — she holds his debt and once, in the winter Garl died, gave him a warm stone for the loft which he has kept in his clay pot for two years. He cannot reconcile the stone with the debt, and has stopped trying, which is not the same as forgiving it. \emph{Ch 2:} he calls her \emph{house}, says \emph{I owe} where another man says \emph{I am}, and agreed with everything so thoroughly that there was nowhere for an argument to go. She gave the foundling a bed and him the loft, and he put the one sentence into the room that nobody had asked for.

\textbf{Jorra Vail} — son of the house, who treats him as a brother when the wheel needs him and as a servant when it does not. Tarn has decided to hate him for it and keeps failing, because Jorra once sat up all night with him through a fever and has never mentioned it. \emph{Ch 2:} Jorra noticed a rasp left out of its place and Tarn answered him; the kindness and the command arrived in the same voice and Tarn did not separate them.

\textbf{Corin} — the foundling who arrived and was given a bed in the house. Tarn was the first to notice that the boy turns his head toward a pot before anyone taps it, and has told nobody, because a clay-bound hand who can hear the full note is worth more than his debt and he does not want to be bought. \emph{Ch 2:} he said \emph{He’s a foundling. I’m a debt. Don’t put us in the same sentence, Sella.} His resentment is not envy of the bed; it is a man keeping the one thing he owns, which is the exact account of what he is owed and what he owes.

\textbf{Dessa Silt} — he has carried messages between her and the house, and barley from her to Sella, and has hated being the thin part of that chain. He knows her cough is not the pits and has never said so.

\textbf{Mira Silt} — she has suspected for a year that he is more than he says, and has not asked. He has noticed her noticing, and has answered her with the weather every time.

\textbf{Garl Vail} — he was sent for a collector and turned back. The dead man’s coat still hangs behind the door and Tarn has never once touched it.

\CRsheetheading{Arc So Far}
He enters already sweeping, before anyone wakes, a bent man of twenty-two with grey Ash in the creases of his knuckles and a pale ring round his left wrist where the clip sat in Y.H. 806. Six years of a seven-year term lie behind him and he has counted every one and will not count the next. In Chapter 2 he pockets the small bench knife and then says of his own accord that he took it and put it back and they can check; he stops his broom for two strokes when a pot is set down and gives its flat note, turns his head a degree and says nothing; he banks the coals without looking at the fire. When the foundling is given a bed and he is given the loft he makes one short sentence of it and goes back to sweeping. His deference has a nail in it, and the nail is the one possession the house cannot check.

\CRsheetheading{Chapter Log}
\begin{itemize}[leftmargin=1.4em,itemsep=1pt,topsep=3pt,parsep=0pt]

\item Ch 1: Does not appear; the broom, the knife and the loft are held for Chapter 2.

\item Ch 2: Swept the kiln-house before anyone woke and stopped for two strokes at a pot’s flat note without turning his head or naming the sound. Put the small bench knife in his pocket, then took it out and put it back where it could be checked and said so unasked. Told Sella \emph{Six years. One to go. I’ve counted every one, and I won’t count the next.} Carried up from the Debt Loft with a broom on the third day, banked the coals, and said \emph{He’s a foundling. I’m a debt. Don’t put us in the same sentence, Sella.} when the boy was given a bed. Heard the second flat note of the three days from a pot set down on the bench and did not say what it was.

\end{itemize}
}

\chapter*{The Atlas}
\markboth{The Atlas}{The Atlas}
\addcontentsline{toc}{chapter}{The Atlas}
\CRmap*{Sill and the Lower Ilm — the country of the birth-town}{build/atlas/home-region.jpg}{The clay country about Sill, drawn out a few days’ walking: the Ilm from Nine Pots and the Ferr to Greybank, the weir, the Old Kiln, Cold Acre, Mud End and the Reedholt marsh. Surveyed and inked in Y.H. 812, the Thin Year.}

\CRmap{Sill and the Two Kilns — the birth-town}{build/atlas/home-settlement.jpg}{A detailed plan of Sill on the Ilm, drawn to the ground a child can walk: the four-foot fall and the Singing Sill, the two bridges, the Fishstairs and the Cross, Slate Row, Low Kiln with the Clay Yard and the Vail kiln-house, High Kiln and the Glaze Room, the family houses, the smithy, the meal mill, the wells and fields, the Bone Orchard and the edge of the reedbeds. Surveyed and inked in Y.H. 812, the Thin Year.}

\CRmap*{Sarn — the bowl between Hearth and Waste}{build/atlas/world.jpg}{The known world as its own surveyors would draw it: the Rim and the cold beyond, the Sallow Sea, the Keel and the Hearth, and the nations between. Surveyed in Y.H. 812.}

\chapter*{The Reckoning of Days}
\markboth{The Reckoning of Days}{The Reckoning of Days}
\addcontentsline{toc}{chapter}{The Reckoning of Days}
\begin{CRtimeline}
\CRday{Day 1}{Stone 7, Y.H. 812 (Seventhday)  —  the birth room off the Clay Yard in Low Kiln, Sill, Grath  —  Corin is born two hours old with Daniel Marchetti’s whole memory and a branching red flower on his ribs; the stone Dessa warmed at the Vail fire goes cold at his belly while a name is entered in a neat dark hand on Alda’s slate before anyone has looked at his face, Mira gives bread and keeps her back to the writing, and Sella’s door-burned hand feels for the warmth, hovers over the bundle and withdraws with the third day still to come.}
\CRday{Day 2}{Stone 8, Y.H. 812 (Eighthday)  —  the Vail kiln-house and wheel-shed, and the birth room of House Silt, Low Kiln, Sill, Grath  —  In the kiln-house the fire keeps its four-and-four-and-a-gap breath around Garl’s hollow wheel-seat, salt coat and unswept Ash ridge while Tarn pockets and returns the bench knife and Jorra declares the foundling is not Vail to a Sella who answers the pot instead, and three doors down Corin, two days old, hears the flat note of a cup he cannot reach, feels the warmed stone go past temperature, endures Alda’s ordered washing and unasked-for dates and Des’s laugh from a Chicago car park, while Dessa calls her cough weather, counts nineteen steps and refuses his name as the room arranges an unspoken question about who pays to keep him.}
\CRday{Day 3}{Stone 9, Y.H. 812 (Ninthday)  —  the Clay Yard and the Vail kiln-house with its wheel-shed and Debt Loft, Low Kiln, Sill, Grath  —  Sella carries the covered bundle up the icy yard with the burn-door closing on its weight and the boy counting her steps, while Dessa stands at the end of nineteen steps she does not take with her thumb on the sieve-wire scar and Alda tells her he went in the levy among the children taken for the account and lets unasked-for dates settle over an unspoken figure, leaving the cold stone behind on the abandoned bed, and under the Vail roof the boy is laid on fresh straw with heat returning painfully to his hands and feet, hears Jorra’s sentence sealed and Tarn’s I’m-a-debt-not-a-foundling and a second flat note from a pot, and sees the charred edge of a slate in the kiln ash that he cannot reach, the market and the Cross still ahead.}
\end{CRtimeline}

\printindex[persons]
\printindex[places]
\printindex[lore]
\end{document}
